%%%%%%%%%%%%%%%%%%%%%%%%%%%%%%%%%%%%%%%%%%%%%%%%%%%%%%%%%%%%%%%%%%%%%%
% Lớp 11 | Hàm số mũ và hàm số lôgarít | Phép tính luỹ thừa
% Dạng 1D6?1-2: Biến đổi, rút gọn biểu thức chứa lũy thừa
% Số câu: 184  (sinh tự động bởi tach_cau.py; từ nay thầy sửa trực tiếp file này)
%%%%%%%%%%%%%%%%%%%%%%%%%%%%%%%%%%%%%%%%%%%%%%%%%%%%%%%%%%%%%%%%%%%%%%

%%%%%%%%%%%%%%%% TRẮC NGHIỆM 4 PHƯƠNG ÁN %%%%%%%%%%%%%%%%

% nguồn: Ok_Toan11_TracNghiem_Lan2
\begin{ex}%[1D6N1-2]%[ma:a6cfe051]
	Với $a$ là số thực dương tùy  ý, $a^{\tfrac{5}{3}}$ bằng
	 \choice
		 {$\sqrt[5]{a^3}$}
		 {\True $\sqrt[3]{a^5}$}
		 {$a^8$}
		{$a^2$} 
	 \loigiai{
	Với $a>0$  ta có $a^{\tfrac{5}{3}}=\sqrt[3]{a^5}$.}
\end{ex}

% nguồn: Ok_Toan11_TracNghiem_Lan2
\begin{ex}%[1D6N1-2]%[ma:dbc42f18]
	Với $\alpha $ là số thực bất kì, mệnh đề  nào sau đây \textbf{sai}?
	 \choice
		 {$\sqrt{10^{\alpha}}=\left(\sqrt{10}\right)^{\alpha}$}
		{$\sqrt{10^{\alpha}}=10^{\tfrac{\alpha}{2}}$} 
		{$\left(10^{\alpha}\right)^2=\left(100\right)^{\alpha}$} 
		{\True $\left(10^{\alpha}\right)^2=\left(10\right)^{\alpha^2}$} 
	\loigiai{ 
		Ta thấy $\sqrt{10^{\alpha}}=\left(\sqrt{10}\right)^{\alpha}$, $\sqrt{10^{\alpha}}=10^{\tfrac{\alpha}{2}}$, $\left(10^{\alpha}\right)^2=\left(100\right)^{\alpha}$  là các mệnh đề đúng.\\
	Xét mệnh đề $\left(10^{\alpha}\right)^2=\left(10\right)^{\alpha^2}$: với $\alpha =1$, ta có $\left(10^1\right)^2=100\ne{\left(10\right)^{1^2}}=10$ nên mệnh đề  $\left(10^{\alpha}\right)^2=\left(10\right)^{\alpha^2}$ sai.}
\end{ex}

% nguồn: Ok_Toan11_TracNghiem_Lan2
\begin{ex}%[1D6N1-2]%[ma:c2d7769b]
	 Cho $a>0$, khi đó $\sqrt[4]{a}$ bằng
	 \choice
		{\True $a^{\tfrac{1}{4}}$} 
		{$a^4$} 
		 {$\dfrac{1}{a^4}$}
		{$a^{-4}$} 
	\loigiai{ 
	Ta  có $\sqrt[4]{a}=a^{\tfrac{1}{4}}$.}
\end{ex}

% nguồn: Ok_Toan11_TracNghiem_Lan2
\begin{ex}%[1D6N1-2]%[ma:3155eb62]
	Cho hàm số  $f(a)=\dfrac{a^{-\tfrac{1}{3}}\cdot\left(\sqrt[3]{a}-\sqrt[3]{a^4}\right)}{a^{\tfrac{1}{8}}\cdot\left(\sqrt[8]{a^3}-\sqrt[8]{a^{-1}}\right)}$ với $a>0$, $a\neq 1$. Tính giá trị $M=f\left(2\,017^{2\,016}\right)$.
	\choice 
		 {$M=2\,017^{1\,008}-1$}
		 {\True $M=-2\,017^{1\,008}-1$}
		{$M=2\,017^{2\,016}-1$} 
		 {$M=1-2\,017^{2\,016}$}
	 \loigiai{
		Ta có $f(a)=\dfrac{a^{-\tfrac{1}{3}}\cdot\left(a^{\tfrac{1}{3}}-a^{\tfrac{4}{3}}\right)}{a^{\tfrac{1}{8}}\cdot\left(a^{\tfrac{3}{8}}-a^{\tfrac{-1}{8}}\right)} =\dfrac{1-a}{a^{\tfrac{1}{2}}-1}=\dfrac{1-a}{\sqrt{a}-1}  =-1-\sqrt{a}$.\\
		Nên $M=f\left(2\,017^{2\,016}\right) =-1-\sqrt{2\,017^{2\,016}} =-1-2\,017^{1\,008}$. 
	} 
\end{ex}

% nguồn: Ok_Toan11_TracNghiem_Lan2
\begin{ex}%[1D6N1-2]%[ma:6858df9a]
	 Với $a$ là số thực dương tuỳ ý, $\sqrt{a^3}$ bằng
	 \choice
		{$a^{\tfrac{1}{6}}$} 
		{$a^{\tfrac{2}{3}}$} 
		{$a^6$} 
		{\True $a^{\tfrac{3}{2}}$} 
	\loigiai{ 
	 Ta có $\sqrt{a^3}=a^{\tfrac{3}{2}}$.}
\end{ex}

% nguồn: Ok_Toan11_TracNghiem_Lan2
\begin{ex}%[1D6N1-2]%[ma:2963d216]
	Cho hàm số $f(a)=\dfrac{a^{\tfrac{2}{3}}\cdot\left(\sqrt[3]{a^{-2}}-\sqrt[3]{a}\right)}{a^{\tfrac{1}{8}}\cdot\left(\sqrt[8]{a^3}-\sqrt[8]{a^{-1}}\right)}$ với $a>0$, $a\neq 1$. Tính giá trị  $M=f\left(2\,017^{2\,018}\right)$.
	 \choice
		{$2\,017^{2\,018}+1$} 
		{\True $-2\,017^{1\,009}-1$} 
		{$2\,017^{1\,009}$} 
		{$2\,017^{1\,009}+1$} 
	 \loigiai{
		Ta có $f(a)=\dfrac{a^{\tfrac{2}{3}}\cdot\left(a^{-\tfrac{2}{3}}-a^{\tfrac{1}{3}}\right)}{a^{\tfrac{1}{8}}\cdot\left(a^{\tfrac{3}{8}}-a^{-\tfrac{1}{8}}\right)} =\dfrac{1-a}{a^{\tfrac{1}{2}}-1} =-1-a^{\tfrac{1}{2}}$.\\ 
		Do đó  $M=f\left(2\,017^{2\,018}\right)=-1-2\,017^{1\,009}$.
	} 
\end{ex}

% nguồn: Ok_TracNghiemLop11_Lan1
\begin{ex}%[1D6N1-2]%[ma:87e02717]
	Với $x>0$, đặt $P=x\sqrt[5]{x^2}$. Khẳng định nào sau đây  {\bf đúng}?
	 \choice
		{$P=x^6$} 
		 {\True $P=x^{\frac{7}{5}}$}
		{$P=x^{\frac{5}{6}}$} 
		{$P=x^{\frac{6}{5}}$} 
	\loigiai{ 
		Ta có $P=x\sqrt[5]{x^2}=x\cdot x^{\frac{2}{5}}=x^{1+\frac{2}{5}}=x^{\frac{7}{5}}$. 
	 }
\end{ex}

% nguồn: Ok_TracNghiemLop11_Lan1
\begin{ex}%[1D6N1-2]%[ma:ecc1f871]
	Cho $a>0$, $m$, $n\in \mathbb{R}$. Khẳng định nào  sau đây đúng?
	\choice 
		 {$a^m+a^n=a^{m+n}$}
		{$a^m\cdot  a^n=a^{m-n}$}
		{\True $(a^m)^n=(a^n)^m$} 
		 {$\dfrac{a^m}{a^n}=a^{n-m}$}
	 \loigiai{
		Tính chất lũy thừa,  ta có $(a^m)^n=a^{m\cdot n}=(a^n)^m$.
	 }
\end{ex}

% nguồn: Ok_TracNghiemLop11_Lan1
\begin{ex}%[1D6N1-2]%[ma:eef79377]
	Rút gọn biểu  thức $P=x^{\frac{1}{3}}\cdot\sqrt[6]x$ với $x>0$.
	 \choice
		{\True  $P=\sqrt{x}$}
		 {$P=x^{\frac{1}{8}}$}
		 {$P=x^{\frac{2}{9}}$}
		{$P=x^2$} 
	 \loigiai{
	Ta có  $P=x^{\frac{1}{3}}\cdot\sqrt[6]x=x^{\frac{1}{3}}.x^{\frac{1}{6}}=x^{\frac{1}{3}+\frac{1}{6}}=x^{\frac{1}{2}}=\sqrt{x}$.}
\end{ex}

% nguồn: Ok_TracNghiemLop11_Lan1
\begin{ex}%[1D6N1-2]%[ma:6bbeaf2f]
	Cho $ a,b>0$;  $\alpha,\beta \in \mathbb{R}$. Phát biểu nào đúng?
	\choice 
		{$a^{\alpha  +\beta}=a^{\alpha}+a^{\beta}$}
		{$a^{\alpha}b^{\beta}=(ab)^{\alpha  \beta}$}
		{\True $\left(a^{\alpha}\right)^{\beta}=a^{\alpha \beta}$} 
		{$\dfrac{a^{\alpha}}{a^{\beta}}=a^{\alpha  +\beta}$}
	 \loigiai{
		Ta có  $\left(a^{\alpha}\right)^{\beta}=a^{\alpha \beta}$.
	 }
\end{ex}

% nguồn: Ok_TracNghiemLop11_Lan1
\begin{ex}%[1D6N1-2]%[ma:9c8a6ee5]
	Với $a$ là  số thực dương tùy ý, $\sqrt{a^3}$ bằng
	 \choice
		{\True  $a^{\tfrac{3}{2}}$}
		 {$a^6$}
		 {$a^{\tfrac{1}{6}}$}
		{$a^{\tfrac{2}{3}}$} 
	\loigiai{ 
	 Với $a>0$ ta có $\sqrt{a^3}=a^{\tfrac{3}{2}}$.}
\end{ex}

% nguồn: Ok_TracNghiemLop11_Lan1
\begin{ex}%[1D6N1-2]%[ma:1ed4ffc9]
	Biến đổi $\sqrt[5]{2^3}$ bằng 
	\choice 
		{$2^{\tfrac{3}{2}}$} 
		 {$2^{\tfrac{2}{5}}$}
		{\True  $2^{\tfrac{3}{5}}$}
		{$2^{\tfrac{5}{3}}$} 
	 \loigiai{
		Biểu  thức $\sqrt[5]{2^3}=2^{\tfrac{3}{5}}$.
	} 
\end{ex}

% nguồn: Ok_TracNghiemLop11_Lan1
\begin{ex}%[1D6N1-2]%[ma:dc624c28]
	Với $a$ là số thực dương tùy  ý, biểu thức $a^{\tfrac{5}{3}}\cdot a^{\tfrac{1}{3}}$ bằng
	\choice 
		{$a^{5}$} 
		{$a^{\tfrac{5}{9}}$} 
		 {$a^{\tfrac{4}{3}}$}
		{\True $a^2$} 
	 \loigiai{
		Ta có $a^{\tfrac{5}{3}}\cdot a^{\tfrac{1}{3}}=a^{\tfrac{5}{3}+\tfrac{1}{3}}={a^2}$. 
	 }
\end{ex}

% nguồn: Ok_TracNghiemLop11_Lan1
\begin{ex}%[1D6N1-2]%[5-CTST]%[ma:31855f04]
	Với  $\alpha$ là một số thực bất kì, mệnh đề nào sau đây sai?
	\choice 
		{$\sqrt{10^\alpha} = \left(\sqrt{10  }\right)^\alpha$}
		{$\sqrt{10^\alpha} =  10^{\tfrac{\alpha}{2}}$}
		{$\left(10^\alpha \right)^2 =  100^\alpha$}
		{\True $\left(10^\alpha  \right)^2 = 10^{\alpha^2}$}
	 \loigiai{
		Xét mệnh đề $\left(10^\alpha \right)^2 = 10^{\alpha^2}$ với $\alpha =  1$ ta có $\left(10^1 \right)^2 = 10 \neq \left(10\right)^{1^2} = 10$. \\
	Do đó $\left(10^\alpha \right)^2  = 10^{\alpha^2}$ là một mệnh đề sai.}
\end{ex}

% nguồn: Ok_TracNghiemLop11_Lan1
\begin{ex}%[1D6N1-2]%[ma:167312b1]
	Cho $x$ là số thực dương, khác $1$, $\sqrt{x^3}$  bằng
	 \choice
		{\True $x^{\frac{3}{2}}$} 
		 {$x$}
		 {$x^{\frac{1}{3}}$}
		{$x^{\frac{2}{3}}$} 
	\loigiai{ 
	Ta có $\sqrt{x^3}=\left(x^3\right)^{\frac{1}{2}}=x^{\frac{3}{2}}.$} 
\end{ex}

% nguồn: Ok_TracNghiemLop11_Lan1
\begin{ex}%[1D6N1-2]%[ma:c17ca034]
	Cho $x, y$ là hai số thực dương và $m, n$ là hai số thực tùy ý. Đẳng thức nào sau đây là {\bfseries sai}? 
	\choice 
		 {$\left(x^{m}\right)^{n}=(x)^{m n}$}
		 {$x^{m}\cdot x^{n}=x^{m+n}$}
		{$(x  y)^{n}=x^{n}y^{n}$}
		{\True$x^{m}y^{n}=(x y)^{m+n}$} 
	\loigiai{Đẳng  thức $x^{m}y^{n}=(x y)^{m+n}$ sai.
	} 
\end{ex}

% nguồn: Ok_TracNghiemLop11_Lan1
\begin{ex}%[1D6N1-2]%[ma:fe81e0a0]
	Cho $a$ là  số thực khác $0$ và $n$ là một số nguyên dương, $a^{-n}$ bằng
	 \choice
		{\True $\dfrac{1}{a^n}$} 
		 {$-a^n$}
		 {$\dfrac{n}{a}$}
		 {$a^{\frac{1}{n}}$}
	 \loigiai{
	Ta có  $a^{-n}=\dfrac{1}{a^n}$.}
\end{ex}

% nguồn: Ok_TracNghiemLop11_Lan1
\begin{ex}%[1D6N1-2]%[ma:3b40cbe7]
	Cho  $a>0, m,n \in \mathbb{R}$. Khẳng định nào sau đây đúng?
	\choice 
		 {$a^m+a^n=a^{m+n}$}
		 {$a^m\cdot a^n=a^{mt-n}$}
		{$\dfrac{a^m}{a^n}=a^{n-m}$} 
		{\True $\left(a^{mi}\right)^n=\left(a^n\right)^{mi}$} 
	\loigiai{  Ta có $\left(a^{mi}\right)^n=\left(a^n\right)^{mi}=a^{mni}$.
	} 
\end{ex}

% nguồn: Ok_TracNghiemLop11_Lan1
\begin{ex}%[1D6N1-2]%[ma:feff18fd]
	Biếu  thức $P=\sqrt[3]{x^2 \sqrt{x^3}}$ với $x>0$ được rút gọn bằng
	\choice 
		{\True  $x^{\frac{7}{6}}$}
		{$x^{\frac{1}{3}}$} 
		 {$x^{\frac{5}{6}}$}
		 {$x^{\frac{9}{3}}$}
	\loigiai{ 
		 Ta có $P=\sqrt[3]{x^2 \sqrt{x^3}}=\sqrt[3]{x^{2}\cdot x^{\frac{3}{2}}}=\sqrt[3]{x^{\frac{7}{2}}}=x^{\frac{7}{6}}$.
	 }
\end{ex}

% nguồn: Ok_TracNghiemLop11_Lan1
\begin{ex}%[1D6N1-2]%[ma:6c938fba]
	Cho $a$ là số thực dương, $m$ và $n$ là những số  thực bất kì. Khẳng định nào sau đây \textbf{đúng}?
	\choice 
		{\True $a^m\cdot a^n=a^{m+n}$} 
		{$a\cdot a\cdot a\cdot  a\cdot a=5a$}
		 {$a^m\cdot a^n=a^{mn}$}
		 {$a^1=1$}
	 \loigiai{
		Khẳng định  đúng là $a^m\cdot a^n=a^{m+n}$.
	 }
\end{ex}

% nguồn: Ok_TracNghiemLop11_Lan1
\begin{ex}%[1D6N1-2]%[ma:145e2e98]
	Rút gọn biểu thức $P=x^{\tfrac{1}{6}} \cdot \sqrt[{3}]{x}$  với $x > 0$.
	 \choice
		 {$P=x^{\tfrac{1}{8}}$}
		{\True $P=\sqrt{x}$} 
		{$P=x^{\tfrac{2}{9}}$} 
		{$P=x^2$} 
	 \loigiai{
		Với $x > 0$, ta có $P=x^{\tfrac{1}{6}}\cdot  x^{\tfrac{1}{3}}=x^{\tfrac{1}{6}+\tfrac{1}{3}}=x^{\tfrac{1}{2}}=\sqrt{x}$.
	} 
\end{ex}

% nguồn: Ok_TracNghiemLop11_Lan1
\begin{ex}%[1D6N1-2]%[ma:4532118a]
	Với $a>0$, $b>0$, $\alpha$, $\beta $ là các  số thực bất kì, đẳng thức nào sau đây sai?
	\choice 
		 {$\dfrac{a^{\alpha }}{a^{\beta }}=a^{\alpha -\beta }$}
		{$a^{\alpha }\cdot a^{\beta  }=a^{\alpha +\beta }$}
		{\True $\dfrac{a^{\alpha }}{b^{\beta }}={\left( \dfrac{a}{b} \right)}^{\alpha  -\beta }$}
		{$a^{\alpha }\cdot{b^{\alpha }}=\left( ab \right)^{\alpha  }$}
	\loigiai{ 
		Với $a>0$, $b>0$, $\alpha$, $\beta $ là các  số thực bất kì, các mệnh đề đúng là
		\begin{itemize}
			 \item $\dfrac{a^{\alpha }}{a^{\beta }}=a^{\alpha -\beta }$;
			\item $a^{\alpha  }\cdot a^{\beta }=a^{\alpha +\beta }$;
			\item $a^{\alpha }\cdot{b^{\alpha }}=\left( ab  \right)^{\alpha }$.
		\end{itemize}
	 }
\end{ex}

% nguồn: Ok_TracNghiemLop11_Lan1
\begin{ex}%[1D6N1-2]%[ma:8586481c]
	Cho $x$, $y$ là hai  số thực dương và $m$, $n$ là hai số thực tùy ý. Đẳng thức nào sau đây \textbf{sai}?
	 \choice
		{\True  $x^my^n=(xy)^{m+n}$}
		{$\left(x^m\right)^n=x^{mn}$} 
		{$x^m:x^n=x^{m-n}$} 
		 {$x^mx^n=x^{m+n}$}
	\loigiai{ 
		Ta có $(xy)^{m+n}=x^{m+n}y^{m+n}$. 
	 }
\end{ex}

% nguồn: Ok_TracNghiemLop11_Lan1
\begin{ex}%[1D6N1-2]%[ma:c5d96bf2]
	Cho các số nguyên dương $m$, $n$ và số thực dương $a$. Mệnh đề nào  sau đây \textbf{sai}?
	\choice 
		{\True $\sqrt{a}  \cdot \sqrt{a}=\sqrt[n+m]{a}$}
		{$\sqrt[m]{\sqrt{a}}=\sqrt[m \cdot  n]{a}$}
		{$\sqrt{a} \cdot \sqrt[m]{a}=\sqrt[m  n]{a^{m+n}}$}
		 {$(\sqrt{a})^m=\sqrt{a^m}$}
	\loigiai{ 
		Theo tính chất của lũy thừa thì mệnh đề sai là  $\sqrt{a} \cdot \sqrt{a}=\sqrt[n+m]{a}$.
	 }
\end{ex}

% nguồn: Ok_TracNghiemLop11_Lan1
\begin{ex}%[1D6N1-2]%[ma:4defe82c]
	Cho $x$, $y>0$ và $\alpha$, $\beta \in \mathbb{R}$. Tìm  đẳng thức {\bf sai} dưới đây.
	 \choice
		{$(xy)^\alpha = x^\alpha \cdot  y^\alpha$}
		{\True $x^\alpha + y^\alpha =  (x+y)^\alpha$}
		{$(x^\alpha)^\beta = x^{\alpha \beta}$} 
		{$x^\alpha  \cdot x^\beta = x^{\alpha+\beta}$}
	 \loigiai{
		Theo tính chất của lũy thừa thì đẳng thức $x^\alpha +  y^\alpha = (x+y)^\alpha$ {\bf sai}.
	} 
\end{ex}

% nguồn: Ok_TracNghiemLop11_Lan1
\begin{ex}%[1D6N1-2]%[ma:3916648b]
	Cho biểu thức $P=x^{-\tfrac{3}{4}}\cdot \sqrt{\sqrt{x^5}}$, với $x > 0$. Khẳng định nào sau  đây là đúng?
	\choice 
		{$P=x^{-2}$} 
		{$P=x^{-\tfrac{1}{2}}$} 
		 {\True$P=x^{\tfrac{1}{2}}$}
		 {$P=x^2$}
	 \loigiai{
		Ta có  $P=x^{-\tfrac{3}{4}}\cdot \sqrt{\sqrt{x^5}}=x^{-\tfrac{3}{4}}\cdot x^{\tfrac{5}{4}}=x^{-\tfrac{3}{4}+\tfrac{5}{4}}=x^{\tfrac{1}{2}}$.
	 }
\end{ex}

% nguồn: Ok_TracNghiemLop11_Lan1
\begin{ex}%[1D6N1-2]%[ma:b82648d9]
	Cho $a$ là số thực dương. Giá  trị rút gọn của biểu thức $ P=a^{\frac{4}{3}}\sqrt{a}$ bằng
	 \choice
		{$a^{\frac{7}{3}}$} 
		{$a^{\frac{5}{6}}$} 
		 {\True $a^{\frac{11}{6}}$}
		{$a^{\frac{10}{3}}$} 
	\loigiai{ 
		Ta có $ P=a^{\frac{4}{3}}\sqrt{a} = a^{\frac{4}{3}}  \cdot a^{\frac{1}{2}} = a^{\frac{4}{3}+\frac{1}{2}} = a^{\frac{11}{6}}$.
	} 
\end{ex}

% nguồn: Ok_TracNghiemLop11_Lan1
\begin{ex}%[1D6N1-2]%[ma:589d5591]
	 Cho $a$ là số thực dương và $m$, $n$ là các số thực tùy ý. Trong các khẳng định sau, khẳng định nào là \textbf{đúng}?
	\choice 
		{\True  $a^m \cdot a^n=a^{m+n}$}
		 {$a^m+a^n=a^{m+n}$}
		{$a^m \cdot a^n=a^{m \cdot n}$} 
		{$a^m+a^n=a^{m \cdot n}$} 
	\loigiai{ 
	Ta có  $a^m\cdot{a^n}=a^{m+n}$}
\end{ex}

% nguồn: Ok_Toan11_TracNghiem_Lan2
\begin{ex}%[1D6H1-2]%[ma:d821926e]
	[Mã 110  2017] Rút gọn biểu thức $P=x^{\frac{1}{3}}\cdot \sqrt[6]{x}$ với $x>0$.
	 \choice
		{\True $P=\sqrt{x}$} 
		{$P=x^{\frac{1}{8}}$} 
		 {$P=x^{\frac{2}{9}}$}
		 {$P=x^2$}
	\loigiai{ 
		Ta có $P=x^{\frac{1}{3}}\cdot \sqrt[6]{x} 
		=x^{\frac{1}{3}}\cdot{x^{\frac{1}{6}}} 
		 =x^{\frac{1}{3}+\dfrac{1}{6}}
		 =x^{\frac{1}{2}}=\sqrt{x}$.
	 }
\end{ex}

% nguồn: Ok_Toan11_TracNghiem_Lan2
\begin{ex}%[1D6H1-2]%[ma:1a02f40a]
	Cho $x$ là số thực dương. Biết $\sqrt{x\cdot \sqrt[3]{x\sqrt{x\sqrt[3]{x}}}} = x^{\tfrac{a}{b}}$ với $a$, $b$ là các số tự nhiên và $\dfrac{a}{b}$  là phân số tối giản. Tính $a+b$.
	 \choice
		{\True  $16$}
		{$15$} 
		{$14$} 
		 {$17$}
	\loigiai{ 
		 Ta có
		\[ 
		 \sqrt{x\cdot \sqrt[3]{x\sqrt{x\sqrt[3]{x}}}}
		=\sqrt{x\cdot \sqrt[3]{x\sqrt{x\cdot x^{\tfrac13}}}} 
		 =\sqrt{x\cdot \sqrt[3]{x\sqrt{x^{\tfrac43}}}}
		=\sqrt{x\cdot \sqrt[3]{x\cdot x^{\tfrac23}}} 
		 =\sqrt{x\cdot x^{\tfrac59}}
		 = x^{\tfrac79}.
		\] 
		Khi đó $a=7$, $b=7$ nên  $a+b=16$.
	 }
\end{ex}

% nguồn: Ok_Toan11_TracNghiem_Lan2
\begin{ex}%[1D6H1-2]%[ma:cd79cdbf]
	Cho  hàm số $f\left(a\right)=\dfrac{a^{\tfrac{2}{3}}\left(\sqrt[3]{a^{-2}}-\sqrt[3]{a}\right)}{a^{\tfrac{1}{8}}\left(\sqrt[8]{a^{3}}-\sqrt[8]a^{-1}\right)}$ với $a>0,a\ne 1$. Giá trị của $M=f\left(2\,021^{2\,022}\right)$ là
	 \choice
		{$2\,021^{1\,011}$} 
		{$2\,021^{1\,011}+1$} 
		 {$-2\,021^{1\,011}+1$}
		 {\True $-2\,021^{1\,011}-1$}
	\loigiai{ 
		Ta  có
		\begin{eqnarray*}
			 f\left(a\right)&=&\dfrac{a^{\tfrac{2}{3}}\left(\sqrt[3]{a^{-2}}-\sqrt[3]{a}\right)}{a^{\tfrac{1}{8}}\left(\sqrt[8]{a^{3}}-\sqrt[8]a^{-1}\right)}\\
			 &=&\dfrac{a^{\tfrac{2}{3}}\left(a^{\tfrac{-2}{3}}-a^{\tfrac{1}{3}}\right)}{a^{\tfrac{1}{8}}\left(a^{\tfrac{3}{8}}-a^{\tfrac{1}{8}}\right)}=\dfrac{1-a}{a^{\tfrac{1}{2}}-1}\\
			&=&\dfrac{-\left(a^{\tfrac{1}{2}}-1\right)\left(a^{\tfrac{1}{2}}+1\right)}{a^{\tfrac{1}{2}}-1}\\ 
			&=&-a^{\tfrac{1}{2}}-1. 
		\end{eqnarray*}
	Khi  đó $M=f\left(2\,021^{2\,022}\right)=-\left(2\,021^{2\,022}\right)^{\tfrac{1}{2}}-1=-2\,021^{1\,011}-1$.}
\end{ex}

% nguồn: Ok_Toan11_TracNghiem_Lan2
\begin{ex}%[1D6H1-2]%[ma:1b88559a]
	Với $a$ là số thực dương,  biểu thức $P=a^{\tfrac{1}{3}}\cdot \sqrt{a}$ bằng
	 \choice
		{$a^{\tfrac{1}{6}}$} 
		 {$a^{\tfrac{2}{5}}$}
		{\True $a^{\tfrac{5}{6}}$} 
		{$a^{\tfrac{4}{3}}$} 
	\loigiai{ 
	Ta có $P=a^{\tfrac{1}{3}}\cdot  \sqrt{a}=a^{\tfrac{1}{3}}\cdot a^{\tfrac{1}{2}}=a^{\tfrac{5}{6}}$.}
\end{ex}

% nguồn: Ok_Toan11_TracNghiem_Lan2
\begin{ex}%[1D6H1-2]%[ma:80d28641]
	Với $a$ là số thực  dương tùy ý, $\sqrt{\dfrac{1}{a^3}}$ bằng
	\choice 
		 {$a^{-3}$}
		 {$a^{\tfrac{3}{2}}$}
		{$a^{\tfrac{1}{6}}$} 
		 {\True $a^{-\tfrac{3}{2}}$}
	 \loigiai{
	Ta có $\sqrt{\dfrac{1}{a^3}}=\left(\dfrac{1}{a^3}\right)^{\tfrac{1}{2}}=\left(a^{-3}\right)^{\tfrac{1}{2}}=a^{-\tfrac{3}{2}}$.} 
\end{ex}

% nguồn: Ok_Toan11_TracNghiem_Lan2
\begin{ex}%[1D6H1-2]%[ma:180d6566]
	Rút gọn biểu thức  $Q=a^{\tfrac{7}{3}}:\sqrt[3]{a}$ với $a>0$.
	\choice 
		 {$Q=a^{-\tfrac{4}{3}}$}
		{\True $Q=a^2$} 
		{$Q=a^{\tfrac{8}{3}}$} 
		 {$Q=a^4$}
	\loigiai{ 
	Ta có  $Q=a^{\tfrac{7}{3}}:\sqrt[3]{a}=a^{\tfrac{7}{3}}:a^{\tfrac{1}{3}}=a^{\tfrac{7}{3}-\tfrac{1}{3}}=a^2$.}
\end{ex}

% nguồn: Ok_Toan11_TracNghiem_Lan2
\begin{ex}%[1D6H1-2]%[ma:a075f5dd]
	[TK 2020-2021] Với $a$ là số thực dương tùy ý,  $\sqrt{a^3}$ bằng
	\choice 
		{$a^6$} 
		 {\True $a^{\frac{3}{2}}$}
		 {$a^{\frac{2}{3}}$}
		{$a^{\frac{1}{6}}$} 
	 \loigiai{
		Ta có $\sqrt[m]{a^n}=a^{\frac{n}{m}}$ với mọi $a>0$ và $m, n\in \mathbb{Z}^+ \Rightarrow  \sqrt{a^3}=a^{\frac{3}{2}}$.
	 }
\end{ex}

% nguồn: Ok_Toan11_TracNghiem_Lan2
\begin{ex}%[1D6H1-2]%[ma:968e3769]
	Rút gọn biểu  thức $P=x^{\tfrac{1}{6}}\cdot \sqrt[3]{x}$ với $x>0$.
	\choice 
		{$P=x^{\tfrac{1}{8}}$} 
		 {\True $P=\sqrt{x}$}
		{$P=x^{\tfrac{2}{9}}$} 
		 {$P=x^2$}
	 \loigiai{
	Ta có $P=x^{\tfrac{1}{6}}\cdot  x^{\tfrac{1}{3}}=x^{\tfrac{1}{6}+\tfrac{1}{3}}=x^{\tfrac{1}{2}}=\sqrt{x}$.}
\end{ex}

% nguồn: Ok_Toan11_TracNghiem_Lan2
\begin{ex}%[1D6H1-2]%[ma:532b404a]
	Cho biểu thức $P = \sqrt[4]{x \cdot  \sqrt[3]{x^2 \cdot \sqrt{x^3}}}
	$, với $x>0$. Mệnh  đề nào dưới đây đúng?
	 \choice
		{$P=x^{\tfrac{2}{3}}$} 
		 {$P=x^{\tfrac{1}{4}}$}
		 {\True $P=x^{\tfrac{13}{24}}$}
		 {$P=x^{\tfrac{1}{2}}$}
	 \loigiai{
		Ta  có
		\allowdisplaybreaks 
		\begin{eqnarray*}
			P&=&\sqrt[4]{x \cdot  \sqrt[3]{x^2 \cdot \sqrt{x^3}}}
			= \sqrt[4]{x \cdot \sqrt[3]{x^2 \cdot x^{\tfrac{3}{2}}}}\\ 
			&=&\sqrt[4]{x \cdot \sqrt[3]{x^{\tfrac{7}{2}}}} 
			=  \sqrt[4]{x \cdot x^{\tfrac{7}{6}}}\\
			 &=&\sqrt[4]{x^{\tfrac{13}{6}}}
			= x^{\tfrac{13}{24}}. 
		\end{eqnarray*}
	 }
\end{ex}

% nguồn: Ok_Toan11_TracNghiem_Lan2
\begin{ex}%[1D6H1-2]%[ma:e6e3739a]
	 Cho biểu thức $P=\dfrac{a^{\sqrt{7}+1}\cdot a^{2-\sqrt{7}}}{\left(a^{\sqrt{2}-2}\right)^{\sqrt{2}+2}}$ với $a>0$. Rút gọn biểu thức $P$ được kết quả là
	\choice 
		 {\True $P=a^5$}
		 {$P=a^3$}
		{$P=a$} 
		{$P=a^4$} 
	\loigiai{ 
	Ta  có $P=\dfrac{a^{\sqrt{7}+1}\cdot a^{2-\sqrt{7}}}{\left(a^{\sqrt{2}-2}\right)^{\sqrt{2}+2}}=\dfrac{a^3}{a^{-2}}=a^5$.}
\end{ex}

% nguồn: Ok_Toan11_TracNghiem_Lan2
\begin{ex}%[1D6H1-2]%[ma:24ca75f5]
	Cho $a$ là số thực  dương. Rút gọn biểu thức $A=a\sqrt{a^3\sqrt{a\sqrt{a}}}$ về dạng $a^{\tfrac{m}{n}}$ trong đó $\dfrac{m}{n}$ là phân số tối giản và $m,n\in{\mathbb{N}^{*}}$. Tính giá trị của biểu thức $T=m^2+n^2$.
	 \choice
		{$2425$} 
		 {$539$}
		{\True  $593$}
		{$1369$} 
	 \loigiai{
		Ta có $A=a\sqrt{a^3\sqrt{a\sqrt{a}}}=a\cdot a^{\tfrac{3}{2}}\cdot a^{\tfrac{1}{4}}\cdot a^{\tfrac{1}{8}}=a^{1+\tfrac{3}{2}+\tfrac{1}{4}+\tfrac{1}{8}}=a^{\tfrac{23}{8}}$.\\ 
	Do đó $m=23$; $n=8$ nên $T=m^2+n^2=23^2+8^2=593$.} 
\end{ex}

% nguồn: Ok_Toan11_TracNghiem_Lan2
\begin{ex}%[1D6H1-2]%[ma:dfe4ba93]
	Giả sử $a$ là  số thực dương, khác $1$. Biểu thức $\sqrt{a\sqrt[3]{a}}$ được viết dưới dạng $a^{\alpha}$. Khi đó
	\choice 
		{\True  $\alpha =\dfrac{2}{3}$}
		{$\alpha  =\dfrac{5}{3}$}
		 {$\alpha =\dfrac{1}{6}$}
		{$\alpha =\dfrac{11}{6}$} 
	 \loigiai{
	Ta có $\sqrt{a\sqrt[3]{a}}=\sqrt{a^{1+\tfrac{1}{3}}}=a^{\tfrac{2}{3}}=a^{\alpha}\Rightarrow \alpha =\dfrac{2}{3}$.} 
\end{ex}

% nguồn: Ok_Toan11_TracNghiem_Lan2
\begin{ex}%[1D6H1-2]%[ma:98bf0d2d]
	Với $a$ là số thực dương tùy ý khi đó  $\sqrt{a^2\cdot \sqrt[5]{a}}$ bằng
	 \choice
		{\True $a^{\tfrac{11}{10}}$} 
		{$a^{\tfrac{1}{10}}$} 
		{$a^{\tfrac{22}{5}}$} 
		 {$a^{\tfrac{10}{11}}$}
	\loigiai{ 
	Với $a>0$  ta có $\sqrt{a^2\cdot \sqrt[5]{a}}=\sqrt{a^2\cdot a^{\tfrac{1}{5}}}=\sqrt{a^{\tfrac{11}{5}}}=a^{\tfrac{11}{10}}$.}
\end{ex}

% nguồn: Ok_Toan11_TracNghiem_Lan2
\begin{ex}%[1D6H1-2]%[ma:ac898b8b]
	Biểu thức $P = \sqrt[3]{x\sqrt[5]{x^2\sqrt{x}}}  = x^{\alpha}
	$, giá trị của  $\alpha $ là
	\choice 
		{\True $\dfrac{1}{2}$} 
		 {$\dfrac{5}{2}$}
		 {$\dfrac{9}{2}$}
		 {$\dfrac{3}{2}$}
	 \loigiai{
		Ta có 
		\allowdisplaybreaks 
		\begin{eqnarray*}
			P&=&\sqrt[3]{x\sqrt[5]{x^2\sqrt{x}}} 
			= \sqrt[3]{x\sqrt[5]{x^2 \cdot  x^{\tfrac{1}{2}}}}\\
			&=&\sqrt[3]{x  \cdot \left(x^{\tfrac{5}{2}}\right)^{\tfrac{1}{5}}}
			= \left(x^{\tfrac{3}{2}}\right)^{\tfrac{1}{3}}\\ 
			 &=&x^{\tfrac{1}{2}}.
		\end{eqnarray*}
	Suy  ra $\alpha = \dfrac{1}{2}$.}
\end{ex}

% nguồn: Ok_Toan11_TracNghiem_Lan2
\begin{ex}%[1D6H1-2]%[ma:4ffaaac6]
	Rút gọn biểu thức $A=x^{\tfrac{1}{3}}\cdot \sqrt[6]{x}$,  $x>0$ ta được
	 \choice
		{\True $A=\sqrt{x}$} 
		 {$A=x^{\tfrac{2}{9}}$}
		 {$A=x^2$}
		 {$A=x^{\tfrac{81}{2}}$}
	\loigiai{ 
	Ta  có $A=x^{\tfrac{1}{3}}\cdot \sqrt[6]x=x^{\tfrac{1}{2}}=\sqrt{x}$.}
\end{ex}

% nguồn: Ok_Toan11_TracNghiem_Lan2
\begin{ex}%[1D6H1-2]%[ma:e8f377b7]
	Cho biểu thức $P=\sqrt[6]{x\cdot \sqrt[4]{x^5\cdot \sqrt{x^3}}},$ với $x>0$. Mệnh đề nào dưới  đây đúng?
	 \choice
		{$P=x^{\tfrac{15}{16}}$} 
		{\True $P=x^{\tfrac{7}{16}}$} 
		 {$P=x^{\tfrac{5}{42}}$}
		 {$P=x^{\tfrac{47}{48}}$}
	 \loigiai{
	Ta có $P=\sqrt[6]{x\cdot  \sqrt[4]{x^5\cdot \sqrt{x^3}}}=x^{\left[\left(\tfrac{3}{2}+5\right)\tfrac{1}{4}+1\right]\tfrac{1}{6}}=x^{\tfrac{7}{16}}$.}
\end{ex}

% nguồn: Ok_Toan11_TracNghiem_Lan2
\begin{ex}%[1D6H1-2]%[ma:583b8bc4]
	[Mã 105 2017] Rút  gọn biểu thức $Q=b^{\frac{5}{3}}\colon \sqrt[3]{b}$ với $b>0$.
	\choice 
		 {$Q=b^{-\frac{4}{3}}$}
		{\True  $Q=b^{\frac{4}{3}}$}
		{$Q=b^{\frac{5}{9}}$} 
		{$Q=b^2$} 
	 \loigiai{
		 $Q=b^{\frac{5}{3}}\colon \sqrt[3]{b}=b^{\frac{5}{3}}\colon{b^{\frac{1}{3}}}=b^{\frac{4}{3}}$.
	 }
\end{ex}

% nguồn: Ok_Toan11_TracNghiem_Lan2
\begin{ex}%[1D6H1-2]%[ma:f7750901]
	Với $a$ là số thực dương tuỳ  ý, $\sqrt{a^3\sqrt[4]{a}}$ bằng
	 \choice
		{$a^{\tfrac{17}{6}}$} 
		{\True  $a^{\tfrac{13}{8}}$}
		{$a^{\tfrac{13}{6}}$} 
		 {$a^{\tfrac{17}{4}}$}
	 \loigiai{
	Ta có $\sqrt{a^3\sqrt[4]{a}}=\sqrt{a^3\cdot  a^{\tfrac{1}{4}}}=\sqrt{a^{\tfrac{13}{4}}}=a^{\tfrac{13}{8}}$.}
\end{ex}

% nguồn: Ok_Toan11_TracNghiem_Lan2
\begin{ex}%[1D6H1-2]%[ma:de25d748]
	Biểu thức $K=\sqrt[2]{2\sqrt[3]{2}}$ viết  dưới dạng lũy thừa với số mũ hữu tỉ là
	\choice 
		{$2^{\tfrac{4}{3}}$} 
		 {$2^{\tfrac{5}{3}}$}
		{$2^{\tfrac{1}{3}}$} 
		 {\True $2^{\tfrac{2}{3}}$}
	 \loigiai{
	Ta có $K=\sqrt[2]{2\sqrt[3]2}=\sqrt{2\cdot  2^{\tfrac{1}{3}}}=\sqrt{2^{\tfrac{4}{3}}}=\left(2^{\tfrac{4}{3}}\right)^{\tfrac{1}{2}}=2^{\tfrac{2}{3}}$.}
\end{ex}

% nguồn: Ok_Toan11_TracNghiem_Lan2
\begin{ex}%[1D6H1-2]%[ma:b11168ce]
	Cho $x$, $y$ là hai số nguyên  thỏa mãn $3^x\cdot6^y
	=\dfrac{2^{15}\cdot6^{40}}{9^{50}\cdot12^{25}}$. Tính  $x y$?
	 \choice
		{$-445$} 
		{$-755$} 
		{\True  $-450$}
		{$-425$} 
	\loigiai{ 
		Ta có 
		\allowdisplaybreaks 
		\begin{eqnarray*}
			&&3^x\cdot6^y=\dfrac{2^{15}\cdot6^{40}}{9^{50}\cdot12^{25}}\\ 
			&\Leftrightarrow& 
			3^x\cdot3^y\cdot2^y=\dfrac{2^{15}\cdot2^{40}\cdot3^{40}}{3^{100}\cdot3^{25}\cdot2^{50}}\\ 
			&\Leftrightarrow&3^{x+y}\cdot2^y=3^{-85}\cdot2^5\\ 
			 &\Rightarrow& \heva{&x=-90\\&y=5}\\
			 &\Rightarrow& xy=-450.
		\end{eqnarray*}
	 }
\end{ex}

% nguồn: Ok_Toan11_TracNghiem_Lan2
\begin{ex}%[1D6H1-2]%[ma:f3618071]
	Rút gọn biểu thức  $P=\dfrac{x^{\tfrac{1}{3}}\sqrt[6]{x}}{\sqrt[4]{x}}$, với $x>0$.
	\choice 
		{\True  $P=\sqrt[4]{x}$}
		{$P=x^{-\tfrac{1}{6}}$} 
		{$P=\sqrt{x}$} 
		 {$P=x^{\tfrac{1}{6}}$}
	 \loigiai{
	Ta có  $P=\dfrac{x^{\tfrac{1}{3}}\sqrt[6]x}{\sqrt[4]x}=\dfrac{x^{\tfrac{1}{3}}\cdot x^{\tfrac{1}{6}}}{x^{\tfrac{1}{4}}}=x^{\tfrac{1}{3}+\tfrac{1}{6}-\tfrac{1}{4}}=x^{\tfrac{1}{4}}=\sqrt[4]{x}$.}
\end{ex}

% nguồn: Ok_Toan11_TracNghiem_Lan2
\begin{ex}%[1D6H1-2]%[ma:f45f0ccb]
	Cho $a$ là số thực dương. Rút gọn biểu thức  $A=a\sqrt{a^{3}\sqrt{a\sqrt{a}}}$ về dạng $a^{\tfrac{m}{n}}$ trong đó $\dfrac{m}{n}$ là phân số tối giản và $m$, $n\in\mathbb{N}^{*}$. Tính giá trị của biểu thức $T=m^2+n^2$.
	\choice 
		{$2425$} 
		 {$539$}
		 {\True $593$}
		{$1369$} 
	\loigiai{ 
		Ta có 
		\[ 
		 A=a\sqrt{a^{3}\sqrt{a\sqrt{a}}}
		 =a\cdot a^{\tfrac32}\cdot a^{\tfrac14}\cdot a^{\tfrac18}
		 =a^{1+\tfrac32+\tfrac14+\tfrac18}
		=a^{\tfrac{23}{8}}. 
		\] 
		Suy ra  $m=23$, $n=8$.\\
		Vậy $T=m^2+n^2=23^2+8^2=593$. 
	 }
\end{ex}

% nguồn: Ok_Toan11_TracNghiem_Lan2
\begin{ex}%[1D6H1-2]%[ma:ecff3332]
	Cho $a$ là  số thực dương khác $1$. Khi đó $\sqrt[4]{a^{\tfrac{2}{3}}}
	$ bằng 
	\choice 
		 {$\sqrt[3]{a^2}$}
		{$a^{\tfrac{8}{3}}$} 
		 {$a^{\tfrac{3}{8}}$}
		{\True  $\sqrt[6]a$}
	 \loigiai{
	Ta  có $\sqrt[4]{{{a}^{\tfrac{2}{3}}}}=\left(a^{\tfrac{2}{3}}\right)^{\tfrac{1}{4}}=a^{\tfrac{2}{3}\cdot \tfrac{1}{4}}=a^{\tfrac{1}{6}}=\sqrt[6]a$.}
\end{ex}

% nguồn: Ok_Toan11_TracNghiem_Lan2
\begin{ex}%[1D6H1-2]%[ma:e2d2eeeb]
	Cho $4^x+4^{-x}=7$.  Khi đó biểu thức $P=\dfrac{5-2^x-2^{-x}}{3+2^{x+1}+2^{1-x}}=\dfrac{a}{b}$
	với $\dfrac{a}{b}$ tối giản  và $a\in\mathbb{Z}$, $b\in\mathbb{Z}^+$. Tính tổng $a+b$.
	\choice 
		 {$8$}
		 {\True $11$}
		{$17$} 
		{$4$} 
	 \loigiai{
		Ta  có
		 \[
		\left(2^x+2^{-x}\right)^2=4^x+4^{-x}+2\cdot 2^x\cdot 2^{-x}=7+2=9. 
		 \]
		Suy ra  $2^x+2^{-x}=3$.\\
		Khi  đó
		\[ 
		 P=\dfrac{5-2^x-2^{-x}}{3+2^{x+1}+2^{1-x}}
		=\dfrac{5-\left(2^x+2^{-x}\right)}{3+2\left(2^x+2^{-x}\right)} 
		=\dfrac{5-3}{3+2\cdot 3} 
		 =\dfrac{2}{9}.
		 \]
		Vậy $a=2$, $b=9\Rightarrow a+b=11$. 
	 }
\end{ex}

% nguồn: Ok_Toan11_TracNghiem_Lan2
\begin{ex}%[1D6H1-2]%[ma:2de4542a]
	Cho $x$ là  số thực dương. Biểu thức $\sqrt[4]{x^2\sqrt[3]{x}}$ được viết dưới dạng lũy thừa với số mũ hữu tỉ là
	\choice 
		 {$x^{\tfrac{12}{7}}$}
		 {$x^{\tfrac{5}{6}}$}
		{\True $x^{\tfrac{7}{12}}$} 
		{$x^{\tfrac{6}{5}}$} 
	\loigiai{ 
	Ta có $\sqrt[4]{x^2\cdot\sqrt[3]{x}}=\sqrt[4]{x^2\cdot x^{\tfrac{1}{3}}}=\sqrt[4]{x^{\tfrac{7}{3}}}=x^{\tfrac{7}{12}}$  với $x>0$.}
\end{ex}

% nguồn: Ok_Toan11_TracNghiem_Lan2
\begin{ex}%[1D6H1-2]%[ma:2758ae44]
	Cho $a \neq 1$ là số thực dương. Viết và rút gọn biểu thức $a^{\tfrac{3}{2\,022}} \cdot \sqrt[2\,022]{a}$ dưới dạng lũy thừa với số mũ hữu tỉ. Tìm số mũ của biểu thức rút  gọn đó.
	\choice 
		{$\dfrac{1}{1\,011}$} 
		 {$\dfrac{3}{2\,022^2}$}
		{\True  $\dfrac{2}{1\,011}$}
		 {$\dfrac{3}{1\,011}$}
	 \loigiai{
	Ta có  $a^{\tfrac{3}{2\,022}}\cdot \sqrt[2022]a=a^{\tfrac{3}{2\,022}}\cdot{a^{\tfrac{1}{2\,022}}}=a^{\tfrac{3}{2\,022}+\tfrac{1}{2\,022}}=a^{\tfrac{2}{1\,011}}$.}
\end{ex}

% nguồn: Ok_Toan11_TracNghiem_Lan2
\begin{ex}%[1D6H1-2]%[ma:6d0403e7]
	Cho biểu thức $P=\sqrt[3]{x\cdot \sqrt[4]{x^3\sqrt{x}}}$, với $x>0$.  Mệnh đề nào dưới đây đúng?
	 \choice
		 {$P=x^{\tfrac{1}{2}}$}
		{$P=x^{\tfrac{7}{12}}$} 
		{\True  $P=x^{\tfrac{5}{8}}$}
		 {$P=x^{\tfrac{7}{24}}$}
	\loigiai{ 
	Ta có $P=\sqrt[3]{x\cdot  \sqrt[4]{x^{\tfrac{7}{2}}}}=\sqrt[3]{x\cdot x^{\tfrac{7}{8}}}=\sqrt[3]{x^{\tfrac{15}{8}}}=x^{\tfrac{5}{8}}$.}
\end{ex}

% nguồn: Ok_Toan11_TracNghiem_Lan2
\begin{ex}%[1D6H1-2]%[ma:e3f5e66b]
	Cho hai số thực dương $a$, $b$. Rút gọn biểu thức $A=\dfrac{a^{\tfrac14}\sqrt[3]{b}+b^{\tfrac14}\sqrt[3]{a}}{\sqrt[12]{a}+\sqrt[12]{b}}$ ta thu được  $A=a^{m}b^{n}$. Tích của $mn$ là
	\choice 
		{$\dfrac19$} 
		{\True $\dfrac1{16}$} 
		 {$\dfrac1{18}$}
		{$\dfrac18$} 
	 \loigiai{
		Ta  có
		\[ 
		A=\dfrac{a^{\tfrac14}b^{\tfrac13}+b^{\tfrac14}a^{\tfrac13}} 
		 {a^{\tfrac1{12}}+b^{\tfrac1{12}}}
		=\dfrac{a^{\tfrac14}b^{\tfrac14}\left(b^{\tfrac1{12}}+a^{\tfrac1{12}}\right)} 
		 {a^{\tfrac1{12}}+b^{\tfrac1{12}}}
		=a^{\tfrac14}b^{\tfrac14}. 
		 \]
		Suy ra  $m n=\dfrac14\cdot\dfrac14=\dfrac1{16}$.
	} 
\end{ex}

% nguồn: Ok_Toan11_TracNghiem_Lan2
\begin{ex}%[1D6H1-2]%[ma:1b8229e2]
	 Cho $a$, $b$ là các số thực dương. Rút gọn $P=\dfrac{a^{\tfrac{4}{3}}b+ab^{\tfrac{4}{3}}}{\sqrt[3]{a}+\sqrt[3]{b}}$ ta được
	\choice 
		 {\True $P=ab$}
		 {$P=a+b$}
		{$P=a^4b+ab^4$} 
		{$P=ab\left(a+b\right)$} 
	 \loigiai{
	Ta  có $P=\dfrac{a^{\tfrac{4}{3}}b+ab^{\tfrac{4}{3}}}{\sqrt[3]{a}+\sqrt[3]{b}}=\dfrac{a\cdot a^{\tfrac{1}{3}}b+ab\cdot b^{\tfrac{1}{3}}}{a^{\tfrac{1}{3}}+b^{\tfrac{1}{3}}}=\dfrac{ab\left(a^{\tfrac{1}{3}}+b^{\tfrac{1}{3}}\right)}{a^{\tfrac{1}{3}}+b^{\tfrac{1}{3}}}=ab$.}
\end{ex}

% nguồn: Ok_Toan11_TracNghiem_Lan2
\begin{ex}%[1D6H1-2]%[ma:493ade79]
	Cho $a$  là một số thực dương. Giá trị của biểu thức $P=\left(\sqrt{2^a}\right)^{\tfrac{4}{a}}$ bằng
	\choice 
		 {\True $4$}
		{$2$} 
		 {$8$}
		 {$1$}
	\loigiai{ 
	 Ta có $P=\left(\sqrt{2^a}\right)^{\tfrac{4}{a}}=\left(2^{\tfrac{a}{2}}\right)^{\tfrac{4}{a}}=2^{^{\tfrac{4}{a}\cdot \tfrac{a}{2}}}=2^2=4$.}
\end{ex}

% nguồn: Ok_Toan11_TracNghiem_Lan2
\begin{ex}%[1D6H1-2]%[ma:f8c519a9]
	 Với $x\ge 0$ thì $\sqrt{x\sqrt{x\sqrt{x^2}}}$ bằng
	 \choice
		 {\True $x$}
		 {$x^2$}
		{$\sqrt{x}$} 
		 {$x^4$}
	 \loigiai{
	Ta có $\sqrt{x\sqrt{x\sqrt{x^2}}}=\sqrt{x\sqrt{x\cdot x}}=\sqrt{x\cdot  x}=x$.}
\end{ex}

% nguồn: Ok_Toan11_TracNghiem_Lan2
\begin{ex}%[1D6H1-2]%[ma:352b0b6a]
	Rút gọn biểu thức  $P=x^{\tfrac{2}{5}}\cdot \sqrt[6]{x}$ với $x>0$.
	\choice 
		 {$P=x^{\tfrac{1}{15}}$}
		{$P=x^{\tfrac{17}{15}}$} 
		{\True  $P=x^{\tfrac{17}{30}}$}
		{$P=\sqrt{x}$} 
	\loigiai{ 
	Ta có $P=x^{\tfrac{2}{5}}\cdot \sqrt[6]{x}=x^{\tfrac{2}{5}}\cdot x^{\tfrac{1}{6}}= x^{\tfrac{2}{5}+\tfrac{1}{6}}= x^{\tfrac{17}{30}}$.} 
\end{ex}

% nguồn: Ok_Toan11_TracNghiem_Lan2
\begin{ex}%[1D6H1-2]%[ma:979d0f76]
	Cho hai số thực dương $a$, $b$. Rút gọn biểu thức $A=\dfrac{a^{\tfrac{1}{4}}\sqrt[3]{b}+b^{\tfrac{1}{4}}\sqrt[3]{a}}{\sqrt[12]a+\sqrt[12]b}$ ta thu được $A=a^m\cdot b^n$. Tích của $m\cdot  n$ là
	\choice 
		 {$\dfrac{1}{9}$}
		{\True $\dfrac{1}{16}$} 
		{$\dfrac{1}{18}$} 
		{$\dfrac{1}{8}$} 
	\loigiai{ 
		Ta  có $A=\dfrac{a^{\tfrac{1}{4}}\sqrt[3]{b}+b^{\tfrac{1}{4}}\sqrt[3]{a}}{\sqrt[12]a+\sqrt[12]b}=\dfrac{a^{\tfrac{1}{4}}b^{\tfrac{1}{4}}\left(b^{\tfrac{1}{12}}+a^{\tfrac{1}{12}}\right)}{\left(a^{\tfrac{1}{12}}+b^{\tfrac{1}{12}}\right)}=a^{\tfrac{1}{4}}b^{\tfrac{1}{4}}$.\\
	Suy ra $m\cdot n=\dfrac{1}{4}\cdot  \dfrac{1}{4}=\dfrac{1}{16}$.}
\end{ex}

% nguồn: Ok_Toan11_TracNghiem_Lan2
\begin{ex}%[1D6H1-2]%[ma:f1cd35bd]
	Biết  biểu thức $P=\sqrt[6]{x^3\sqrt[3]{x^2\sqrt{x}}}, \left(x>0\right)$ được viết dưới dạng lũy thừa với số mũ hữu tỷ là $x^{\alpha}$. Khi đó, giá trị của $\alpha $ bằng
	 \choice
		 {$\dfrac{37}{15}$}
		 {\True $\dfrac{23}{36}$}
		{$\dfrac{23}{30}$} 
		 {$\dfrac{53}{30}$}
	\loigiai{Ta có  $P = \sqrt[6]{x^3 \sqrt[3]{x^2 \sqrt{x}}}
		= \sqrt[6]{x^3 \sqrt[3]{x^2 \cdot  x^{\tfrac{1}{2}}}}
		= \sqrt[6]{x^3  \cdot x^{\tfrac{5}{6}}}
		= x^{\tfrac{23}{36}} 
		$.\\ 
	Suy ra $\alpha  =\dfrac{23}{36}$.}
\end{ex}

% nguồn: Ok_Toan11_TracNghiem_Lan2
\begin{ex}%[1D6H1-2]%[ma:9a751cd1]
	[Mã 102 2017] Cho biểu  thức $P=\sqrt[4]{x\cdot \sqrt[3]{x^2\cdot \sqrt{x^3}}}$, với $x>0$. Mệnh đề nào dưới đây \textbf{đúng}?
	 \choice
		 {$P=x^{\frac{2}{3}}$}
		{$P=x^{\frac{1}{2}}$} 
		{\True  $P=x^{\frac{13}{24}}$}
		{$P=x^{\frac{1}{4}}$} 
	\loigiai{ 
		Với $x>0$,  ta có
		\begin{eqnarray*}
			 P&=&\sqrt[4]{x\cdot \sqrt[3]{x^2\cdot \sqrt{x^3}}}
			=\sqrt[4]{x\cdot \sqrt[3]{x^2\cdot{x^{\frac{3}{2}}}}}  \\
			 &=&\sqrt[4]{x\cdot \sqrt[3]{x^{\frac{7}{2}}}}
			 =\sqrt[4]{x\cdot{x^{\frac{7}{6}}}}
			 =\sqrt[4]{x^{\frac{13}{6}}}=x^{\frac{13}{24}}.
		\end{eqnarray*}
	} 
\end{ex}

% nguồn: Ok_Toan11_TracNghiem_Lan2
\begin{ex}%[1D6H1-2]%[ma:323da31c]
	Cho $x$ là số thực dương. Biết  $\sqrt{x\cdot \sqrt[3]{x\sqrt{x\sqrt[3]{x}}}}=x^{\tfrac{b}{a}}$ với $a$, $b$ là các số tự nhiên và $\dfrac{a}{b}$ là phân số tối giản. Tính $a+b$.
	 \choice
		{\True $16$} 
		{$15$} 
		 {$14$}
		 {$17$}
	 \loigiai{
		Ta có $\sqrt{x\cdot \sqrt[3]{x\sqrt{x\sqrt[3]{x}}}}=\sqrt{x\cdot \sqrt[3]{x\sqrt{x\cdot x^{\tfrac{1}{3}}}}}=\sqrt{x\cdot x^{\tfrac{1}{3}}\cdot x^{\tfrac{2}{9}}}=\sqrt{x\cdot  x^{\tfrac{5}{9}}}=x^{\tfrac{7}{9}}$.\\
	Khi đó  $a=7$; $b=7$ nên $a+b=16$.}
\end{ex}

% nguồn: Ok_Toan11_TracNghiem_Lan2
\begin{ex}%[1D6H1-2]%[ma:d177179b]
	Cho hai số thực dương $a$, $b$. Rút gọn biểu thức  $A=\dfrac{a^{\tfrac{1}{3}}\sqrt{b}+b^{\tfrac{1}{3}}\sqrt{a}}{\sqrt[6]a+\sqrt[6]b}$ ta thu được $A=a^m\cdot b^n$. Tích của $m\cdot n$ là
	\choice 
		{$\dfrac{1}{8}$} 
		 {$\dfrac{1}{21}$}
		 {\True $\dfrac{1}{9}$}
		{$\dfrac{1}{18}$} 
	\loigiai{Ta có 
		\allowdisplaybreaks 
		\begin{eqnarray*}
			A&=&\dfrac{a^{\tfrac{1}{3}}\sqrt{b}+b^{\tfrac{1}{3}}\sqrt{a}}{\sqrt[6]a+\sqrt[6]b}=\dfrac{a^{\tfrac{1}{3}}\cdot b^{\tfrac{1}{2}}+b^{\tfrac{1}{3}}\cdot a^{\tfrac{1}{2}}}{a^{\tfrac{1}{6}}+b^{\tfrac{1}{6}}}\\ 
			&=&\dfrac{a^{\tfrac{1}{3}}\cdot b^{\tfrac{1}{3}}\left(b^{\tfrac{1}{6}}+a^{\tfrac{1}{6}}\right)}{a^{\tfrac{1}{6}}+b^{\tfrac{1}{6}}}\\ 
			&=&a^{\tfrac{1}{3}}\cdot b^{\tfrac{1}{3}}. 
		\end{eqnarray*}
	Do đó $m=\dfrac{1}{3}$, $n=\dfrac{1}{3}\Rightarrow m\cdot n=\dfrac{1}{9}$.} 
\end{ex}

% nguồn: Ok_Toan11_TracNghiem_Lan2
\begin{ex}%[1D6H1-2]%[ma:465fffc8]
	 Đơn giản biểu thức $P=a^{\sqrt{2}}\cdot \left(\dfrac{1}{a}\right)^{\sqrt{2}-1}$ với $a>0$, được kết quả là
	\choice 
		{$a^{\sqrt{2}}$} 
		{$a^{2\sqrt{2}-1}$} 
		{$a^{1-\sqrt{2}}$} 
		{\True  $a$}
	 \loigiai{
		Ta có 
		\allowdisplaybreaks 
		\begin{eqnarray*}
			P&=&a^{\sqrt{2}}\cdot \left(\dfrac{1}{a}\right)^{\sqrt{2}-1}=a^{\sqrt{2}}\cdot \left(a^{-1}\right)^{\sqrt{2}-1}\\ 
			 &=&a^{\sqrt{2}}\cdot a^{-\sqrt{2}+1}=a^{\sqrt{2}-\sqrt{2}+1}=a^1=a.
		\end{eqnarray*}
	Vậy $P=a$.} 
\end{ex}

% nguồn: Ok_Toan11_TracNghiem_Lan2
\begin{ex}%[1D6H1-2]%[ma:dd8b8a0a]
	Biểu thức $\sqrt{x}\cdot \sqrt[3]{x}\cdot \sqrt[6]{x^5}$, $(x>0)$ viết dưới dạng luỹ thừa với số mũ  hữu tỷ là
	\choice 
		{$x^{\tfrac{2}{3}}$} 
		 {$x^{\tfrac{5}{2}}$}
		 {$x^{\tfrac{7}{3}}$}
		{\True $x^{\tfrac{5}{3}}$} 
	\loigiai{ 
	 Ta có $\sqrt{x}\cdot \sqrt[3]{x}\cdot \sqrt[6]x^5 =x^{\tfrac{1}{2}}\cdot x^{\tfrac{1}{3}}\cdot x^{\tfrac{5}{6}}=x^{\tfrac{10}{6}}=x^{\tfrac{5}{3}}$.}
\end{ex}

% nguồn: Ok_Toan11_TracNghiem_Lan2
\begin{ex}%[1D6H1-2]%[ma:ae2d885d]
	Cho biểu thức $P=x^{-\tfrac{3}{4}}\cdot \sqrt{\sqrt{x^5}}$,  $x>0$. Khẳng định nào sau đây là đúng?
	\choice 
		{$P=x^{-2}$} 
		{$P=x^{-\tfrac{1}{2}}$} 
		{\True  $P=x^{\tfrac{1}{2}}$}
		{$P=x^2$} 
	 \loigiai{
	Ta  có $P=x^{-\tfrac{3}{4}}\cdot \sqrt{\sqrt{x^5}}=x^{-\tfrac{3}{4}}\cdot x^{\tfrac{5}{4}}=x^{-\tfrac{3}{4}+\tfrac{5}{4}}=x^{\tfrac{1}{2}}$.}
\end{ex}

% nguồn: Ok_Toan11_TracNghiem_Lan2
\begin{ex}%[1D6H1-2]%[ma:21ceb1e3]
	Rút gọn biểu thức $A=\dfrac{\sqrt[3]{a^7}\cdot a^{\tfrac{11}{3}}}{a^4\cdot \sqrt[7]{a^{-5}}}$ với $a>0$ ta được  kết quả $A=a^{\tfrac{m}{n}}$ trong đó $m$, $n\in{N^{*}}$ và $\dfrac{m}{n}$ là phân số tối giản. Khẳng định nào sau đây đúng?
	 \choice
		 {\True $m^2-n^2=312$}
		{$m^2+n^2=543$} 
		 {$m^2-n^2=-312$}
		 {$m^2+n^2=409$}
	\loigiai{ 
		Ta có $A=\dfrac{\sqrt[3]{a^7}\cdot a^{\tfrac{11}{3}}}{a^4\cdot  \sqrt[7]{a^{-5}}}=\dfrac{a^{\tfrac{7}{3}}\cdot a^{\tfrac{11}{3}}}{a^4\cdot a^{\tfrac{-5}{7}}}=\dfrac{a^6}{a^{\tfrac{23}{7}}}=a^{\tfrac{19}{7}}$.\\
		Mà $A=a^{\tfrac{m}{n}}$, $m$, $n\in{N^{*}}$ và  $\dfrac{m}{n}$ là phân số tối giản nên $m=19$, $n=7$.\\
	Do đó $m^2-n^2=312$.} 
\end{ex}

% nguồn: Ok_Toan11_TracNghiem_Lan2
\begin{ex}%[1D6H1-2]%[ma:e28469b3]
	Rút gọn biểu thức $A=\dfrac{\sqrt[3]{a^5}\cdot a^{\tfrac{7}{3}}}{a^4\cdot \sqrt[7]{a^{-2}}}$ với $a>0$ ta được kết quả $A=a^{\tfrac{m}{n}}$, trong đó $m$, $n\in{\mathbb{N}^{*}}$ và $\dfrac{m}{n}$ là phân số tối giản. Khẳng định  nào sau đây đúng?
	 \choice
		{$3m^2-2n=2$} 
		{$m^2+n^2=43$} 
		{\True $2m^2+n=15$} 
		 {$m^2+n^2=25$}
	 \loigiai{
		Ta  có
		 \allowdisplaybreaks
		\begin{eqnarray*}
			A&=&\dfrac{\sqrt[3]{a^5}\cdot a^{\tfrac{7}{3}}}{a^4\cdot \sqrt[7]{a^{-2}}}=\dfrac{a^{\tfrac{5}{3}}\cdot a^{\tfrac{7}{3}}}{a^4\cdot  a^{\tfrac{-2}{7}}}\\
			 &=&\dfrac{a^{\tfrac{5}{3}+\tfrac{7}{3}}}{a^{4-\tfrac{2}{7}}}=\dfrac{a^4}{a^{\tfrac{26}{7}}}\\
			&=&a^{4-\tfrac{26}{7}}=a^{\tfrac{2}{7}}. 
		\end{eqnarray*}
	Suy ra $m=2;n=7\Rightarrow  2m^2+n=15$.}
\end{ex}

% nguồn: Ok_Toan11_TracNghiem_Lan2
\begin{ex}%[1D6H1-2]%[ma:9f172d61]
	Rút gọn biểu thức $A=\dfrac{\sqrt[3]{a^5}\cdot a^{\tfrac73}}{a^{4}\cdot\sqrt[7]{a^{-2}}}$ với $a>0$ ta được kết quả  $A=a^{\tfrac{m}{n}}$, trong đó $m$, $n\in\mathbb{N}^{*}$ và $\tfrac{m}{n}$ là phân số tối giản. Khẳng định nào sau đây đúng?
	 \choice
		 {$3m^2-2n=2$}
		 {$m^2+n^2=43$}
		 {\True $2m^2+n=15$}
		{$m^2+n^2=25$} 
	 \loigiai{
		Ta có 
		\[ 
		 A=\dfrac{\sqrt[3]{a^5}\cdot a^{\tfrac73}}{a^{4}\cdot\sqrt[7]{a^{-2}}}
		=\dfrac{a^{\tfrac53}\cdot  a^{\tfrac73}}{a^{4-\tfrac27}}
		=a^{\tfrac53+\tfrac73}\cdot  a^{-\left(4-\tfrac27\right)}
		=a^{4-\tfrac{26}{7}}=a^{\tfrac{2}{7}}. 
		 \]
		Suy ra $m=2$, $n=7$. Do đó $2m^2+n=2\cdot  2^2+7=15$.
	 }
\end{ex}

% nguồn: Ok_Toan11_TracNghiem_Lan2
\begin{ex}%[1D6H1-2]%[ma:7e18bbc7]
	Cho $a$ là số thực dương. Đơn giản  biểu thức $P=\dfrac{a^{\tfrac{4}{3}}\left(a^{\tfrac{-1}{3}}+a^{\tfrac{2}{3}}\right)}{a^{\tfrac{1}{4}}\left(a^{\tfrac{3}{4}}+a^{\tfrac{-1}{4}}\right)}$.
	 \choice
		 {$P=a\left(a+1\right)$}
		{$P=a-1$} 
		{\True $P=a$} 
		{$P=a+1$} 
	 \loigiai{
		Ta  có
		 \allowdisplaybreaks
		\begin{eqnarray*}
			 P&=&\dfrac{a^{\tfrac{4}{3}}\left(a^{\tfrac{-1}{3}}+a^{\tfrac{2}{3}}\right)}{a^{\tfrac{1}{4}}\left(a^{\tfrac{3}{4}}+a^{\tfrac{-1}{4}}\right)}=\dfrac{a^{\tfrac{4}{3}}\cdot a^{\tfrac{-1}{3}}+a^{\tfrac{4}{3}}\cdot a^{\tfrac{2}{3}}}{a^{\tfrac{1}{4}}\cdot a^{\tfrac{3}{4}}+a^{\tfrac{1}{4}}\cdot a^{\tfrac{-1}{4}}}\\
			&=&\dfrac{a+a^2}{a+1}=\dfrac{a\left(a+1\right)}{a+1}=a. 
	\end{eqnarray*}}
\end{ex}

% nguồn: Ok_Toan11_TracNghiem_Lan2
\begin{ex}%[1D6H1-2]%[ma:8b4ae244]
	 [Đề Tham Khảo 2017] Tính giá trị của biểu thức $P=(7+4\sqrt{3})^{2017}(4\sqrt{3}-7)^{2016}$.
	 \choice
		{$P=(7+4\sqrt{3})^{2016}$} 
		 {$P=1$}
		{$P=7-4\sqrt{3}$} 
		 {\True $P=7+4\sqrt{3}$}
	\loigiai{ 
		Ta có $\begin{aligned}
		P&=(7+4\sqrt{3})^{2017}(4\sqrt{3}-7)^{2016} \\ 
		 &=(7+4\sqrt{3})\cdot{\left[(7+4\sqrt{3})(4\sqrt{3}-7)\right]^{2016}} \\
		&  =(7+4\sqrt{3})(-1)^{2016}=7+4\sqrt{3}.
	\end{aligned}$
 }
\end{ex}

% nguồn: Ok_Toan11_TracNghiem_Lan2
\begin{ex}%[1D6H1-2]%[ma:442b6909]
	Cho biểu thức $P=x^{\tfrac{1}{2}}\cdot x^{\tfrac{1}{3}}\cdot \sqrt[6]{x}$ với $x>0$. Mệnh  đề nào dưới đây đúng?
	\choice 
		{\True $P=x$} 
		 {$P=x^{\tfrac{11}{6}}$}
		 {$P=x^{\tfrac{7}{6}}$}
		 {$P=x^{\tfrac{5}{6}}$}
	\loigiai{ 
	Ta có  $P=x^{\tfrac{1}{2}}\cdot x^{\tfrac{1}{3}}\cdot \sqrt[6]{x}=x^{\tfrac{1}{2}+\tfrac{1}{3}+\tfrac{1}{6}}=x$.}
\end{ex}

% nguồn: Ok_Toan11_TracNghiem_Lan2
\begin{ex}%[1D6H1-2]%[ma:d59bdb70]
	Viết biểu thức $P=\sqrt[3]{x\cdot \sqrt[4]x}$ ($x>0$) dưới  dạng luỹ thừa với số mũ hữu tỷ.
	\choice 
		{$P=x^{\tfrac{5}{4}}$} 
		{\True $P=x^{\tfrac{5}{12}}$} 
		{$P=x^{\tfrac{1}{7}}$} 
		{$P=x^{\tfrac{1}{12}}$} 
	\loigiai{ 
	Ta có $P=\left(x\cdot x^{\tfrac{1}{4}}\right)^{\tfrac{1}{3}}=\left(x^{\tfrac{5}{4}}\right)^{\tfrac{1}{3}}=x^{\tfrac{5}{12}}$.} 
\end{ex}

% nguồn: Ok_Toan11_TracNghiem_Lan2
\begin{ex}%[1D6H1-2]%[ma:efdfbc56]
	Biết biểu thức $P=\sqrt[6]{x^{3}\sqrt[3]{x^2\sqrt{x}}}$, $(x>0)$ được viết dưới dạng lũy thừa với số mũ hữu tỷ là  $x^{\alpha}$. Khi đó, giá trị của $\alpha$ bằng
	 \choice
		{$\dfrac{37}{15}$} 
		 {\True $\dfrac{23}{36}$}
		{$\dfrac{23}{30}$} 
		 {$\dfrac{53}{30}$}
	\loigiai{ 
		Ta có 
		 \[
		 P=\sqrt[6]{x^{3}\sqrt[3]{x^2\sqrt{x}}}
		 =\sqrt[6]{x^{3}\sqrt[3]{x^2x^{\tfrac{1}{2}}}}
		=\sqrt[6]{x^{3}\sqrt[3]{x^{\tfrac{5}{2}}}} 
		 =\sqrt[6]{x^3x^{\tfrac{5}{6}}}
		 =\sqrt[6]{x^{\tfrac{23}{6}}}
		=x^{\tfrac{23}{36}}. 
		\] 
	 }
\end{ex}

% nguồn: Ok_Toan11_TracNghiem_Lan2
\begin{ex}%[1D6H1-2]%[ma:cb722fb4]
	Cho $a$ là số thực dương. Viết và rút gọn biểu thức $a^{\tfrac{3}{2\,018}}\cdot \sqrt[2\,018]{a}$ dưới dạng lũy thừa  với số mũ hữu tỉ. Tìm số mũ của biểu thức rút gọn đó.
	 \choice
		 {\True $\dfrac{2}{1\,009}$}
		 {$\dfrac{1}{1\,009}$}
		 {$\dfrac{3}{1\,009}$}
		 {$\dfrac{3}{2\,018^2}$}
	 \loigiai{
		Ta có $a^{\tfrac{3}{2\,018}}\cdot  \sqrt[2018]a=a^{\tfrac{3}{2\,018}}\cdot a^{\tfrac{1}{2\,018}}=a^{\tfrac{4}{2\,018}}=a^{\tfrac{2}{1\,009}}$.\\
	Vậy số mũ  của biểu thức rút gọn bằng $\dfrac{2}{1\,009}$.}
\end{ex}

% nguồn: Ok_Toan11_TracNghiem_Lan2
\begin{ex}%[1D6H1-2]%[ma:e3e950be]
	Cho số thực $a$ dương tùy ý. Đặt $a^{\tfrac{5}{4}}\cdot  \sqrt{a\sqrt[3]{a}}=a^p$. Khẳng định đúng là
	\choice 
		{$p=\dfrac{19}{12}$} 
		 {\True $p=\dfrac{23}{12}$}
		 {$p=\dfrac{13}{12}$}
		{$p=\dfrac{23}{24}$} 
	\loigiai{ 
		Ta có $a^{\tfrac{5}{4}}\cdot \sqrt{a\sqrt[3]{a}}=a^{\tfrac{5}{4}}\cdot  a^{\tfrac{1}{2}}\cdot a^{\tfrac{1}{2\cdot 3}}=a^{\tfrac{5}{4}+\tfrac{1}{2}+\tfrac{1}{6}}=a^{\tfrac{23}{12}}$.\\
	Suy  ra $p=\dfrac{23}{12}$.}
\end{ex}

% nguồn: Ok_Toan11_TracNghiem_Lan2
\begin{ex}%[1D6H1-2]%[ma:df5dfb26]
	Rút gọn biểu thức $Q=b^{\tfrac{5}{3}}:\sqrt[3]{b}$ với  $b>0$.
	\choice 
		 {$Q=b^{-\tfrac{4}{3}}$}
		{\True $Q=b^{\tfrac{4}{3}}$} 
		 {$Q=b^{\tfrac{5}{9}}$}
		 {$Q=b^2$}
	 \loigiai{
	Ta có  $Q=b^{\tfrac{5}{3}}:\sqrt[3]{b}=b^{\tfrac{5}{3}}:b^{\tfrac{1}{3}}=b^{\tfrac{4}{3}}$.}
\end{ex}

% nguồn: Ok_TracNghiemLop11_Lan1
\begin{ex}%[1D6H1-2]%[ma:fa08cd8f]
	Cho biểu thức $P=\sqrt x \cdot \sqrt[3]{x}\cdot\sqrt[6]{x^4}$, $\left(x > 0\right)$. Mệnh đề nào dưới  đây đúng?
	 \choice
		 {$P=x^{\tfrac{2}{3}}$}
		 {$P=x^{\tfrac{5}{2}}$}
		{\True $P=x^{\tfrac{3}{2}}$} 
		{$P=x^{\tfrac{2}{5}}$} 
	 \loigiai{
		Ta có $P=\sqrt x  \cdot \sqrt[3]{x}\cdot\sqrt[6]{x^4} = x^{\tfrac{1}{2}}\cdot x^{\tfrac{1}{3}}\cdot \tfrac{2}{3}=x^{\tfrac{3}{2}}$.
	 }
\end{ex}

% nguồn: Ok_TracNghiemLop11_Lan1
\begin{ex}%[1D6H1-2]%[ma:f21f6576]
	Cho biểu thức $P=x^{-\tfrac{3}{4}}\cdot\sqrt{\sqrt{x^5}}$ với $x>0$. Khẳng định nào sau đây là  đúng?
	\choice 
		{$x^{-2}$} 
		 {$x^2$}
		 {$x^{-\tfrac{1}{2}}$}
		{\True  $x^{\tfrac{1}{2}}$}
	 \loigiai{
	Ta có $P=x^{-\tfrac{3}{4}}\cdot x^{\tfrac{5}{4}}=x^{-\tfrac{3}{4}+\tfrac{5}{4}}=x^{\tfrac{1}{2}}$.} 
\end{ex}

% nguồn: Ok_TracNghiemLop11_Lan1
\begin{ex}%[1D6H1-2]%[ma:db892e80]
	Cho số thực dương $ a>0$ và khác $1$. Rút gọn biểu thức $P=\dfrac{a^{\frac{1}{3}}\left(a^{\frac{1}{2}} - a^{\frac{5}{2}}\right)}{a^{\frac{1}{4}}\left(a^{\frac{7}{12}} - a^{\frac{19}{12}}\right)}$ có kết quả dạng $ m+{a^n}$ với $ m$, $ n$ là 2 số nguyên. Tổng của $ m$ và $  n$ bằng
	 \choice
		 {$3$}
		{\True $2$} 
		 {$4$}
		 {$ 5$}
	\loigiai{ 
		Ta  có:
		\begin{eqnarray*}
			P&=&\dfrac{a^{\frac{1}{3}}\left(a^{\frac{1}{2}}-a^{\frac{5}{2}}\right)}{a^{\frac{1}{4}}\left(a^{\frac{7}{12}}-a^{\frac{19}{12}}\right)}\\ 
			&=&\dfrac{a^{\frac{1}{3}}\cdot a^{\frac{1}{2}}\left(1-a^2\right)}{a^{\frac{1}{4}}\cdot a^{\frac{7}{12}}\left(1-a\right)}\\ 
			 &=&\dfrac{a^{\frac{5}{6}}\left(1+a\right)}{a^{\frac{5}{6}}}=1+a\\
			\Rightarrow  m&=&1,n=1\\
			 \Rightarrow m+n&=&2
		\end{eqnarray*}
	 }
\end{ex}

% nguồn: Ok_TracNghiemLop11_Lan1
\begin{ex}%[1D6H1-2]%[ma:f3174303]
	Viết biểu thức $P=\dfrac{a^2a^{\tfrac{5}{2}}\sqrt[3]{a^4}}{\sqrt[6]{a^7}}$, $\left(a > 0\right)$ dưới dạng lũy thừa với số  mũ hữu tỉ.
	\choice 
		{$P=a^5$} 
		{\True $P=a^{\tfrac{14}{3}}$} 
		{$P=a^4$} 
		 {$P=a^{\tfrac{13}{3}}$}
	\loigiai{ 
		Ta có $a^2a^{\tfrac{5}{2}}\sqrt[3]{a^4}=a^2a^{\tfrac{5}{2}}a^{\tfrac{4}{3}} =  a^{\tfrac{35}{6}}$ và $\sqrt[6]{a^7}=a^{\tfrac{7}{6}}$.\\
		 Do đó $P = a^{\tfrac{14}{3}}$.
	} 
\end{ex}

% nguồn: Ok_TracNghiemLop11_Lan1
\begin{ex}%[1D6H1-2]%[ma:b3e866b0]
	Viết biểu thức $P=\sqrt[3]{x\cdot\sqrt[4]{x}}$, $(x>0)$ dưới dạng  lũy thừa với số mũ hữu tỷ.
	\choice 
		{$P=x^{\tfrac{5}{4}}$} 
		{$P=x^{\tfrac{1}{12}}$} 
		 {\True $P=x^{\tfrac{5}{12}}$}
		{$P=x^{\tfrac{1}{7}}$} 
	 \loigiai{
	Ta có $P=\sqrt[3]{x\cdot\sqrt[4]{x}}=\sqrt[3]{x\cdot x^{\tfrac{1}{4}}}=\sqrt[3]{x^{\tfrac{5}{4}}}=x^{\tfrac{5}{12}}$.} 
\end{ex}

% nguồn: Ok_TracNghiemLop11_Lan1
\begin{ex}%[1D6H1-2]%[ma:c8ab8057]
	Với $\alpha$ là số thực bất kì, mệnh đề nào sau đây sai? 
	 \choice
		 {$\sqrt{10^\alpha}=(\sqrt{10})^\alpha$}
		 {$\sqrt{10^\alpha}=10^{\tfrac{\alpha}{2}}$}
		 {$\left(10^\alpha\right)^2=(100)^\alpha$}
		{\True $\left(10^\alpha\right)^2=(10)^{\alpha^2}$} 
	 \loigiai{
		Với $\alpha$ là số thực bất kì thì $\left(10^\alpha\right)^2=(10)^{\alpha \cdot 2} \ne (10)^{\alpha^2}$. 
	} 
\end{ex}

% nguồn: Ok_TracNghiemLop11_Lan1
\begin{ex}%[1D6H1-2]%[ma:a3742df0]
	 Rút gọn biểu thức
	$P=\dfrac{a^{\sqrt 5 + 1} \cdot a^{2 - \sqrt 5}}{\left(a^{\sqrt 2 - 2}\right)^{\sqrt 2 + 2}}$ với $ a$ là số thực  dương khác $1$.
	 \choice
		{\True $a^5$} 
		{$ a$} 
		{$a^3$} 
		{$a^4$} 
	 \loigiai{
		Ta có $P=\dfrac{a^{\sqrt 5+1}\cdot a^{2-\sqrt 5}}{\left(a^{\sqrt 2-2}\right)^{\sqrt  2+2}}=\dfrac{a^{\sqrt 5+1+2-\sqrt 5}}{a^{\left(\sqrt 2-2\right)\left(\sqrt 2+2\right)}}=\dfrac{a^3}{a^{-2}}=a^5$.
	 }
\end{ex}

% nguồn: Ok_TracNghiemLop11_Lan1
\begin{ex}%[1D6H1-2]%[ma:620e55d4]
	Cho hai số thực dương $ x$ và $ y$. Rút gọn  biểu thức
	$A=\left(\dfrac{x^{\sqrt 5}}{y^{\sqrt 5 - 2}}\right)^{\sqrt 5 + 2}\cdot \left(\dfrac{x^{- 2\sqrt 5 - 2}}{y^{- 3}}\right)$ được kết quả có  dạng ${x^a}{y^b}$ với $ a$, $ b$ là 2 số nguyên. Khẳng định nào đúng?
	\choice 
		 {\True $5b+a=13$}
		{$5a+b=12$} 
		 {$ a+b=3$}
		{$2a+b=1$} 
	\loigiai{ 
		Với $ x>0;y>0$, ta có 
		\begin{eqnarray*}
			A&=&\left(\dfrac{x^{\sqrt 5}}{y^{\sqrt 5-2}}\right)^{\sqrt 5+2}\cdot \left(\dfrac{x^{-2\sqrt  5-2}}{y^{-3}}\right)\\
			&=&\left(\dfrac{x^{\sqrt 5\left(\sqrt 5+2\right)}}{y^{\left[\left(\sqrt 5-2\right)\cdot  \left(\sqrt 5+2\right) \right]}}\right)\cdot \left(\dfrac{x^{-2\sqrt 5-2}}{y^{-3}}\right)\\
			&=&\left(\dfrac{x^{5+2\sqrt 5}}{y}\right)\cdot  \left(\dfrac{x^{-2\sqrt 5-2}}{y^{-3}}\right)\\
			&=&\dfrac{x^{5+2\sqrt  5-2\sqrt 5-2}}{y^{1-3}}\\
			&=&x^3{y^2}  \\
			\Rightarrow a&=&3, b=2 
		\end{eqnarray*}
		 Vậy, $5b+a=5\cdot 2+3=13$.
	 }
\end{ex}

% nguồn: Ok_TracNghiemLop11_Lan1
\begin{ex}%[1D6H1-2]%[ma:f349bd67]
	 Cho $x$ là số thực dương, viết biểu thức $P=x^{\tfrac{1}{3}}\cdot \sqrt[4]{x}$ dưới dạng lũy thừa với số mũ hữu tỉ
	 \choice
		{$P=x^{\tfrac{1}{12}}$} 
		{\True $P=x^{\tfrac{7}{12}}$} 
		{$P=x^{\tfrac{2}{7}}$} 
		 {$P=x^{\tfrac{2}{3}}$.}
	\loigiai{ 
		Ta có  $P=x^{\tfrac{1}{3}}\cdot \sqrt[4]{x}=x^{\tfrac{1}{3}}\cdot x^{\tfrac{1}{4}}=x^{\tfrac{7}{12}}$.
	} 
\end{ex}

% nguồn: Ok_TracNghiemLop11_Lan1
\begin{ex}%[1D6H1-2]%[ma:1d78d177]
	 Cho biểu thức $P=\dfrac{b^{\sqrt 3+1}\cdot b^{2-\sqrt 3}}{\left(b^{\sqrt 2-2} \right)^{\sqrt 2+2}}$, với $b>0$. Khẳng định nào sau đây {\bf đúng}?
	 \choice
		 {$P=b^4$}
		{\True  $P=b^5$}
		 {$P=b^3$}
		 {$P=b$}
	\loigiai{ 
		Ta có $P=\dfrac{b^{\sqrt 3+1}\cdot b^{2-\sqrt  3}}{\left(b^{\sqrt 2-2} \right)^{\sqrt 2+2}}=\dfrac{b^{\sqrt 3+1+2-\sqrt 3}}{b^{(\sqrt 2-2)(\sqrt 2+2)}}=\dfrac{b^3}{b^{-2}}=b^3\cdot b^2=b^5$.
	} 
\end{ex}

% nguồn: Ok_TracNghiemLop11_Lan1
\begin{ex}%[1D6H1-2]%[ma:d5776bdf]
	Cho $x  > 0$. Viết biểu thức $\sqrt[5]{x^7}:\sqrt{x}$ dưới dạng lũy thừa với số mũ hữu tỉ.
	\choice 
		 {$x^{\tfrac{9}{7}}$}
		{$x^{\tfrac{9}{5}}$} 
		 {$x^{\tfrac{7}{10}}$}
		 {\True $x^{\tfrac{9}{10}}$}
	 \loigiai{
		Ta có  $\sqrt[5]{x^7}:\sqrt{x}=x^{\tfrac{7}{5}}:x^{\tfrac{1}{2}}=x^{\tfrac{9}{10}}$.
	 }
\end{ex}

% nguồn: Ok_TracNghiemLop11_Lan1
\begin{ex}%[1D6H1-2]%[ma:7a50a335]
	Rút gọn biểu thức $A=\dfrac{\sqrt[3]{a^7}\cdot a^{\frac{11}{3}}}{a^4\cdot\sqrt[7]{a^{-5}}}$, với $a>0$ ta được kết  quả $A=a^{\frac{m}{n}}$ trong đó $m,n\in\mathbb{N}^*$ và $\dfrac{m}{n}$ là phân số tối giản. Khẳng định nào sau đây \textbf{đúng}?
	 \choice
		{\True$m^2-n^2=312$} 
		{$m^2+n^2=543$} 
		{$m^2-n^2=-312$} 
		{$m^2+n^2=409$} 
	 \loigiai{
		Ta  có $A=\dfrac{\sqrt[3]{a^7}\cdot a^{\frac{11}{3}}}{a^4\cdot\sqrt[7]{a^{-5}}} = \dfrac{a^{\frac{7}{3}}\cdot a^{\frac{11}{3}}}{a^4\cdot a^{\frac{-5}{7}}} = \dfrac{a^{\frac{18}{3}}}{a^{\frac{23}{7}}}=a^{\frac{18}{3}-\frac{23}{7}}=a^{\frac{19}{7}}$.\\
		 Khẳng định đúng là $m^2-n^2=19^2-7^2=312$.
	} 
\end{ex}

% nguồn: Ok_TracNghiemLop11_Lan1
\begin{ex}%[1D6H1-2]%[ma:a034264b]
	Biểu thức $P=\sqrt[{3}]{x\sqrt[{5}]{x^2  \sqrt{x}}}=x^{\alpha}$ (với $x > 0$), giá trị của $\alpha$ là
	\choice 
		{\True  $\dfrac{1}{2}$}
		 {$\dfrac{5}{2}$}
		 {$\dfrac{9}{2}$}
		{$\dfrac{3}{2}$} 
	\loigiai{ 
		$P=\sqrt[{3}]{x\sqrt[{5}]{x^2  \sqrt{x}}}=\sqrt[{3}]{x\sqrt[{5}]{x^2\cdot x^{\tfrac{1}{2}}}}=\sqrt[{3}]{x\cdot\left(x^{\tfrac{5}{2}} \right)^{\tfrac{1}{5}}}=\left(x^{\tfrac{3}{2}} \right)^{\tfrac{1}{3}}=x^{\tfrac{1}{2}} \Rightarrow \alpha=\dfrac{1}{2}$.
	} 
\end{ex}

% nguồn: Ok_TracNghiemLop11_Lan1
\begin{ex}%[1D6H1-2]%[ma:edb548bc]
	Cho $a$ là số thực dương khác $1$. Giá trị của  biểu thức $P=a^{\tfrac{2}{3}} \sqrt{a}$ bằng
	 \choice
		 {$a^{\tfrac{2}{3}}$}
		{$a^{\tfrac{5}{6}}$} 
		{$a^3$} 
		{\True $a^{\tfrac{7}{6}}$} 
	\loigiai{ 
	Ta có  $P=a^{\tfrac{2}{3}} \sqrt{a}=a^{\tfrac{2}{3}} \cdot a^{\tfrac{1}{2}}=a^{\tfrac{7}{6}}$.
	} 
\end{ex}

% nguồn: Ok_TracNghiemLop11_Lan1
\begin{ex}%[1D6H1-2]%[ma:3e1ef6f2]
	Biểu thức $K=\sqrt[3]{\dfrac{2}{3}\sqrt[3]{\dfrac{2}{3}\sqrt{\dfrac{2}{3}}}}$ viết dưới dạng luỹ thừa với số mũ hữu  tỉ là
	 \choice
		{$\left(\dfrac{2}{3}\right)^{\tfrac{5}{18}}$} 
		{\True  $\left(\dfrac{2}{3}\right)^{\tfrac{1}{2}}$}
		{$\left(\dfrac{2}{3}\right)^{\tfrac{1}{8}}$} 
		{$\left(\dfrac{2}{3}\right)^{\tfrac{1}{6}}$} 
	 \loigiai{
		Coi $X=\dfrac{2}{3}$. Theo nguyên tắc \lq\lq Chia cộng\rq\rq \,ta  có\\
		$K=\sqrt[3]{X\sqrt[3]{X\sqrt X}}=\sqrt[3]{X\sqrt[3]{X.X^{\tfrac{1}{2}}}}=\sqrt[3]{X\sqrt[3]{X^{\tfrac{3}{2}}}}=\sqrt[3]{X.X^{\tfrac{1}{2}}}=\sqrt[3]{X^{\tfrac{3}{2}}}=X^{\tfrac{1}{2}}$. 
	} 
\end{ex}

% nguồn: Ok_TracNghiemLop11_Lan1
\begin{ex}%[1D6H1-2]%[ma:4af970d8]
	Viết biểu  thức $Q=b^{\frac{7}{3}}:\sqrt[3]{b}(b>0)$ dưới dạng một lũy thừa có mũ hữu ti.
	 \choice
		{$Q=b^{\frac{2}{3}}$} 
		{\True  $Q=b^2$}
		{$Q=b^{\frac{16}{3}}$} 
		{$Q=b^{\frac{8}{3}}$} 
	 \loigiai{
	 $Q=b^{\frac{7}{3}}:\sqrt[3]{b}=b^{\frac{7}{3}}:b^{\frac{1}{3}}=b^{\frac{7}{3}-\frac{1}{3}}=b^2 $.
	 }
\end{ex}

% nguồn: Ok_TracNghiemLop11_Lan1
\begin{ex}%[1D6H1-2]%[ma:28b8baae]
	Biết $4^a+4^{-a}=5^{1000}$.  Lúc đó giá trị của $A=2^a+2^{-a}$ là
	 \choice
		 {$A=\sqrt{5^{1000}-2}$}
		{\True $A=\sqrt{5^{1000}+2}$} 
		{$A=5^{1000}+2$} 
		 {$A=5^{1000}-2$}
	 \loigiai{Ta có $\left( 2^a+2^{-a}\right)^2=(2^a)^2+2+(2^{-a})^2=4^a+2+4^{-a}=5^{1000}+2$.\\
	 $\Rightarrow A=2^a+2^{-a}=\sqrt{5^{1000}+2}$.}
\end{ex}

% nguồn: Ok_TracNghiemLop11_Lan1
\begin{ex}%[1D6H1-2]%[ma:5643b93d]
	Tính giá  trị của biểu thức $P=\left(7+4\sqrt 3\right)^{2017}\left(4\sqrt 3-7\right)^{2016}$.
	 \choice
		{$P=\left(7+4\sqrt 3\right)^{2016}$} 
		{$P=1$} 
		{$P=7-4\sqrt  3$}
		{\True $P=7+4\sqrt  3$}
	\loigiai{ 
		\begin{eqnarray*}
			P&=&\left(7+4\sqrt 3\right)^{2017}\left(4\sqrt  3-7\right)^{2016}\\
			&=&\left(7+4\sqrt  3\right)\cdot \left[\left(7+4\sqrt 3\right)\left(4\sqrt 3-7\right) \right]^{2016} \\
			&=&\left(7+4\sqrt 3\right)\left(-1\right)^{2016}\\ 
			 &=&7+4\sqrt 3
		\end{eqnarray*}
	 }
\end{ex}

% nguồn: Ok_TracNghiemLop11_Lan1
\begin{ex}%[1D6H1-2]%[5-CTST]%[ma:fc64a5df]
	Cho biểu thức $P =  x^{\tfrac{1}{2}} \cdot x^{\tfrac{1}{3}} \cdot \sqrt[6]{x}$ với $x > 0$. Mệnh đề nào dưới đây đúng?
	 \choice
		 {\True $P = x$}
		 {$P = x^{\tfrac{11}{6}}$}
		{$P =  x^{\tfrac{7}{6}}$}
		 {$P = x^{\tfrac{5}{6}}$}
	 \loigiai{
		Ta có $P = x^{\tfrac{1}{2}} \cdot x^{\tfrac{1}{3}} \cdot \sqrt[6]{x} = x^{\tfrac{1}{2} + \tfrac{1}{3} + \tfrac{1}{6}} =  x$.
	 }
\end{ex}

% nguồn: Ok_TracNghiemLop11_Lan1
\begin{ex}%[1D6H1-2]%[ma:1fd5cbb9]
	Với mọi  số thực dương $x$ tùy ý, biểu thức $P=x^{\frac{1}{2}}\cdot x^{\frac{1}{3}}\cdot\sqrt[6]{x}$ bằng
	 \choice
		{$x^{\frac{5}{6}}$} 
		 {\True $x$}
		 {$x^{\frac{11}{6}}$}
		{$x^{\frac{7}{6}}$} 
	 \loigiai{
		Ta có $P=x^{\frac{1}{2}}\cdot x^{\frac{1}{3}}\cdot\sqrt[6]{x}=x^{\frac{1}{2}}\cdot x^{\frac{1}{3}}\cdot  x^{\frac{1}{6}}=x^{\frac{1}{2}+\frac{1}{3}+\frac{1}{6}}=x$.
	} 
\end{ex}

% nguồn: Ok_Toan11_TracNghiem_Lan2
\begin{ex}%[1D6V1-2]%[ma:1b00d811]
	Cho $m$, $n$ là hai số dương không  đồng thời bằng $1$, biểu thức
	$\dfrac{m^{2\sqrt{2}}-n^{2\sqrt{3}}}{\left(m^{\sqrt{2}}-n^{\sqrt{3}}\right)^2}-1$ bằng 
	 \choice
		{\True  $\dfrac{2n^{\sqrt{3}}}{m^{\sqrt{2}}-n^{\sqrt{3}}}$}
		{$\dfrac{-2n^{\sqrt{3}}}{m^{\sqrt{2}}-n^{\sqrt{3}}}$} 
		 {$\dfrac{2m^{\sqrt{2}}}{m^{\sqrt{2}}-n^{\sqrt{3}}}$}
		{$\dfrac{-2m^{\sqrt{2}}}{m^{\sqrt{2}}-n^{\sqrt{3}}}$} 
	 \loigiai{
		 Ta có
		\allowdisplaybreaks 
		\begin{align*}
			\displaystyle 
			 &\dfrac{m^{2\sqrt{2}}-n^{2\sqrt{3}}}{\left(m^{\sqrt{2}}-n^{\sqrt{3}}\right)^2}-1\\
			 &=\dfrac{m^{2\sqrt{2}}-n^{2\sqrt{3}}-\left(m^{\sqrt{2}}-n^{\sqrt{3}}\right)^2}
			 {\left(m^{\sqrt{2}}-n^{\sqrt{3}}\right)^2}\\
			&=\dfrac{m^{2\sqrt{2}}-n^{2\sqrt{3}}-m^{2\sqrt{2}}-n^{2\sqrt{3}}+2m^{\sqrt{2}}n^{\sqrt{3}}} 
			{\left(m^{\sqrt{2}}-n^{\sqrt{3}}\right)^2}\\ 
			 &=\dfrac{2n^{\sqrt{3}}\left(m^{\sqrt{2}}-n^{\sqrt{3}}\right)}
			{\left(m^{\sqrt{2}}-n^{\sqrt{3}}\right)^2}\\ 
			&=\dfrac{2n^{\sqrt{3}}}{m^{\sqrt{2}}-n^{\sqrt{3}}}. 
		\end{align*}
	} 
\end{ex}

% nguồn: Ok_Toan11_TracNghiem_Lan2
\begin{ex}%[1D6V1-2]%[ma:4ce26f11]
	Cho $a>0$, $b>0$,  giá trị của biểu thức $T=2\left(a+b\right)^{-1}\cdot \left(ab\right)^{\tfrac{1}{2}}\cdot \left[1+\dfrac{1}{4}\left(\sqrt{\dfrac{a}{b}}-\sqrt{\dfrac{b}{a}}\right)^2\right]^{\tfrac{1}{2}}$ bằng
	\choice 
		 {\True $1$}
		 {$\dfrac{1}{2}$}
		 {$\dfrac{2}{3}$}
		{$\dfrac{1}{3}$} 
	 \loigiai{
		 Ta có
		\allowdisplaybreaks 
		\begin{eqnarray*}
			 T&=&2\left(a+b\right)^{-1}\cdot \left(ab\right)^{\tfrac{1}{2}}\cdot \left[1+\dfrac{1}{4}\left(\sqrt{\dfrac{a}{b}}-\sqrt{\dfrac{b}{a}}\right)^2\right]^{\tfrac{1}{2}}\\
			&=&2\left(a+b\right)^{-1}\cdot \left(ab\right)^{\tfrac{1}{2}}\cdot \left[1+\dfrac{1}{4}\left(\dfrac{a-b}{\sqrt{ab}}\right)^2\right]^{\tfrac{1}{2}}\\ 
			 &=&2\left(a+b\right)^{-1}\cdot \left(ab\right)^{\tfrac{1}{2}}\cdot \left[1+\dfrac{\left(a-b\right)^2}{4ab}\right]^{\tfrac{1}{2}}\\
			 &=&2\left(a+b\right)^{-1}\cdot \left(ab\right)^{\tfrac{1}{2}}\cdot \left[\dfrac{\left(a+b\right)^2}{4ab}\right]^{\tfrac{1}{2}}\\
			 &=&2\dfrac{1}{a+b}\cdot \left(ab\right)^{\tfrac{1}{2}}\cdot \dfrac{\left(a+b\right)}{2\left(ab\right)^{\tfrac{1}{2}}}=1.
		\end{eqnarray*}
	 }
\end{ex}

% nguồn: Ok_Toan11_TracNghiem_Lan2
\begin{ex}%[1D6V1-2]%[ma:e8c977e3]
	Cho biểu thức $\sqrt[5]{8\sqrt{2\sqrt[3]{2}}}={{2}^{\tfrac{m}{n}}}$, trong đó $\dfrac{m}{n}$ là phân số tối giản.  Gọi $P=m^2+n^2$. Khẳng định nào sau đây đúng?
	\choice 
		{$P\in \left(330;340\right)$} 
		 {$P\in \left(350;360\right)$}
		{$P\in  \left(260;370\right)$}
		 {\True $P\in \left(340;350\right)$}
	 \loigiai{
		Ta có  $\sqrt[5]{8\sqrt{2\sqrt[3]{2}}}=\sqrt[5]{{{2}^{3}}\sqrt{2\sqrt[3]{2}}}={{2}^{\tfrac{3}{5}}}\cdot {{2}^{\tfrac{1}{10}}}\cdot {{2}^{\tfrac{1}{30}}}={{2}^{\tfrac{3}{5}+\tfrac{1}{10}+\tfrac{1}{30}}}={{2}^{\tfrac{11}{15}}}$.\\
		 Suy ra $\dfrac{m}{n}=\dfrac{11}{15}\Rightarrow \heva{&m=11 \\&n=15.}$\\
	 Do đó $P=m^2+n^2=11^2+15^2=346$.}
\end{ex}

% nguồn: Ok_Toan11_TracNghiem_Lan2
\begin{ex}%[1D6V1-2]%[ma:04fa8449]
	Cho hàm số $f(a)=\dfrac{a^{\tfrac23}\left(\sqrt[3]{a^{-2}}-\sqrt[3]{a}\right)}{a^{\tfrac18}\left(\sqrt[8]{a^{3}}-\sqrt[8]{a^{-1}}\right)}$ với $a>0$, $a\ne 1$. Giá trị của $M=f\left(2021^{2022}\right)$ là 
	\choice 
		 {$2021^{1011}$}
		 {$2021^{1011}+1$}
		{$-2021^{1011}+1$} 
		{\True $-2021^{1011}-1$} 
	\loigiai{ 
		Ta  có
		\[ 
		 f(a)=\dfrac{a^{\tfrac23}\left(\sqrt[3]{a^{-2}}-\sqrt[3]{a}\right)}{a^{\tfrac18}\left(\sqrt[8]{a^{3}}-\sqrt[8]{a^{-1}}\right)}
		 =\dfrac{a^{\tfrac23}\left(a^{-\tfrac23}-a^{\tfrac13}\right)}{a^{\tfrac18}\left(a^{\tfrac38}-a^{-\tfrac18}\right)}
		 =\dfrac{1-a}{a^{\tfrac12}-1}
		 =-a^{\tfrac12}-1.
		 \]
		 Khi đó
		\[ 
		M=f\left(2021^{2022}\right) 
		= -\left(2021^{2022}\right)^{\tfrac12}-1 
		 =-2021^{1011}-1.
		 \]
	 }
\end{ex}

% nguồn: Ok_TracNghiemLop11_Lan1
\begin{ex}%[1D6B1-2]%[ma:b808d88a]
	Rút gọn biểu thức $P=x^{\tfrac{1}{3}}\cdot \sqrt[6]{x}$  với $x>0$.
	\choice 
		{\True $P=\sqrt{x}$} 
		{$P=x^{\tfrac{1}{8}}$} 
		{$P=x^{\tfrac{2}{9}}$} 
		{$P=x^2$} 
	\loigiai{ 
	 Ta có
	\begin{align*}
		 P=x^{\tfrac{1}{3}} \cdot \sqrt[6]{x}=x^{\tfrac{1}{3}} \cdot x^{\tfrac{1}{6}}=x^{\tfrac{1}{3}+\tfrac{1}{6}}=x^{\tfrac{1}{2}}=\sqrt{x}.
	\end{align*}
	} 
\end{ex}


%%%%%%%%%%%%%%%% ĐÚNG/SAI %%%%%%%%%%%%%%%%

% nguồn: Ok_dungsai_lop11_Lan1
\begin{ex}%[1D6H1-2]%[ma:dc038574]
 Cho biểu thức $Q=(\sqrt{x}-\sqrt[4]x+1)(\sqrt{x}+\sqrt[4]x+1)(x-\sqrt{x}+1)$ với  $x \geq 0$.
 \choiceTF 
  {\True Khi  $x=2$ thì $Q=7$}
  {Phương  trình $Q=0\Rightarrow x=1+\sqrt{2}$}
  {Phương trình  $Q=3$ có hai nghiệm dương phân biệt}
  {\True  Khi $x=3$ thì $Q$ là một số nguyên tố}
 \loigiai{ 
 Ta  có
 \allowdisplaybreaks 
 \begin{eqnarray*}
  Q&=&(\sqrt{x}-\sqrt[4]x+1)(\sqrt{x}+\sqrt[4]x+1)(x-\sqrt{x}+1)\\ 
   &=&\left[(\sqrt{x}+1)^2-(\sqrt[4]{x})^2\right](x+1-\sqrt{x}) \\
   &=&[(x+2 \sqrt{x}+1)-\sqrt{x}](x+1-\sqrt{x}) \\
   &=&(x+1+\sqrt{x})(x+1-\sqrt{x})=(x+1)^2-(\sqrt{x})^2\\
  &=&x^2+x+1. 
 \end{eqnarray*}
 Vậy $Q=x^2+x+1$. 
 \begin{itemchoice}
  \itemch  \textbf{Đúng}.\\
  Với $x=2\Rightarrow  Q=2^2+2+1=7$
  \itemch \textbf{Sai}.\\ 
  Khi $Q=0\Leftrightarrow x^2+x+1=0$ vô  nghiệm.
   \itemch \textbf{Sai}.\\
   Khi $Q=3\Leftrightarrow x^2+x+1=3\Leftrightarrow x^2+x-2=0\Leftrightarrow\hoac{&x=1\\
   &x=-2.}$\\
   Phương trình có $2$ nghiệm trái dấu nhau.
  \itemch  \textbf{Đúng}.\\
  Với $x=3\Rightarrow  Q=3^2+3+1=13$
 \end{itemchoice}
 } 
\end{ex}

% nguồn: Ok_dungsai_lop11_Lan1
\begin{ex}%[1D6H1-2]%[ma:8b598fb0]
 Cho biểu thức $A=(\sqrt{2}-1)^{2  x}+(3+2 \sqrt{2})^x$.
  \choiceTF
  {\True Cho  $\left(\sqrt{2}+1\right)^x=3$ thì $A=\dfrac{82}{9}$}
   {\True Cho $\left(\sqrt{2}+1\right)^x=1$ thì $A=2$}
  {Cho  $\left(\sqrt{2}+1\right)^x=2$ thì $A=\dfrac{17}{9}$}
  {Cho  $\left(\sqrt{2}+1\right)^x=m^2$ thì $A=\dfrac{1+m^6}{m^4}$}
  \loigiai{
  Ta có $\sqrt{2}-1=\dfrac{(\sqrt{2}-1)(\sqrt{2}+1)}{\sqrt{2}+1}=\dfrac{1}{\sqrt{2}+1}=(\sqrt{2}+1)^{-1}$;\\
 Và $3+2  \sqrt{2}=(\sqrt{2})^2+2 \sqrt{2}+1=(\sqrt{2}+1)^2$.\\
 Do đó, 
  \allowdisplaybreaks
 \begin{eqnarray*}
   A&=&\left[(\sqrt{2}+1)^{-1}\right]^{2 x}+\left[(\sqrt{2}+1)^2\right]^x\\
  &=&(\sqrt{2}+1)^{-2 x}+(\sqrt{2}+1)^{2 x}\\ 
   &=&\left[(\sqrt{2}+1)^x\right]^{-2}+\left[(\sqrt{2}+1)^x\right]^2.
 \end{eqnarray*}
 Vậy $A=\left[(\sqrt{2}+1)^x\right]^{-2}+\left[(\sqrt{2}+1)^x\right]^2$. 
 \begin{itemchoice}
  \itemch  \textbf{Đúng}.\\
  Với $\left(\sqrt{2}+1\right)^x=3 \Rightarrow A=3^{-2}+3^2=\dfrac{82}{9}$. 
  \itemch \textbf{Đúng}.\\ 
  Với $\left(\sqrt{2}+1\right)^x=1\Rightarrow A=1^{-2}+1^2=2$. 
  \itemch \textbf{Sai}.\\ 
  Với  $\left(\sqrt{2}+1\right)^x=2 \Rightarrow A=2^{-2}+2^2=\dfrac{17}{4}$.
  \itemch \textbf{Sai}.\\ 
  Với $\left(\sqrt{2}+1\right)^x=m \Rightarrow A=m^{-4}+m^4=\dfrac{1+m^8}{m^4}$. 
 \end{itemchoice}
 } 
\end{ex}

% nguồn: Ok_dungsai_lop11_Lan1
\begin{ex}%[1D6H1-2]%[ma:81d5704a]
 Cho các biểu thức $A=\dfrac{\sqrt{a}+\sqrt[4]{a b}}{\sqrt[4]a+\sqrt[4]b}-\dfrac{\sqrt{a}-\sqrt{b}}{\sqrt[4]a-\sqrt[4]b}$; $B=-\sqrt[4]b$; với $a > 0$, $b  > 0$, $a \neq b$.
 \choiceTF 
  {\True Sau  khi rút gọn, thì $A$ chỉ chứa biến $b$}
  {Biểu thức $A>0$ luôn đúng} 
   {$A=B+\sqrt{a}$}
   {\True $\dfrac{A-1}{B}=1-\dfrac{1}{B}$}
  \loigiai{
  Ta có
 \allowdisplaybreaks 
 \begin{eqnarray*}
  \dfrac{\sqrt{a}+\sqrt[4]{a  b}}{\sqrt[4]a+\sqrt[4]b}-\dfrac{\sqrt{a}-\sqrt{b}}{\sqrt[4]a-\sqrt[4]b}&=&\dfrac{(\sqrt[4]a)^2+\sqrt[4]{a b}}{\sqrt[4]a+\sqrt[4]b}-\dfrac{(\sqrt[4]a)^2-(\sqrt[4]b)^2}{\sqrt[4]a-\sqrt[4]b}\\ &=&\dfrac{\sqrt[4]a(\sqrt[4]a+\sqrt[4]b)}{\sqrt[4]a+\sqrt[4]b}-\dfrac{(\sqrt[4]a-\sqrt[4]b)(\sqrt[4]a+\sqrt[4]b)}{\sqrt[4]a-\sqrt[4]b}\\
   &=&\sqrt[4]a-(\sqrt[4]a+\sqrt[4]b)=-\sqrt[4]b.
 \end{eqnarray*}
 \begin{itemchoice}
  \itemch \textbf{Đúng}.\\ 
  Ta có  $A=-\sqrt[4]b$.
   \itemch \textbf{Sai}.\\
  Ta có $A=-\sqrt[4]b<0$  với mọi $b>0$.
  \itemch  \textbf{Sai}.\\
  Ta có  $A=B=\sqrt[4]{b}$.
   \itemch \textbf{Đúng}.\\
   Do $A=B$ nên $\dfrac{A-1}{B}=\dfrac{B-1}{B}=1-\dfrac{1}{B}$.
 \end{itemchoice}
 } 
\end{ex}

% nguồn: Ok_dungsai_lop11_Lan2
\begin{ex}%[1D6H1-2]%[ma:d5d559b1]
	Cho $ a>0$,  $ b>0$ và hàm số $ f(x)=\dfrac{9^x}{9^x+3}$ . Biết $m$, $n$ là các số thực thỏa mãn
	$$\sqrt{a}.\sqrt[5]{a^2\sqrt{a}}=a^m,\,  \dfrac{b^{\tfrac{4}{3}}\left(b^{\tfrac{-1}{3}}+b^{\tfrac{2}{3}}\right)}{b^{\tfrac{1}{4}}\left(b^{\tfrac{3}{4}}+b^{\tfrac{-1}{4}}\right)}=b^n.$$
	 \choiceTF
		 {$ f(0)=0$}
		{\True $ m\in\mathbb{Z}$} 
		{\True $m=n$} 
		{Nếu  $ f\left(m-a\right)+f\left(n-b\right)=2$ thì $ f\left(a^m\right)+f\left(b^n\right)=0$}
	\loigiai{ 
		\begin{itemchoice}
			\itemch $  f(0)=\dfrac{9^0}{9^0+3}=\dfrac{1}{4}$.
			\itemch $\sqrt{a}\cdot\sqrt[5]{a^2\sqrt{a}}=a^{\tfrac{1}{2}}{\left(a^2a^{\tfrac{1}{2}}\right)^{\tfrac{1}{5}}}=a\Rightarrow  m=1\in\mathbb{Z}$.
			\itemch $\dfrac{b^{\tfrac{4}{3}}\left(b^{\tfrac{-1}{3}}+b^{\tfrac{2}{3}}\right)}{b^{\tfrac{1}{4}}\left(b^{\tfrac{3}{4}}+b^{\tfrac{-1}{4}}\right)}=\dfrac{b+b^2}{b+1}=b\Rightarrow n=1\Rightarrow m=n$. 
			 \itemch $ f(x)=\dfrac{9^x}{9^x+3}\Rightarrow f\left(1-x\right)=\dfrac{9^{1-x}}{9^{1-x}+3}=\dfrac{3}{9^x+3}=1-f(x)$.\\
			$ f\left(m-a\right)+f\left(n-b\right)=2\Rightarrow  2=f\left(1-a\right)+f\left(1-b\right)=2-f(a)-f(b)=2-f\left(a^m\right)-f\left(b^n\right)$.\\
			 Suy ra: $ f\left(a^m\right)+f\left(b^n\right)=0$.
		\end{itemchoice}
	 }
\end{ex}

% nguồn: Ok_dungsai_lop11_Lan2
\begin{ex}%[1D6H1-2]%[ma:4493927a]
	Cho biểu thức $9^5\cdot 27^{\tfrac25}=A$ và $144^{\tfrac34}:9^{\tfrac34}=B$. Khi đó 
	\choiceTF 
		{$9^5\cdot 27^{\tfrac25}=\left(9\cdot  27\right)^{\tfrac25}$}
		 { $9^5\cdot 27^{\tfrac25}=3^{k}$ thì $k=3$ }
		{  \True $144^{\tfrac34}:9^{\tfrac34}=2^{k}$ thì $k=3$ }
		{\True Phép toán $A-B$ thu được kết quả là một số tự nhiên} 
	\loigiai{ 
		\begin{itemchoice}
			\itemch $9^5\cdot 27^{\tfrac25}\ne \left(9\cdot  27\right)^{\tfrac25}$.
			\itemch  $9^5\cdot 27^{\tfrac25}=\left(3^2\cdot 3^{3}\right)^{\tfrac25}
			=\left(3^5\right)^{\tfrac25} 
			 =3^2=9$.\\
			Khi  đó $3^k=9\Rightarrow k=2$.
			\itemch  $144^{\tfrac34}:9^{\tfrac34}
			 =\left(\dfrac{144}{9}\right)^{\tfrac34}
			=16^{\tfrac34} 
			=\left(2^{4}\right)^{\tfrac34} 
			=2^{3}=8$. 
			\itemch Ta  có
			\begin{itemize}
				\item $A=9^5\cdot  27^{\tfrac25}
				=\left(3^2\cdot  3^{3}\right)^{\tfrac25}
				=\left(3^5\right)^{\tfrac25} 
				=3^2=9$. 
				\item $B=144^{\tfrac34}:9^{\tfrac34} 
				 =\left(\dfrac{144}{9}\right)^{\tfrac34}
				=16^{\tfrac34} 
				=\left(2^{4}\right)^{\tfrac34} 
				 =2^{3}=8$.
			\end{itemize}
			Khi  đó $A-B=1$.
		\end{itemchoice}
	} 
\end{ex}

% nguồn: Ok_dungsai_lop11_Lan2
\begin{ex}%[1D6H1-2]%[ma:a45eb48e]
	Với mọi số thực $a>0$, $b>0$. Giả sử  các biểu thức xuất hiện trong các công thức của mỗi mệnh đề đều có nghĩa. Xét tính đúng sai của các khẳng định sau
	\choiceTF 
		{\True $\sqrt[3]{a^2}=a^{\tfrac23}$} 
		{ $a^{\tfrac13}\sqrt{a}=a^{\tfrac16}$} 
		 { $\dfrac{\sqrt[3]{b^7}}{b^3}=b^{\tfrac79}$}
		 {$\dfrac{\left(\sqrt{ab^2}\right)^4}{\sqrt[3]{a^6b^{12}}}=ab^2$}
	\loigiai{ 
		\begin{itemchoice}
			\itemch $\sqrt[3]{a^2}=a^{\tfrac23}$,  với $a>0$.
			 \itemch $a^{\tfrac13}\sqrt{a}=a^{\tfrac13}a^{\tfrac12}=a^{\tfrac56}\ne a^{\tfrac16}$.
			\itemch $\dfrac{\sqrt[3]{b^7}}{b^3}=b^{\tfrac73} \cdot  b^{-3}=b^{-\tfrac23}\ne b^{\tfrac79}$.
			\itemch  $\dfrac{\left(\sqrt{ab^2}\right)^4}{\sqrt[3]{a^6b^{12}}}
			 =\dfrac{\left(ab^2\right)^2}{a^2b^4}
			=1$. 
		\end{itemchoice}
	} 
\end{ex}

% nguồn: Ok_dungsai_lop11_Lan2
\begin{ex}%[1D6H1-2]%[ma:37fb7523]
	Cho $a$, $b$ là số thực dương  và $x$, $y\in\mathbb{R}$.
	\choiceTF 
		{$\left(a^x\right)^y=a^{x^y}$} 
		{Nếu $\left(\sqrt{2}-1\right)^x<\left(\sqrt{2}-1\right)^y$  thì $x<y$}
		{\True  Rút gọn $ P=\dfrac{a^{\sqrt{5}+1}\cdot a^{2-\sqrt{5}}}{\left(a^{\sqrt{2}-2}\right)^{\sqrt{2}+2}}$ ta được $ P=a^5$}
		{\True  Nếu $f(x)=\dfrac{9^x-2}{9^x+3}$ thì $S=f\left(\dfrac{1}{2\,023}\right)+f\left(\dfrac{2}{2\,023}\right)+\cdots +f\left(\dfrac{2\,022}{2\,023}\right)=\dfrac{1\,011}{3}$}
	 \loigiai{
		\begin{itemchoice}
			\itemch  $\left(2^1\right)^3=2^3\ne{2^{1^3}}$.
			\itemch  $\heva{&
			{\left(\sqrt{2}-1\right)^x}<\left(\sqrt{2}-1\right)^y\\ 
			&0<\sqrt{2}-1<1}\Rightarrow x>y$. 
			 \itemch $P=\dfrac{a^{\sqrt{5}+1}\cdot a^{2-\sqrt{5}}}{\left(a^{\sqrt{2}-2}\right)^{\sqrt{2}+2}}=\dfrac{a^{\sqrt{5}+1+2-\sqrt{5}}}{a^{\left(\sqrt{2}-2\right)\left(\sqrt{2}+2\right)}}=\dfrac{a^3}{a^{-2}}=a^5$.
			\itemch  \allowdisplaybreaks
			\begin{eqnarray*}
				 f(x)+f\left(1-x\right)&=&\dfrac{9^x-2}{9^x+3}+\dfrac{9^{1-x}-2}{9^{1-x}+3}\\
				&=&\dfrac{9^x-2}{9^x+3}+\dfrac{9-2\cdot9^x}{9+3\cdot9^x}\\ 
				&=&\dfrac{9^x-2}{9^x+3}+\dfrac{3-\dfrac{2}{3}\cdot{9^x}}{3+9^x}=\dfrac{1+\dfrac{1}{3}\cdot{9^x}}{9^x+3}=\dfrac{3+9^x}{3\left(9^x+3\right)}=\dfrac{1}{3}. 
			\end{eqnarray*}
			Do đó $S=\underset{k=1}{\overset{1\,011}{\mathop\sum}}\,\left[f\left(\dfrac{k}{2\,023}\right)+f\left(1-\dfrac{k}{2\,023}\right)\right]=\underset{k=1}{\overset{1\,011}{\mathop\sum}}\,\dfrac{1}{3}=\dfrac{1\,011}{3}$. 
	\end{itemchoice}}
\end{ex}

% nguồn: Ok_dungsai_lop11_Lan2
\begin{ex}%[1D6H1-2]%[ma:d0340134]
	Cho biểu  thức $A=3^{2 x-1}\cdot \left(\dfrac{1}{3}\right)^{2 x-1}+9^{x+1}$.
	 \choiceTF
		{\True Nếu $3^x=2$  thì $A=37$}
		{\True Nếu $3^x=1$ thì $A=10$} 
		 {Nếu $3^x=3$ thì $A=80$}
		{Nếu $3^x=4$ thì $A=144$} 
	 \loigiai{
		 Ta có $A=\left(3\cdot \dfrac{1}{3}\right)^{2x-1}+9\cdot {9^x}=1+9\cdot(3^x)^2$.
		\begin{itemchoice}
			\itemch  Nếu $3^x=2$ thì $A=1+9\cdot 2^2=37$.
			\itemch Nếu  $3^x=1$ thì $A=1+9\cdot 1^2=10$.
			\itemch Nếu  $3^x=3$ thì $A=1+9\cdot 3^2=82$.
			\itemch Nếu  $3^x=4$ thì $A=1+9\cdot 4^2=145$.
		\end{itemchoice}
	 }
\end{ex}

% nguồn: Ok_dungsai_lop11_Lan2
\begin{ex}%[1D6H1-2]%[ma:2a2018d4]
	Cho $a$, $b$ là các số thực dương. Xét tính đúng sai của các khẳng định  sau
	\choiceTF 
		{ $a^m b^n  = (ab)^{m+n}$ }
		{\True $\dfrac{a^m}{a^n} = a^{m-n}$} 
		{\True  $\left(a^m\right)^n = a^{m\cdot n}$}
		 { $a^m + a^n = a^{m+n}$ }
	\loigiai{ 
		\begin{itemchoice}
			\itemch $a^m  \cdot b^n \neq (ab)^{m+n}$.
			\itemch  $\dfrac{a^m}{a^n} = a^{m-n}$.
			\itemch $\left(a^m\right)^n = a^{m\cdot n}$. 
			\itemch $a^m +  a^n \neq a^{m+n}$.
		\end{itemchoice}
	} 
\end{ex}

% nguồn: Ok_dungsai_lop11_Lan2
\begin{ex}%[1D6H1-2]%[ma:3045ce60]
	Với mọi $a>0$, $b>0$ và $m$, $n$ là các số thực tùy ý. Giả sử các biểu  thức xuất hiện trong các công thức của mỗi mệnh đề đều có nghĩa.
	\choiceTF 
		{\True  $a^m\cdot a^n = a^{m+n}$}
		{\True $a^{-n} = \dfrac{1}{a^n}$  với $n$ là số nguyên dương}
		{\True Nếu  $\left(\dfrac{5}{2}\right)^m > \left(\dfrac{5}{2}\right)^n$ thì $m>n$}
		{Nếu $\left(\dfrac{\pi}{4}\right)^m > \left(\dfrac{\pi}{4}\right)^{3m-2}$ thì $m<1$} 
	 \loigiai{
		\begin{itemchoice}
			\itemch $a^m\cdot a^n = a^{m+n}$, với  $a>0$ và $m$, $n$ là các số thực tùy ý.
			\itemch $a^{-n} = \dfrac{1}{a^n}$ với $n$ là số  nguyên dương.
			\itemch Do $\dfrac52>1$ nên $\left(\dfrac52\right)^m > \left(\dfrac52\right)^n \Leftrightarrow  m>n$.
			\itemch Do $0<\dfrac{\pi}{4}<1$ nên $\left(\dfrac{\pi}{4}\right)^m > \left(\dfrac{\pi}{4}\right)^{3m-2} \Leftrightarrow  m<3m-2 \Leftrightarrow m>1$.
		\end{itemchoice}
	} 
\end{ex}

% nguồn: Ok_dungsai_lop11_Lan2
\begin{ex}%[1D6V1-2]%[ma:7ddeace7]
	Cho biểu thức $A=4^{\tfrac{2}{5}}  \cdot 8^{\tfrac{2}{5}}$ và $B=9^{\tfrac{3}{4}}:144^{\tfrac{3}{4}}$.
	\choiceTF 
		{\True $4^{\tfrac{2}{5}} \cdot 8^{\tfrac{2}{5}}=\left(4\cdot 8\right)^{\tfrac{2}{5}}$} 
		{$4^{\tfrac{2}{5}}  \cdot 8^{\tfrac{2}{5}}=2^k$ thì $k=3$}
		{$9^{\tfrac{3}{4}}:144^{\tfrac{3}{4}}=2^m$  thì $m=3$}
		 {\True $A:B=32$}
	 \loigiai{
		\begin{itemchoice}
			\itemch $A=4^{\tfrac{2}{5}} \cdot 8^{\tfrac{2}{5}}=\left(4\cdot 8\right)^{\tfrac{2}{5}}=\left(2^2\cdot 2^3\right)^{\tfrac{2}{5}}=\left(2^{2+3}  \right)^{\tfrac{2}{5}}=(2^5)^{\tfrac{2}{5}}=2^{5\cdot\tfrac{2}{5}}=2^2$.
			\itemch Ta có  $A=2^2$. Suy ra $k=2$.
			\itemch $B=9^{\tfrac{3}{4}}:144^{\tfrac{3}{4}}=\left(\dfrac{9}{144} \right)^{\tfrac{3}{4}}=\left(\dfrac{1}{16} \right)^{\tfrac{3}{4}}=\left(\dfrac{1}{2^4}  \right)^{\tfrac{3}{4}}=\left(2^{-4} \right)^{\tfrac{3}{4}}=2^{-4\cdot\dfrac{3}{4}}=2^{-3}$.\\
			Ta có $B=2^{-3}$. Suy  ra $m=-3$.
			\itemch Ta có $A:B=2^2:2^{-3}=2^{5}$. Do  đó $A:B=2^2:2^{-3}=2^{2+3}=32$.
		\end{itemchoice}
	 }
\end{ex}

% nguồn: Ok_dungsai_lop11_Lan2
\begin{ex}%[1D6V1-2]%[ma:8b560416]
	Cho biểu thức $P=\dfrac{a^{\tfrac{1}{3}} \sqrt{b}+b^{\tfrac{1}{3}} \sqrt{a}}{\sqrt[6]a+\sqrt[6]b}$ với $a$, $b > 0$. Rút gọn biểu thức $P$ ta được $P=a^m\cdot b^n$. Khi  đó
	\choiceTF 
		{$m$, $n$  là các số nguyên}
		{$P^6=ab$} 
		{\True  $m+2n=1$}
		{\True Với  $a=8$, $b=27$ thì $P=6$}
	\loigiai{ 
		 Ta có \allowdisplaybreaks
		\begin{eqnarray*}
			P&=&\dfrac{a^{\tfrac{1}{3}} b^{\tfrac{1}{2}}+b^{\tfrac{1}{3}}  a^{\tfrac{1}{2}}}{a^{\tfrac{1}{6}}+b^{\tfrac{1}{6}}}\\
			&=&\dfrac{a^{\tfrac{1}{3}} b^{\tfrac{1}{3}} \left(b^{\tfrac{1}{6}}+a^{\tfrac{1}{6}} \right)}{a^{\tfrac{1}{6}}+b^{\tfrac{1}{6}}}\\ 
			&=&a^{\tfrac{1}{3}}  b^{\tfrac{1}{3}}.
		\end{eqnarray*}
		\begin{itemchoice}
			\itemch Ta có $P=a^{\tfrac{1}{3}} b^{\tfrac{1}{3}}$ suy ra $m=n=\dfrac{1}{3}$. Do đó $m$, $n\in\mathbb{Q}$. 
			\itemch Ta có $P=a^{\tfrac{1}{3}} b^{\tfrac{1}{3}}=(a\cdot b)^{\tfrac{1}{3}}$ suy  ra $P^6=(ab)^2$.
			\itemch Ta có $P=a^{\tfrac{1}{3}} b^{\tfrac{1}{3}}$ suy ra $m=n=\dfrac{1}{3}$.  Do đó $m+2n=1$.
			\itemch  Với $a=8$, $b=27$ thì $P=(8\cdot 27)^{\tfrac{1}{3}}=2\cdot 3=6$.
		\end{itemchoice}
	 }
\end{ex}

% nguồn: Ok_dungsai_lop11_Lan2
\begin{ex}%[1D6V1-2]%[ma:fb700512]
	Cho biểu  thức $\sqrt[3]{\dfrac{2}{5}\sqrt[7]{\dfrac{2}{5}}\sqrt[3]{\dfrac{2}{5}}}\cdot \dfrac{2}{5}$.
	 \choiceTF
		{\True  $\sqrt[3]{\dfrac{2}{5}}-\left(\dfrac{2}{5} \right)^{\tfrac{1}{3}}=0$}
		{$\sqrt[7]{\dfrac{2}{5} \cdot \sqrt[3]{\dfrac{2}{5}}}=\left(\dfrac{2}{5} \right)^{\tfrac{a}{b}}$, ($\dfrac{a}{b}$ là phân số tối giản), khi đó $a+b=24$} 
		{\True $\sqrt[3]{\dfrac{2}{5} \cdot \sqrt[7]{\dfrac{2}{5} \cdot \sqrt[3]{\dfrac{2}{5}}}}=\left(\dfrac{2}{5} \right)^{\tfrac{c}{d}}$, ($\dfrac{c}{d}$ là phân số tối giản),  khi đó $c+d=88$}
		{\True $\sqrt[3]{\dfrac{2}{5} \cdot \sqrt[7]{\dfrac{2}{5} \cdot  \sqrt[3]{\dfrac{2}{5}}}} \cdot \dfrac{2}{5}=\left(\dfrac{2}{5} \right)^{\tfrac{m}{n}}$, ($\dfrac{m}{n}$ là phân số tối giản), khi đó $m+n=151$}
	\loigiai{ 
		\begin{itemchoice}
			\itemch Ta có $\sqrt[3]{\dfrac{2}{5}}-\left(\dfrac{2}{5}  \right)^{\tfrac{1}{3}}=\left(\dfrac{2}{5} \right)^{\tfrac{1}{3}}-\left(\dfrac{2}{5} \right)^{\tfrac{1}{3}}=0$.
			\itemch Ta có  \allowdisplaybreaks
			\begin{eqnarray*}
				 \sqrt[7]{\dfrac{2}{5} \cdot \sqrt[3]{\dfrac{2}{5}}}&=&\left(\dfrac{2}{5} \right)^{\tfrac{1}{7}}\cdot\left[\left(\dfrac{2}{5} \right)^{\tfrac{1}{3}}\right]^{\tfrac{1}{7}}\\
				&=&\left(\dfrac{2}{5} \right)^{\tfrac{1}{7}}\cdot  \left(\dfrac{2}{5} \right)^{\tfrac{1}{21}}\\
				&=&\left(\dfrac{2}{5}  \right)^{\tfrac{1}{7}+\tfrac{1}{21}}\\
				 &=&\left(\dfrac{2}{5} \right)^{\tfrac{4}{21}}.
			\end{eqnarray*}
			Mà $\sqrt[7]{\dfrac{2}{5} \cdot \sqrt[3]{\dfrac{2}{5}}}=\left(\dfrac{2}{5} \right)^{\tfrac{a}{b}}$ nên $a=4$; $b=21$.  \\
			 Do đó $a+b=25$.
			 \itemch Ta có \allowdisplaybreaks
			\begin{eqnarray*}
				 \sqrt[3]{\dfrac{2}{5} \cdot \sqrt[7]{\dfrac{2}{5} \cdot \sqrt[3]{\dfrac{2}{5}}}}&=&\left(\dfrac{2}{5} \right)^{\tfrac{1}{3}}\cdot\left[\left(\dfrac{2}{5} \right)^{\tfrac{1}{7}}\right]^{\tfrac{1}{3}}\cdot\left[\left(\dfrac{2}{5} \right)^{\tfrac{1}{21}}\right]^{\tfrac{1}{3}}\\
				&=&\left(\dfrac{2}{5} \right)^{\tfrac{1}{3}}\cdot\left(\dfrac{2}{5} \right)^{\tfrac{1}{21}}\cdot\left(\dfrac{2}{5} \right)^{\tfrac{1}{63}}\\ 
				 &=&\left(\dfrac{2}{5} \right)^{\tfrac{1}{3}+\tfrac{1}{21}+\tfrac{1}{63}}\\
				&=&\left(\dfrac{2}{5} \right)^{\tfrac{25}{63}}. 
			\end{eqnarray*}
			Mà $\sqrt[3]{\dfrac{2}{5} \cdot \sqrt[7]{\dfrac{2}{5} \cdot  \sqrt[3]{\dfrac{2}{5}}}}=\left(\dfrac{2}{5} \right)^{\tfrac{c}{d}}$ nên $c=25$; $d=63$.\\
			Do đó  $c+d=88$.
			 \itemch Ta có \allowdisplaybreaks
			\begin{eqnarray*}
				\sqrt[3]{\dfrac{2}{5} \cdot \sqrt[7]{\dfrac{2}{5} \cdot  \sqrt[3]{\dfrac{2}{5}}}} \cdot \dfrac{2}{5}&=&\left(\dfrac{2}{5} \right)^{\tfrac{1}{3}}\cdot\left[\left(\dfrac{2}{5} \right)^{\tfrac{1}{7}}\right]^{\tfrac{1}{3}}\cdot\left[\left(\dfrac{2}{5} \right)^{\tfrac{1}{21}}\right]^{\tfrac{1}{3}}\cdot\dfrac{2}{5}\\
				 &=&\left(\dfrac{2}{5} \right)^{\tfrac{1}{3}}\cdot\left(\dfrac{2}{5} \right)^{\tfrac{1}{21}}\cdot\left(\dfrac{2}{5} \right)^{\tfrac{1}{63}}\cdot\dfrac{2}{5}\\
				&=&\left(\dfrac{2}{5} \right)^{\tfrac{1}{3}+\tfrac{1}{21}+\tfrac{1}{63}+1}\\ 
				&=&\left(\dfrac{2}{5}  \right)^{\tfrac{88}{63}}.
			\end{eqnarray*}
			Mà $\sqrt[3]{\dfrac{2}{5} \cdot \sqrt[7]{\dfrac{2}{5} \cdot \sqrt[3]{\dfrac{2}{5}}}} \cdot \dfrac{2}{5}=\left(\dfrac{2}{5} \right)^{\tfrac{m}{n}}$ nên $m=88$;  $n=63$.\\
			 Do đó $m+n=151$.
		\end{itemchoice}
	} 
\end{ex}

% nguồn: Ok_dungsai_lop11_Lan2
\begin{ex}%[1D6V1-2]%[ma:59cfe525]
	Cho biểu thức $9^{\tfrac{2}{5}} \cdot 27^{\tfrac{2}{5}}=A$ và $144^{\tfrac{3}{4}}:9^{\tfrac{3}{4}}=B$. 
	\choiceTF 
		{\True $9^{\tfrac{2}{5}} \cdot 27^{\tfrac{2}{5}}=(9\cdot  27)^{\tfrac{2}{5}}$}
		 {$9^{\tfrac{2}{5}} \cdot 27^{\tfrac{2}{5}}=3^k$ thì $k=3$}
		{\True $144^{\tfrac{3}{4}}:9^{\tfrac{3}{4}}=2^k$  thì $k=3$}
		 {\True $A-B=1$}
	 \loigiai{
		\begin{itemchoice}
			\itemch Ta  có $9^{\tfrac{2}{5}} \cdot 27^{\tfrac{2}{5}}=(9\cdot 27)^{\tfrac{2}{5}}$.
			 \itemch Ta có $9^{\tfrac{2}{5}} \cdot 27^{\tfrac{2}{5}}=(9\cdot 27)^{\tfrac{2}{5}}=\left(3^2\cdot 3^3\right)^{\tfrac{2}{5}}=(3^5)^{\tfrac{2}{5}}=3^2$.\\
			Mà $9^{\tfrac{2}{5}} \cdot 27^{\tfrac{2}{5}}=3^k$ nên $k=2$. 
			\itemch $144^{\tfrac{3}{4}}:9^{\tfrac{3}{4}}=\left(\dfrac{144}{9} \right)^{\tfrac{3}{4}}=16^{\tfrac{3}{4}}=(2^4)^{\tfrac{3}{4}}=2^3$. 
			\itemch 
			 Ta có $A=9$; $B=8$. Do đó $A-B=1$.
		\end{itemchoice}
	 }
\end{ex}

% nguồn: Ok_dungsai_lop11_Lan2
\begin{ex}%[1D6V1-2]%[ma:06cc73da]
	Cho các biểu thức $A=(a^2)^{3+2\sqrt{2}} \cdot a^{1-\sqrt{2}} \cdot a^{-4-\sqrt{2}}$ với $a > 0$ và $B=\sqrt[4]{x\cdot \sqrt[3]{x^2\cdot \sqrt{x^3}}}$ với $x  > 0$.
	\choiceTF 
		{\True  Với $a=2$ thì $A< 57$}
		{Với $x=2$ thì  $B> 2$}
		{Khi $x=a$  thì $A\cdot B=a^{\tfrac{39+26\sqrt{2}}{24}}$}
		{Khi  $x=a$ thì $A:B=a^{\tfrac{72+48\sqrt{2}}{13}}$}
	 \loigiai{
		Ta có 
		\begin{itemize}
			\item $A=a^{6+4\sqrt{2}} \cdot  a^{1-\sqrt{2}} \cdot a^{-4-\sqrt{2}}=a^{6+4\sqrt{2}+1-\sqrt{2}-4-\sqrt{2}}=a^{3+2\sqrt{2}}$.
			\item $B=x^{\tfrac{1}{4}}\cdot  \left(x^{\tfrac{2}{3}}\right)^{\tfrac{1}{4}}\cdot\left[\left(x^{\tfrac{3}{2}}\right)^{\tfrac{1}{3}}\right]^{\tfrac{1}{4}}=x^{\tfrac{1}{4}}\cdot x^{\tfrac{1}{6}}\cdot x^{\tfrac{1}{8}} =x^{\tfrac{1}{4}+\tfrac{1}{6}+\tfrac{1}{8}}=x^{\tfrac{13}{24}}$.
		\end{itemize}
		\begin{itemchoice}
			\itemch Với $a=2$  ta có $A=2^{3+2\sqrt{2}}\approx 56{,}82< 57$.
			 \itemch Với $x=2$ ta có $B=2^{\tfrac{13}{24}}\approx1{,}46< 2$.
			\itemch Khi $x=a$ thì $A\cdot  B=a^{3+2\sqrt{2}}\cdot a^{\tfrac{13}{24}}=a^{3+2\sqrt{2}+\tfrac{13}{24}} =a^{\tfrac{85+48\sqrt{2}}{24}}$.
			\itemch Khi $x=a$  thì $A:B=a^{3+2\sqrt{2}}: a^{\tfrac{13}{24}}=a^{3+2\sqrt{2}-\tfrac{13}{24}}=a^{\tfrac{59+48\sqrt{2}}{24}}$.
		\end{itemchoice}
	} 
\end{ex}

% nguồn: Ok_dungsai_lop11_Lan2
\begin{ex}%[1D6V1-2]%[ma:ed1dfffc]
	Cho biểu  thức $K=\left(5^{-\tfrac{2}{3}} \right)^{-3}+\left[(0{,}2)^{\tfrac{3}{5}} \right]^{-5}$.
	 \choiceTF
		{\True $\left(5^{-\tfrac{2}{3}}  \right)^{-3}=5^2$}
		{\True $\left[(0{,}2)^{\tfrac{3}{5}}  \right]^{-5}=0{,}2)^{-3}$}
		{$K=5^m+5^n$ với $m,n$ là các số  tự nhiên chẵn}
		{$K$ chia hết  cho $4$}
	 \loigiai{
		\begin{itemchoice}
			\itemch Ta có $\left(5^{-\tfrac{2}{3}} \right)^{-3}=5^2$. 
			 \itemch Ta có $\left[(0{,}2)^{\tfrac{3}{5}} \right]^{-5}=(0{,}2)^{-3}$.
			\itemch Ta có $K=\left(5^{-\tfrac{2}{3}} \right)^{-3}+\left[(0{,}2)^{\tfrac{3}{5}} \right]^{-5}=5^2+(0{,}2)^{-3}=5^2+\left(\dfrac{1}{5}  \right)^{-3}=5^2+5^3$.\\
			Mà  $K=5^m+5^n$ nên $m=2$; $n=3$.
			\itemch Ta có $\left(5^{-\tfrac{2}{3}}  \right)^{-3}+\left[(0{,}2)^{\tfrac{3}{5}} \right]^{-5}=5^2+(0{,}2)^{-3}=5^2+\left(\dfrac{1}{5} \right)^{-3}=5^2+5^3=150$.\\
			Do đó $K$ không chia  hết cho $4$.
		\end{itemchoice}
	} 
\end{ex}

% nguồn: Ok_dungsai_lop11_Lan2
\begin{ex}%[1D6V1-2]%[ma:0ea366aa]
	Cho biểu thức $\sqrt[5]{2\cdot \sqrt[3]{2\cdot  \sqrt{2}}}=2^{\tfrac{a}{b}}$ và $\sqrt[6]{3\cdot \sqrt[3]{3\cdot \sqrt{3}}}=3^{\tfrac{m}{n}}$ trong đó ($\dfrac{a}{b}$, $\dfrac{m}{n}$ là các phân số tối giản).
	\choiceTF 
		{\True $a+b=13$} 
		{$m-n=3$} 
		{\True  $\dfrac{a}{b}+\dfrac{m}{n}=\dfrac{11}{20}$}
		{\True $\dfrac{a}{b}-\dfrac{m}{n}=\dfrac{1}{20}$} 
	 \loigiai{
		 Ta có
		\begin{itemize}
			\item $\sqrt[5]{2\cdot \sqrt[3]{2\cdot  \sqrt{2}}}=\sqrt[5]{2\cdot \sqrt[3]{2\cdot 2^{\tfrac{1}{2}}}}=\sqrt[5]{2\cdot \sqrt[3]{2^{\tfrac{3}{2}}}}=\sqrt[5]{2\cdot 2^{\tfrac{1}{2}}}=\sqrt[5]{2^{\tfrac{3}{2}}}=2^{\tfrac{3}{10}}$.\\
			Do đó $a=3$; $b=10$. 
			\item $\sqrt[6]{3\cdot \sqrt[3]{3\cdot \sqrt{3}}}=\sqrt[6]{3\cdot  \sqrt[3]{3\cdot 3^{\tfrac{1}{2}}}}=\sqrt[6]{3\cdot \sqrt[3]{3^{\tfrac{3}{2}}}}=\sqrt[6]{3\cdot 3^{\tfrac{1}{2}}}=\sqrt[6]{3^{\tfrac{3}{2}}}=3^{\tfrac{1}{4}}$.\\
			Do đó $m=1$; $n=4$. 
		\end{itemize}
		\begin{itemchoice}
			 \itemch Ta có $a+b=13$.
			\itemch Ta  có $m-n=-3$.
			\itemch Ta  có $\dfrac{a}{b}+\dfrac{m}{n}=\dfrac{3}{10}+\dfrac{1}{4}=\dfrac{11}{20}$.
			\itemch Ta có  $\dfrac{a}{b}-\dfrac{m}{n}=\dfrac{3}{10}-\dfrac{1}{4}=\dfrac{1}{20}$.
		\end{itemchoice}
	 }
\end{ex}

% nguồn: Ok_dungsai_lop11_Lan2
\begin{ex}%[1D6V1-2]%[ma:84665f28]
	Cho  các biểu thức $A=\sqrt{2\cdot \sqrt[3]{2\cdot \sqrt[4]{2}}}$, $B=\sqrt[24]{2^5}\cdot \dfrac{1}{\sqrt{2^{-1}}}$.
	\choiceTF 
		{\True Rút gọn biểu thức $A$, ta được $A=2^{\tfrac{m}{n}}$, với $\dfrac{m}{n}$ là phân số tối giản. Khi  đó, $m+n=41$}
		{Rút gọn biểu thức $B$, ta được $B=2^{\tfrac{p}{q}}$, với $\dfrac{p}{q}$ là phân số tối giản. Khi  đó, $p+q=31$}
		{Rút gọn biểu thức $A-B\sqrt{5}$, ta được  $A-B\sqrt{5}=\sqrt{5}$}
		{\True Rút  gọn biểu thức $A\cdot B$, ta được $A\cdot B=2^{\tfrac{r}{s}}$, với $\dfrac{r}{s}$ là phân số tối giản. Khi đó, $r+s=29$}
	\loigiai{ 
		\begin{itemchoice}
			\itemch Ta có $A=\sqrt{2\cdot \sqrt[3]{2\cdot \sqrt[4]{2}}}=\sqrt{2\sqrt[3]{2\cdot 2^{\tfrac{1}{4}}}}=\sqrt{2\cdot  2^{\tfrac{5}{12}}}=2^{\tfrac{17}{24}}$.\hfill $(1)$\\
			 Suy ra $m=17$, $n=24$. Do đó, $m+n=41$.
			\itemch Ta có $B=\sqrt[24]{2^5}\cdot \dfrac{1}{\sqrt{2^{-1}}}=2^{\tfrac{5}{24}}\cdot 2^{\tfrac{1}{2}}=2^{\tfrac{17}{24}}$.\hfill  $(2)$\\
			Suy ra $p=17$, $q=24$.  Do đó, $p+q=41$.
			\itemch Từ $(1)$  và $(2)$ suy ra $A=B$.\hfill $(3)$\\
			Do đó,  $A-B\sqrt{5}=\left(1-\sqrt{5}\right)\cdot A=\left(1-\sqrt{5}\right)\cdot 2^{\tfrac{17}{24}}\ne\sqrt{5}$.
			\itemch Từ $(3)$, ta có $A\cdot  B=A^2=2^{\tfrac{17}{12}}$.\\
			Suy ra  $r=17$, $s=12$. Do đó, $r+s=29$.
		\end{itemchoice}
	} 
\end{ex}

% nguồn: Ok_dungsai_lop11_Lan2
\begin{ex}%[1D6V1-2]%[ma:47fc1224]
	Cho hai biểu thức $P(x)=x^{\tfrac{4}{5}}\cdot\sqrt[6]{x^5\sqrt{x}}$ và $Q(x)=x^{\tfrac{4}{5}}:\sqrt[6]{x^5\sqrt{x}}$, với $x >  0$.
	\choiceTF 
		{$P(x)\cdot Q(x)=x^{\tfrac{11}{6}}$} 
		 {$P(x):Q(x)=x^{\tfrac{8}{5}}$}
		{$\left[P(2)\cdot  Q(2)\right]^5+\left[P(2):Q(2)\right]^6=2\,308$}
		{\True $\left[P(2)\cdot Q(2)\right]^5\cdot  \left[P(2)\colon Q(2)\right]^6=524\,288$}
	\loigiai{ 
		Với $x > 0, y  > 0$, ta có
		\begin{itemize}
			\item  $P(x)=x^{\tfrac{4}{5}}\cdot \sqrt[6]{x^5\sqrt{x}}=x^{\tfrac{4}{5}}\cdot x^{\tfrac{5}{6}}\cdot x^{\tfrac{1}{12}}=x^{\tfrac{4}{5}+\tfrac{5}{6}+\tfrac{1}{12}}=x^{\tfrac{103}{60}}$.
			\item  $Q(x)=x^{\tfrac{4}{5}}:\sqrt[6]{x^5\sqrt{x}}=\dfrac{x^{\tfrac{4}{5}}}{x^{\tfrac{5}{6}}\cdot x^{\tfrac{1}{12}}}=x^{\tfrac{4}{5}-\tfrac{5}{6}-\tfrac{1}{12}}=x^{\tfrac{-7}{60}}$.
		\end{itemize}
		\begin{itemchoice}
			\itemch $P(x)\cdot  Q(x)=x^{\tfrac{8}{5}}$
			 \itemch $P(x):Q(x)=x^{\tfrac{11}{6}}$.
			\itemch  $\left[P(2)\cdot Q(2)\right]^5+\left[P(2): Q(2)\right]^6=\left(2^{\tfrac{8}{5}} \right)^5+\left(2^{\tfrac{11}{6}} \right)^6=2^8+2^{11}=2\,304$.
			 \itemch $\left[P(2)\cdot Q(2)\right]^5\cdot\left[P(2): Q(2)\right]^6=\left(2^{\tfrac{8}{5}} \right)^5\cdot\left(2^{\tfrac{11}{6}} \right)^6=2^8\cdot 2^{11}=2^{19}=524\,288$.
		\end{itemchoice}
	} 
\end{ex}

% nguồn: Ok_dungsai_lop11_Lan2
\begin{ex}%[1D6V1-2]%[ma:d9fb920c]
	Cho biểu thức $A=\left(\dfrac{a^{\sqrt{5}}}{b^{\sqrt{5}-2}} \right)^{\sqrt{5}+2} \cdot \dfrac{a^{-2-\sqrt{5}}}{b^{-1}}$  với $a,b > 0$.
	 \choiceTF
		{Sau khi rút gọn, thì biểu thức $A$ chỉ  chứa biến $b$}
		{Với $a=2$,  $b=1+5\sqrt{2}$ thì $A=\dfrac{113}{3}$}
		{\True Khi  $A=a^m\cdot b^n$ thì $m+n=3+\sqrt{5}$}
		{Khi $A=a^m\cdot b^n$  thì $m-n=2+\sqrt{5}$}
	\loigiai{ 
		Ta có $A=\dfrac{\left(a^{\sqrt{5}} \right)^{\sqrt{5}+2}}{\left(b^{\sqrt{5}-2} \right)^{\sqrt{5}+2}}  \cdot \dfrac{b}{a^{2+\sqrt{5}}}=\dfrac{a^{5+2\sqrt{5}} \cdot b}{b\cdot a^{2+\sqrt{5}}}=a^{3+\sqrt{5}}$.
		\begin{itemchoice}
			\itemch Ta  có $A=a^{3+\sqrt{5}}$. \\
			Do đó, sau khi rút gọn, thì biểu thức $A$  chỉ chứa biến $a$.
			\itemch Với $a=2$, $b=1+5\sqrt{2}$  thì $A=2^{3+\sqrt{5}}$.
			\itemch Ta có $A=a^{3+\sqrt{5}}$. Suy ra  $m=3+\sqrt{5}$, $n=0$.\\
			 Do đó $m+n=3+\sqrt{5}$.
			\itemch  Ta có $A=a^{3+\sqrt{5}}$. Suy ra $m=3+\sqrt{5}$, $n=0$.\\
			Do đó $m-n=3+\sqrt{5}$. 
		\end{itemchoice}
	 }
\end{ex}

% nguồn: Ok_dungsai_lop11_Lan2
\begin{ex}%[1D6V1-2]%[ma:12cfde7d]
	Cho biểu thức $P(a)=\dfrac{a^{-\tfrac{1}{3}} \left(\sqrt[3]a-\sqrt[3]{a^4} \right)}{a^{\tfrac{1}{8}} \left(\sqrt[8]{a^3}-\sqrt[8]{a^{-1}} \right)}$ với $a > 0$, $a\ne  1$.
	\choiceTF 
		{\True  $P(a)+\sqrt{a}=-1$}
		{\True  $\left(1-\sqrt{3} \right)\cdot P(3)=2$}
		{$P\left(2\,023^{2\,024} \right)+2\,023^{1\,012}+10=1$} 
		{$\left(2^{1\,012}-1\right)\cdot P\left(2^{2\,024} \right)=2^{2\,024}-1$} 
	 \loigiai{
		Ta  có \allowdisplaybreaks
		\begin{eqnarray*}
			P(a)&=&\dfrac{a^{-\tfrac{1}{3}} \left(\sqrt[3]a-\sqrt[3]{a^4} \right)}{a^{\tfrac{1}{8}} \left(\sqrt[8]{a^3}-\sqrt[8]{a^{-1}}  \right)}\\
			&=&\dfrac{a^{-\tfrac{1}{3}} \left(a^{\tfrac{1}{3}}-a^{\tfrac{4}{3}}  \right)}{a^{\tfrac{1}{8}} \left(a^{\tfrac{3}{8}}-a^{-\tfrac{1}{8}} \right)}\\
			 &=&\dfrac{1-a}{a^{\tfrac{1}{2}}-1}\\
			 &=&-a^{\tfrac{1}{2}}-1\\
			&=&-\sqrt{a}-1.\quad(*) 
		\end{eqnarray*}
		\begin{itemchoice}
			\itemch Ta  có $P(a)+\sqrt{a}=-\sqrt{a}-1+\sqrt{a}=-1$.
			\itemch Với  $a=3$ thế vào $(*)$ ta được \allowdisplaybreaks
			\begin{eqnarray*}
				\left(1-\sqrt{3} \right)P(3)&=&\left(1-\sqrt{3} \right)\cdot  \left(-\sqrt{3}-1\right)\\
				 &=&2.
			\end{eqnarray*}
			\itemch Với $a=2\,023^{2\,024}$ thế  vào $(*)$ ta được
			\allowdisplaybreaks 
			\begin{eqnarray*}
				P\left(2\,023^{2\,024}  \right)+2\,023^{1\,012}+10&=&-\sqrt{2\,023^{2\,024}}-1+2023^{1\,012}+10\\
				 &=&-2\,023^{1\,012}-1+2\,023^{1\,012}+10\\
				&=&9. 
			\end{eqnarray*}
			\itemch Với $a=2^{2\,024}$ thế vào  $(*)$ ta được
			\allowdisplaybreaks 
			\begin{eqnarray*}
				 \left(2^{1\,012}-1\right)\cdot P\left(2^{2\,024} \right)&=&\left(2^{1\,012}-1\right)\left(-\sqrt{2^{2\,024}}-1\right)\\
				&=&\left(2^{1\,012}-1\right)\left(-2^{1\,012}-1\right)\\ 
				 &=&-\left(2^{1\,012}-1\right)\left(2^{1\,012}+1\right)\\
				 &=&1-2^{2\,024}.
			\end{eqnarray*}
		\end{itemchoice}
	 }
\end{ex}

% nguồn: Ok_dungsai_lop11_Lan2
\begin{ex}%[1D6V1-2]%[ma:bc9fdbd5]
	Cho biểu  thức $Q=\sqrt{\sqrt{5} \cdot \left(\sqrt[4]{\sqrt{5}}:\sqrt{\sqrt[5]5}\right)^{10}}$.
	 \choiceTF
		{\True $\sqrt[4]{\sqrt{5}}=5^{\tfrac{1}{8}}$} 
		{\True $\sqrt[4]{\sqrt{5}}:\sqrt{\sqrt[5]5}=5^{\tfrac{a}{b}}$, $\left(\dfrac{a}{b}\text{ là phân số tối giản}\right)$, khi đó  $a+b=41$}
		{$\sqrt{5} \cdot \left(\sqrt[4]{\sqrt{5}}:\sqrt{\sqrt[5]5}\right)^{10}=5^{\tfrac{c}{d}}$ $\left(\dfrac{c}{d}\text{ là phân số tối giản}\right)$, khi  đó $c+d=6$}
		{$Q=5^{\tfrac{m}{n}}$, $\left(\dfrac{m}{n}\text{ là phân số tối giản}\right)$, khi đó  $m+n=12$}
	 \loigiai{
		\begin{itemchoice}
			\itemch Ta có $\sqrt[4]{\sqrt{5}}=\left(5^{\tfrac{1}{2}}\right)^{\tfrac{1}{4}}=5^{\tfrac{1}{8}}$. 
			\itemch Ta có  \allowdisplaybreaks
			\begin{eqnarray*}
				 \sqrt[4]{\sqrt{5}}:\sqrt{\sqrt[5]5}&=&\left(5^{\tfrac{1}{2}}\right)^{\tfrac{1}{4}}:\left(5^{\tfrac{1}{5}}\right)^{\tfrac{1}{2}}\\
				 &=&5^{\tfrac{1}{8}}:5^{\tfrac{1}{10}}\\
				&=&5^{\tfrac{1}{8}-\tfrac{1}{10}}\\ 
				&=&5^{\tfrac{1}{40}}. 
			\end{eqnarray*}
			 Mà $\sqrt[4]{\sqrt{5}}:\sqrt{\sqrt[5]5}=5^{\tfrac{a}{b}}$ nên $a=1$, $b=40$.\\
			Do đó $a+b=41$. 
			\itemch  Ta có \allowdisplaybreaks
			\begin{eqnarray*}
				\sqrt{5}  \cdot \left(\sqrt[4]{\sqrt{5}}:\sqrt{\sqrt[5]5}\right)^{10}
				&=&5^{\tfrac{1}{2}}\cdot\left[\left(5^{\tfrac{1}{2}}\right)^{\tfrac{1}{4}}:\left(5^{\tfrac{1}{5}}\right)^{\tfrac{1}{2}}\right]\\ 
				 &=&5^{\tfrac{1}{2}}\cdot\left[5^{\tfrac{1}{8}}:5^{\tfrac{1}{10}}\right]\\
				&=&5^{\tfrac{1}{2}}\cdot 5^{\tfrac{1}{8}-\tfrac{1}{10}}\\ 
				 &=&5^{\tfrac{1}{2}}\cdot 5^{\tfrac{1}{40}}\\
				&=&5^{\tfrac{1}{2}+\tfrac{1}{40}}\\ 
				 &=&5^{\tfrac{21}{40}}.
			\end{eqnarray*}
			Mà $\sqrt{5} \cdot \left(\sqrt[4]{\sqrt{5}}:\sqrt{\sqrt[5]5}\right)^{10}=5^{\tfrac{c}{d}}$ nên  $c=21$; $d=40$.\\
			Do  đó $c+d=61$.
			\itemch Ta có  \allowdisplaybreaks
			\begin{eqnarray*}
				 Q&=&\sqrt{5^{\tfrac{1}{2}} \left(5^{\tfrac{1}{8}}:5^{\tfrac{1}{10}} \right)^{10}}\\
				&=&\sqrt{5^{\tfrac{1}{2}} \left(5^{\tfrac{1}{40}}  \right)^{10}}\\
				&=&\sqrt{5^{\tfrac{1}{2}} \cdot 5^{\tfrac{1}{4}}}\\ 
				 &=&\sqrt{5^{\tfrac{3}{4}}}\\
				&=&5^{\tfrac{3}{8}}. 
			\end{eqnarray*}
			Mà $Q=5^{\tfrac{m}{n}}$ nên $m=3$;  $n=8$.\\
			Do đó  $m+n=11$.
		\end{itemchoice}
	 }
\end{ex}

% nguồn: Ok_dungsai_lop11_Lan2
\begin{ex}%[1D6V1-2]%[ma:f9e9ceb2]
	Cho biểu thức $A=\left(\dfrac{a^{\sqrt{5}}}{b^{\sqrt{5}-2}}\right)^{\sqrt{5}+2}\cdot \dfrac{a^{-2-\sqrt{5}}}{b^{-1}}$, với $a$, $b$ là các số thực dương tùy ý. 
	\choiceTF 
		{Sau khi rút gọn, thì  biểu thức $A$ chỉ chứa biến $b$}
		 {Với $a=2$, $b=1+5\sqrt{2}$ thì $A=\dfrac{113}{3}$}
		{\True Khi $A=a^m\cdot b^n$ thì $m+n=3+\sqrt{5}$} 
		{Khi $A=a^m\cdot b^n$  thì $m-n=2+\sqrt{5}$}
	 \loigiai{
		\begin{itemchoice}
			 \itemch Ta có $A=\left(\dfrac{a^{\sqrt{5}}}{b^{\sqrt{5}-2}}\right)^{\sqrt{5}+2}\cdot \dfrac{a^{-2-\sqrt{5}}}{b^{-1}}=\dfrac{a^{\sqrt{5}(\sqrt{5}+2)}}{b^{(\sqrt{5}-2)(\sqrt{5}+2)}}\cdot \dfrac{a^{-2-\sqrt{5}}}{b^{-1}}=\dfrac{a^{5+2\sqrt{5}-2-\sqrt{5}}}{b^{5-4-1}}=a^{3+\sqrt{5}}$.
			 \itemch Với $a=2$, $b=1+5\sqrt{2}$ thì $A=2^{3+\sqrt{5}}\ne\dfrac{113}{3}$.
			\itemch Vì $A=a^{3+\sqrt{5}} b^\circ=a^m\cdot b^n$ nên $\heva{& m=3+\sqrt{5}\\ & n=0 
			}$ $\Rightarrow$ $m+n=3+\sqrt{5}$. 
			\itemch Vì $A=a^{3+\sqrt{5}}b^\circ=a^m\cdot b^n$  nên $\heva{& m=3+\sqrt{5}\\ & n=0
			}$ $\Rightarrow$  $m-n=3+\sqrt{5}$.
		\end{itemchoice}
	} 
\end{ex}

% nguồn: Ok_dungsai_lop11_Lan2
\begin{ex}%[1D6V1-2]%[ma:d6c4c9f0]
	Cho biểu thức $Q=\left(\sqrt{x}-\sqrt[4]{x}+1\right)\left(\sqrt{x}+\sqrt[4]{x}+1\right)\left(x-\sqrt{x}+1\right)$, với  $x\geq 0$.
	 \choiceTF
		{\True Khi $x=2$  thì $Q=7$}
		{Giá trị $x=1+\sqrt{2}$ là một nghiệm của  phương trình $Q=0$}
		{Phương trình $Q=3$  có $2$ nghiệm dương phân biệt}
		{\True Khi $x=3$  thì $Q$ là một số nguyên tố}
	\loigiai{ 
		Với $x\ge0$, ta  có\allowdisplaybreaks
		\begin{align*}
			Q  &=\left(\sqrt{x}+1-\sqrt[4]{x}\right)\left(\sqrt{x}+1+\sqrt[4]{x}\right)\left(x-\sqrt{x}+1\right)\\
			&=\left[\left(\sqrt{x}+1\right)^2-\left(\sqrt[4]{x}\right)^2\right]\left(x+1-\sqrt{x}\right)\\ 
			&=\left[\left(x+2\sqrt{x}+1\right)-\sqrt{x}\right]\left(x+1-\sqrt{x}\right)\\ 
			&=\left(x+1+\sqrt{x}\right)\left(x+1-\sqrt{x}\right)\\ 
			 &=(x+1)^2-\left(\sqrt{x}\right)^2\\
			&=x^2+x+1. 
		\end{align*}
		Do đó, 
		\begin{itemchoice}
			\itemch Khi $x=2$, ta được  $Q=2^2+2+1=7$.
			\itemch Với mọi $x\ge0$, ta có $Q=x^2+x+1\ge1>0$ nên phương  trình $Q=0$ vô nghiệm.
			\itemch Ta có $Q=0\Leftrightarrow x^2+x+1=3\Leftrightarrow x^2+x-2=0\Leftrightarrow\hoac{&x=1>0\\ &x=-2<0\quad (\text{loại}).}$\\ 
			Như vậy, phương trình $Q=3$ có đúng $1$ nghiệm  dương.
			\itemch Khi $x=3$, ta được $Q=3^2+3+1=13$  là một số nguyên tố.
		\end{itemchoice}
	} 
\end{ex}

% nguồn: Ok_dungsai_lop11_Lan2
\begin{ex}%[1D6V1-2]%[ma:9b5a73d2]
	Cho biểu thức $P(x)=\sqrt[3]{x\sqrt[4]{x^3\sqrt{x}}}$, với  $x > 0$.
	 \choiceTF
		{\True  $P(x)=x^{\tfrac{5}{8}}$}
		{$P\left(\sqrt{3}-1\right) < P\left(\sqrt{2}-1\right)$} 
		{$P\left(\sqrt[5]{256} \right)\ne  2$}
		{\True $P(3)\cdot P(3)\cdot P(9)\cdot P(27)\cdot P(81)\cdot P(243)=59\,049$} 
	 \loigiai{
		\begin{itemchoice}
			\itemch Ta có\allowdisplaybreaks 
			\begin{eqnarray*}
				 P&=&\sqrt[3]{x\sqrt[4]{x^3\sqrt{x}}}\\
				 &=&\left[x(x^3\cdot x^{\tfrac{1}{2}})^{\tfrac{1}{4}}\right]^{\tfrac{1}{3}}\\
				&=& \left[x(x^{\tfrac{7}{2}})^{\tfrac{1}{4}}\right]^{\tfrac{1}{3}}\\ 
				 &=& x^{\tfrac{1}{3}}\cdot x^{\tfrac{7}{24}}=x^{\tfrac{5}{8}}.
			\end{eqnarray*}
			\itemch Ta có $\dfrac{5}{8} > 0$;  $\sqrt{3}-1> \sqrt{2}-1$.\\
			Do đó $\left(\sqrt{3}-1\right)^{\tfrac{5}{8}} > \left(\sqrt{2}-1\right)^{\tfrac{5}{8}}$.\\ 
			Vậy  $P\left(\sqrt{3}-1\right) > P\left(\sqrt{2}-1\right)$.
			\itemch  Ta có $P\left(\sqrt[5]{256} \right)=\left(\sqrt[5]{256} \right)^{\tfrac{5}{8}}=\left(2^{\tfrac{8}{5}} \right)^{\tfrac{5}{8}}=2^{\tfrac{8}{5}\cdot \tfrac{5}{8}}=2$.
			\itemch Ta có  \allowdisplaybreaks
			\begin{eqnarray*}
				P(3)\cdot P(3)\cdot P(9)\cdot P(27)\cdot  P(81)\cdot P(243)&=&3^{\tfrac{5}{8}}\cdot 3^{\tfrac{5}{8}}\cdot 9^{\tfrac{5}{8}}\cdot 27^{\tfrac{5}{8}}\cdot 81^{\tfrac{5}{8}}\cdot 243^{\tfrac{5}{8}} \\
				&=&\left(3\cdot  3\cdot 9\cdot 27\cdot 81\cdot 243\right)^{\tfrac{5}{8}}=\left(3^{16} \right)^{\tfrac{5}{8}}\\
				&=&3^{16\cdot  \tfrac{5}{8}}\\
				&=&3^{10}\\ 
				 &=&59\,049.
			\end{eqnarray*}
		\end{itemchoice}
	} 
\end{ex}

% nguồn: Ok_dungsai_lop11_Lan2
\begin{ex}%[1D6V1-2]%[ma:13079b5d]
	Cho biểu thức $P=81^{-0{,}75}+\left(\dfrac{1}{625} \right)^{-\tfrac{1}{4}}-\left(\dfrac{1}{32}  \right)^{-\tfrac{3}{5}}$.
	\choiceTF 
		{\True $81^{-0{,}75}=(3^4)^{-\tfrac{3}{4}}$} 
		{$\left(\dfrac{1}{625}  \right)^{-\tfrac{1}{4}}=\left(5^{-4} \right)^{\tfrac{1}{4}}$}
		{\True $P=3^m+5-2^n$, với  $m+n=0$}
		{$P=-\dfrac{a}{b}$, $\left(a,b\in \mathbb{N}^*\right)$ và  $\dfrac{a}{b}$ là phân số tối giản, khi đó $a-b=52$}
	 \loigiai{
		\begin{itemchoice}
			\itemch  Ta có $81^{-0{,}75}=(3^4)^{-\tfrac{3}{4}}$.
			\itemch Ta  có $\left(\dfrac{1}{625} \right)^{-\tfrac{1}{4}}=\left(5^{-4} \right)^{-\tfrac{1}{4}}$.
			 \itemch Ta có \allowdisplaybreaks
			\begin{eqnarray*}
				P&=&81^{-0{,}75}+\left(\dfrac{1}{625} \right)^{-\tfrac{1}{4}}-\left(\dfrac{1}{32} \right)^{-\tfrac{3}{5}}\\ 
				&=&(3^4)^{-\tfrac{3}{4}}+\left(5^{-4} \right)^{-\tfrac{1}{4}}-\left(2^{-5} \right)^{-\tfrac{3}{5}}\\ 
				 &=&3^{-3}+5-2^3.
			\end{eqnarray*}
			Mà  $P=3^m+5-2^n$ nên $m=-3$; $n=3$. \\
			Do  đó $m+n=0$.
			 \itemch Ta có \allowdisplaybreaks
			\begin{eqnarray*}
				P&=&81^{-0{,}75}+\left(\dfrac{1}{625} \right)^{-\tfrac{1}{4}}-\left(\dfrac{1}{32} \right)^{-\tfrac{3}{5}}\\ 
				&=&(3^4)^{-\tfrac{3}{4}}+\left(5^{-4} \right)^{-\tfrac{1}{4}}-\left(2^{-5} \right)^{-\tfrac{3}{5}}\\ 
				&=&3^{-3}+5-2^3\\ 
				&=&-\dfrac{80}{27}. 
			\end{eqnarray*}
			Mà $P=-\dfrac{a}{b}$  nên $a=80$; $b=27$.\\
			 Do đó $a-b=53$.
		\end{itemchoice}
	} 
\end{ex}

% nguồn: Ok_dungsai_lop11_Lan2
\begin{ex}%[1D6V1-2]%[ma:737abc97]
	 Tính giá trị của biểu thức $P=\left(7+4\sqrt{3}\right)^{2\,025} \left(4\sqrt{3}-7\right)^{2\,024}$.
	 \choiceTF
		{\True  $\left(7+4\sqrt{3}\right)^{2\,025}=\left(7+4\sqrt{3} \right)^{2\,024}\cdot\left(7+4\sqrt{3} \right)$}
		{$P=\left[\left(4\sqrt{3}+7\right)\left(4\sqrt{3}-7\right)\right]^{2\,024}$} 
		{\True  $P=\left(7+4\sqrt{3} \right)\left(-1\right)^{2\,024}$}
		 {$P=7-4\sqrt{3}$}
	 \loigiai{
		Ta  có \allowdisplaybreaks
		\begin{eqnarray*}
			 P&=&\left(7+4\sqrt{3} \right)\cdot\left[\left(4\sqrt{3}+7\right)\left(4\sqrt{3}-7\right)\right]^{2\,024}\\
			&=&\left(7+4\sqrt{3} \right)\cdot\left[\left(4\sqrt{3} \right)^2-7^2\right]^{2\,024}\\ 
			 &=&\left(7+4\sqrt{3} \right)\cdot\left(-1\right)^{2\,024}\\
			 &=&7+4\sqrt{3}.
		\end{eqnarray*}
		Do  đó
		\begin{itemchoice}
			 \itemch
			\itemch 
			 \itemch
			 \itemch
		\end{itemchoice}
	 }
\end{ex}

% nguồn: Ok_dungsai_lop11_Lan2
\begin{ex}%[1D6V1-2]%[ma:3bc570be]
	Cho biểu thức $B=\dfrac{a^{\tfrac{1}{3}} \left(a^{\tfrac{1}{2}}-a^{\tfrac{5}{2}} \right)}{a^{\tfrac{1}{4}}  \left(a^{\tfrac{7}{12}}-a^{\tfrac{19}{12}} \right)}$ với $a > 0$.
	\choiceTF 
		{\True Sau khi rút gọn, thì biểu thức $B=k+a^m$, $k,m >  0$}
		{Với $a=4+\sqrt{5}$ thì $B=6+3\sqrt{5}$} 
		{\True  Với $a=3+\sqrt{5}$ thì $B=4+\sqrt{5}$}
		{Với $a=m+1$  thì $B=3+2m$}
	 \loigiai{
		Ta có \allowdisplaybreaks 
		\begin{eqnarray*}
			B&=&\dfrac{a^{\tfrac{1}{3}}\cdot a^{\tfrac{1}{2}} \left(1-a^{2}  \right)}{a^{\tfrac{1}{4}}\cdot a^{\tfrac{7}{12}} \left(1-a\right)}\\
			&=&\dfrac{a^{\tfrac{5}{6}}\cdot(1-a)(1+a)}{a^{\tfrac{5}{6}}\cdot  (1-a)}\\
			 &=&1+a.
		\end{eqnarray*}
		\begin{itemchoice}
			 \itemch Ta có $B=1+a$.\\
			Do  đó, sau khi rút gọn, thì biểu thức $B=1+a$.
			\itemch Với  $a=4+\sqrt{5}$ thì $B=1+4+\sqrt{5}=5+\sqrt{5}$.
			\itemch Với $a=3+\sqrt{5}$ thì $B=1+3+\sqrt{5}=4+\sqrt{5}$. 
			\itemch Với $a=m+1$ thì $B=1+m+1=2+m$. 
		\end{itemchoice}
	 }
\end{ex}

% nguồn: Ok_dungsai_lop11_Lan2
\begin{ex}%[1D6V1-2]%[ma:20c6154e]
	Cho các biểu  thức
	$A=\dfrac{\sqrt{a}+\sqrt[4]{ab}}{\sqrt[4]{a}+\sqrt[4]{b}}-\dfrac{\sqrt{a}-\sqrt{b}}{\sqrt[4]{a}-\sqrt[4]{b}}$, 
	$B=-\sqrt[4]{b}$; với $a>0$, $b>0$, $a\ne  b$. Khi đó
	 \choiceTF
		{\True  Sau khi rút gọn, thì $A$ chỉ chứa biến $b$}
		{Biểu  thức luôn $A>0$}
		 {$A=B+\sqrt{a}$}
		{\True $\dfrac{A-1}{B}=1-\dfrac{1}{B}$} 
	 \loigiai{
		\begin{itemchoice}
			\itemch Ta  có
			\begin{align*}
				\dfrac{\sqrt{a}+\sqrt[4]{ab}}{\sqrt[4]{a}+\sqrt[4]{b}} 
				-\dfrac{\sqrt{a}-\sqrt{b}}{\sqrt[4]{a}-\sqrt[4]{b}} 
				 &=\dfrac{\left(\sqrt[4]{a}\right)^2+\sqrt[4]{ab}}{\sqrt[4]{a}+\sqrt[4]{b}}
				-\dfrac{\left(\sqrt[4]{a}\right)^2-\left(\sqrt[4]{b}\right)^2}{\sqrt[4]{a}-\sqrt[4]{b}}\\ 
				&=\dfrac{\sqrt[4]{a}\left(\sqrt[4]{a}+\sqrt[4]{b}\right)}{\sqrt[4]{a}+\sqrt[4]{b}} 
				 -\dfrac{\left(\sqrt[4]{a}-\sqrt[4]{b}\right)\left(\sqrt[4]{a}+\sqrt[4]{b}\right)}{\sqrt[4]{a}-\sqrt[4]{b}}\\
				&=\sqrt[4]{a}-\left(\sqrt[4]{a}+\sqrt[4]{b}\right)\\ 
				&=-\sqrt[4]{b}. 
			\end{align*}
			Vậy sau khi rút gọn, thì $A$ chỉ chứa biến $b$. 
			\itemch Ta có $A=-\sqrt[4]{b}<0$, $\forall b>0$. 
			 \itemch $A=B=-\sqrt[4]{b}$.
			 \itemch Vì $A=B$ nên
			 \[
			\dfrac{A-1}{B}=\dfrac{B-1}{B}=1-\dfrac{1}{B}. 
			 \]
		\end{itemchoice}
	} 
\end{ex}

% nguồn: Ok_dungsai_lop11_Lan2
\begin{ex}%[1D6V1-2]%[ma:df6e9ba2]
	Cho số thực  $a$ dương.
	 \choiceTF
		 {$a^{\tfrac{3}{2}}\cdot a^{\tfrac{2}{9}}=\sqrt[3]a$}
		{\True $a^{\tfrac{3}{2}}\cdot\sqrt{a}=a^2$} 
		{\True $\sqrt[3]a\cdot\sqrt[6]a=\sqrt{a}$} 
		{$a^{\sqrt{5}}\cdot\left(\dfrac{1}{a}  \right)^{\sqrt{5}-2}=a^{-2}$}
	\loigiai{ 
		\begin{itemchoice}
			 \itemch
			 $a^{\tfrac{3}{2}}\cdot a^{\tfrac{2}{9}}=a^{\tfrac{3}{2}+\tfrac{2}{9}}=a^{\tfrac{31}{18}}$.
			 \itemch
			$a^{\tfrac{3}{2}}\cdot\sqrt{a}=a^{\tfrac{3}{2}}\cdot  a^{\tfrac{1}{2}}=a^{\tfrac{3}{2}+\tfrac{1}{2}}=a^2$.
			 \itemch
			$\sqrt[3]a\cdot\sqrt[6]a=a^{\tfrac{1}{3}}\cdot a^{\tfrac{1}{6}}=a^{\tfrac{1}{2}}=\sqrt{a}$. 
			 \itemch
			Ta có $a^{\sqrt{5}}  \left(\dfrac{1}{a} \right)^{\sqrt{5}-2}=a^{\sqrt{5}}\cdot a^{-\left(\sqrt{5}-2\right)}=a^{\sqrt{5}}\cdot a^{2-\sqrt{5}}=a^{\sqrt{5}+2-\sqrt{5}}=a^2$.
		\end{itemchoice}
	} 
\end{ex}

% nguồn: Ok_dungsai_lop11_Lan2
\begin{ex}%[1D6V1-2]%[ma:b2de358b]
	Cho các biểu thức $A=\dfrac{\sqrt{a}+\sqrt[4]{ab}}{\sqrt[4]{a}+\sqrt[4]{b}}-\dfrac{\sqrt{a}-\sqrt{b}}{\sqrt[4]{a}-\sqrt[4]{b}}$, $B=-\sqrt[4]{b}$, với $a>0$, $b>0$, $a\neq  b$.
	\choiceTF 
		{\True Sau khi rút gọn, thì $A$ chỉ  chứa biến $b$}
		{Biểu  thức $A$ luôn dương}
		 {$A=B+\sqrt{a}$}
		{\True $A-\sqrt{b}<B$} 
	\loigiai{ 
		\begin{itemchoice}
			\itemch Với $a>0$, $b>0$, $a\neq  b$, ta có
			\allowdisplaybreaks 
			\begin{align*}
				 A&=\dfrac{\sqrt{a}+\sqrt[4]{ab}}{\sqrt[4]{a}+\sqrt[4]{b}}-\dfrac{\sqrt{a}-\sqrt{b}}{\sqrt[4]{a}-\sqrt[4]{b}}\\
				 &=\dfrac{\left(\sqrt[4]{a}\right)^2+\sqrt[4]{ab}}{\sqrt[4]{a}+\sqrt[4]{b}}-\dfrac{\left(\sqrt[4]{a}\right)^2-\left(\sqrt[4]{b}\right)^2}{\sqrt[4]{a}-\sqrt[4]{b}}\\
				 &=\dfrac{\sqrt[4]{a}\left(\sqrt[4]{a}+\sqrt[4]{b}\right)}{\sqrt[4]{a}+\sqrt[4]{b}}-\dfrac{\left(\sqrt[4]{a}-\sqrt[4]{b}\right)\left(\sqrt[4]{a}+\sqrt[4]{b}\right)}{\sqrt[4]{a}-\sqrt[4]{b}}\\
				 &=\sqrt[4]{a}-\left(\sqrt[4]{a}+\sqrt[4]{b}\right)=-\sqrt[4]{b}.\tag{1}
			\end{align*}
			\itemch Từ  $(1)$ suy ra $A<0$.
			 \itemch Từ $(1)$ suy ra $A=B\ne B+\sqrt{a}$, $\forall a>0$.
			\itemch Ta có  $A-\sqrt{b}=B-\sqrt{b}<B$, $\forall b>0$.
		\end{itemchoice}
	 }
\end{ex}

% nguồn: Ok_dungsai_lop11_Lan2
\begin{ex}%[1D6V1-2]%[ma:044772d3]
	Cho biểu thức  $A=\left(\sqrt{2}-1\right)^{2 x}+\left(3+2\sqrt{2}\right)^x$.
	 \choiceTF
		{\True Nếu  $\left(\sqrt{2}+1\right)^x=3$ thì $A=\dfrac{82}{9}$}
		{\True Nếu $\left(\sqrt{2}+1\right)^x=1$ thì $A=2$} 
		{Nếu $\left(\sqrt{2}+1\right)^x=2$ thì  $A=\dfrac{17}{9}$}
		{\True Nếu $\left(\sqrt{2}+1\right)^x=4$ thì $A=\dfrac{257}{16}$} 
	 \loigiai{
		 Ta có\allowdisplaybreaks
		\begin{align*}
			 &\sqrt{2}-1=\dfrac{\left(\sqrt{2}-1\right)\left(\sqrt{2}+1\right)}{\sqrt{2}+1}=\dfrac{1}{\sqrt{2}+1}=\left(\sqrt{2}+1\right)^{-1};\\
			&3+2\sqrt{2}=\left(\sqrt{2}\right)^2+2\sqrt{2}+1=\left(\sqrt{2}+1\right)^2. 
		\end{align*}
		Do  đó,\allowdisplaybreaks
		\begin{align*}
			 A&=\left[\left(\sqrt{2}+1\right)^{-1}\right]^{2x}+\left[\left(\sqrt{2}+1\right)^2\right]^x\\
			&=\left(\sqrt{2}+1\right)^{-2x}+\left(\sqrt{2}+1\right)^{2x}\\ 
			&=\left[\left(\sqrt{2}+1\right)^x\right]^{-2}+\left[\left(\sqrt{2}+1\right)^x\right]^2. 
		\end{align*}
		\begin{itemchoice}
			\itemch Nếu  $\left(\sqrt{2}+1\right)^x=3$ thì $A=3^{-2}+3^2=\dfrac{82}{9}$.
			\itemch Nếu $\left(\sqrt{2}+1\right)^x=1$ thì $A=1^{-2}+1^2=2$. 
			\itemch  Nếu $\left(\sqrt{2}+1\right)^x=2$ thì $A=2^{-2}+2^2=\dfrac{17}{4}$.
			 \itemch Nếu $\left(\sqrt{2}+1\right)^x=4$ thì $A=4^{-2}+4^2=\dfrac{257}{16}$.
		\end{itemchoice}
	} 
\end{ex}


%%%%%%%%%%%%%%%% TRẢ LỜI NGẮN %%%%%%%%%%%%%%%%

% nguồn: Ok_TraLoiNgan_Lop11_Lan1
\begin{ex}%[1D6H1-2]%[ma:1fe6b886]
	Cho $x$,  $y$ là hai số nguyên thỏa mãn $3^x \cdot 6^y=\dfrac{2^{15} \cdot 6^{40}}{9^{50} \cdot 12^{25}}$. Tính $x y$.
	 \shortans{-450}
	\loigiai{ 
		 Ta có
		\[VT=3^x \cdot 6^y=3^x \cdot(2 \cdot 3)^y=2^y \cdot  3^{x+y}.\]
		Lại có \[VP=\dfrac{2^{15} \cdot 6^{40}}{9^{50} \cdot 12^{25}}=\dfrac{2^{15} \cdot(2 \cdot 3)^{40}}{\left(3^2\right)^{50} \cdot\left(2^2 \cdot 3\right)^{25}}=\dfrac{2^{15} \cdot 2^{40} \cdot 3^{40}}{3^{100}  \cdot 2^{50} \cdot 3^{25}}=\dfrac{2^{55} \cdot 3^{40}}{2^{50} \cdot 3^{125}}=2^5 \cdot 3^{-85}.\]
		Suy ra \[2^y \cdot 3^{x+y}=2^5 \cdot 3^{-85} \Leftrightarrow \heva{&y=5 \\&x+y=-8  5} \Leftrightarrow \heva{&y=5 \\
		&x=-90.}\] 
		Vậy $x \cdot  y=5 \cdot(-90)=-450$.
	} 
\end{ex}

% nguồn: Ok_TraLoiNgan_Lop11_Lan1
\begin{ex}%[1D6H1-2]%[ma:480d4bdd]
	Rút gọn biểu thức $E=\dfrac{3^{\sqrt{5}+2} \cdot 3^{1-\sqrt{5}}}{\left(3^{\sqrt{2}+2}\right)^{\sqrt{2}-2}}$ ta được kết quả là  bao nhiêu?
	\shortans{243} 
	\loigiai{ 
		Ta  có $E=\dfrac{3^{\sqrt{5}+2} \cdot 3^{1-\sqrt{5}}}{\left(3^{\sqrt{2}+2}\right)^{\sqrt{2}-2}}=\dfrac{3^3}{3^{-2}}=3^5=243$.
	 }
\end{ex}

% nguồn: Ok_TraLoiNgan_Lop11_Lan1
\begin{ex}%[1D6H1-2]%[ma:546d32f3]
	Cho  $x$, $y$ là các số thực dương. Giả sử $\left(\sqrt[5]{x^3} \cdot y^{\tfrac{2}{5}}\right)^2=x^a \cdot y^b$ với $a$, $b$ là số hữu tỷ. Tính $a+b$.
	\shortans{2} 
	\loigiai{ 
		Ta có $\left(\sqrt[5]{x^3} \cdot y^{\tfrac{2}{5}}\right)^2=\left(x^{\tfrac{3}{5}} \cdot y^{\tfrac{2}{5}}\right)^2=x^{\tfrac{6}{5}} \cdot  y^{\tfrac{4}{5}}=x^a \cdot y^b$.\\
		Suy ra  $a=\dfrac{6}{5}, b=\dfrac{4}{5} \Rightarrow a+b=2$.
	 }
\end{ex}

% nguồn: Ok_TraLoiNgan_Lop11_Lan1
\begin{ex}%[1D6H1-2]%[ma:0e2a2a5a]
	Biết  biểu thức $P=(5+2\sqrt{6})^{2024} \cdot(5-2\sqrt{6})^{2025}=a-2\sqrt{c}$ với $a$; $c$ là số tự nhiên. Tính giá trị $a^{c-2}$.
	\shortans{625} 
	 \loigiai{
		Ta có $(5+2\sqrt{6})(5-2\sqrt{6})=25-24=1$.\\ 
		Do đó, 
		$$ 
		P=(5+2 \sqrt{6})^{2024} \cdot(5-2 \sqrt{6})^{2025}=[(5+2  \sqrt{6})(5-2 \sqrt{6})]^{2024} \cdot(5-2 \sqrt{6})=5-2 \sqrt{6}.
		 $$
		Do đó  $a=5$; $c=6\Rightarrow a^{c-2}=625$.
	} 
\end{ex}

% nguồn: Ok_TraLoiNgan_Lop11_Lan1
\begin{ex}%[1D6H1-2]%[ma:02921889]
	Biết $10^{\alpha}=3$, $10^{\beta}=7$. Kết quả biểu  thức $A=\dfrac{100^{\alpha} \cdot 0{,}001^\beta}{10^{-\alpha} \cdot 10^{2\beta}}$ có dạng $\dfrac{27}{16\,a0b}$, với $a$, $b$ là các số nguyên. Tính giá trị $S=a+b^2$.
	\shortans{57} 
	\loigiai{ 
		 Ta có
		\begin{eqnarray*}
			A&=&\frac{\left(10^2\right)^\alpha \cdot\left(10^{-3}\right)^\beta}{10^{-\alpha} \cdot  10^{2\beta}}\\
			&=&\frac{10^{2  \alpha}}{10^{-\alpha}} \cdot \frac{10^{-3 \beta}}{10^{2\beta}}\\
			&=&10^{2 \alpha-(-\alpha)} \cdot  10^{-3 \beta-2 \beta}\\
			 &=&10^{3\alpha} \cdot 10^{-5\beta}\\
			&=&\left(10^\alpha\right)^3  \cdot\left(10^\beta\right)^{-5}\\
			&=&3^3 \cdot 7^{-5}=\frac{3^3}{7^5}=\frac{27}{16\,807}. 
		\end{eqnarray*}
		Suy ra $a=8$, $b=7$. Vậy  $S=57$.
	} 
\end{ex}

% nguồn: Ok_TraLoiNgan_Lop11_Lan1
\begin{ex}%[1D6H1-2]%[ma:54b6792d]
	Cho hai số  thực dương $a$, $b$. Rút gọn biểu thức $A=\dfrac{a^{\tfrac{1}{3}}\sqrt{b}+b^{\tfrac{1}{3}}\sqrt{a}}{\sqrt[6]{a}+\sqrt[6]{b}}$, ta được $A=a^m\cdot b^n$. Tính $S=3m+6n$.
	 \shortans{3}
	 \loigiai{
		 $A=\dfrac{a^{\tfrac{1}{3}}\sqrt{b}+b^{\tfrac{1}{3}}\sqrt{a}}{\sqrt[6]{a}+\sqrt[6]{b}}=\dfrac{a^{\tfrac{1}{3}}\cdot b^{\tfrac{1}{2}}+b^{\tfrac{1}{3}}\cdot a^{\tfrac{1}{2}}}{a^{\tfrac{1}{6}}+b^{\tfrac{1}{6}}}=\dfrac{a^{\tfrac{1}{3}}\cdot b^{\tfrac{1}{3}}\left(b^{\tfrac{1}{6}}+a^{\tfrac{1}{6}}\right)}{a^{\tfrac{1}{6}}+b^{\tfrac{1}{6}}}=a^{\tfrac{1}{3}}\cdot b^{\tfrac{1}{3}}$.\\
		Suy ra  $m=\dfrac{1}{3}$, $n=\dfrac{1}{3}\Rightarrow 3m+6n=3$.\\
	 }
\end{ex}

% nguồn: Ok_TraLoiNgan_Lop11_Lan2
\begin{ex}%[1D6H1-2]%[ma:601a5791]
	Rút gọn biểu thức $A=\dfrac{\sqrt[3]{a^7}\cdot a^{\tfrac{11}{3}}}{a^4\cdot \sqrt[7]{a^{-5}}}$ với $a>0$ ta được kết  quả $A=a^{\tfrac{m}{n}}$, trong đó $m,n\in \mathbb{N}^*$ và $\dfrac{m}{n}$ là phân số tối giản. Tính $m^2-n^2$.
	 \shortans{312}
	 \loigiai{
		Ta có $A=\dfrac{\sqrt[3]{a^7}\cdot  a^{\tfrac{11}{3}}}{a^4\cdot \sqrt[7]{a^{-5}}}=\dfrac{a^{\tfrac{7}{3}}\cdot a^{\tfrac{11}{3}}}{a^4\cdot a^{\tfrac{-5}{7}}}=\dfrac{a^6}{a^{\tfrac{23}{7}}}=a^{\tfrac{19}{7}}$.\\
		Mà $A=a^{\tfrac{m}{n}}$, $m,n\in \mathbb{N}^*$ và $\dfrac{m}{n}$ là phân số tối  giản nên $m=19$, $n=7$.\\
		Vậy  $m^2-n^2=312$.
	} 
\end{ex}

% nguồn: Ok_TraLoiNgan_Lop11_Lan2
\begin{ex}%[1D6H1-2]%[ma:200450c3]
	Cho $a$ là số thực thỏa mãn $2^a+2^{-a}=5$. Tính giá trị của biểu thức $4^a+4^{-a}$. 
	\shortans{23} 
	 \loigiai{
		\begin{align*}
			 2^a+2^{-a}=5\Rightarrow\left(2^a+2^{-a}\right)^2=25\Rightarrow 4^a+4^{-a}+2=25\Rightarrow 4^a+4^{-a}=23.
		\end{align*}
	 }
\end{ex}

% nguồn: Ok_TraLoiNgan_Lop11_Lan1
\begin{ex}%[1D6V1-2]%[ma:9ac98dcd]
	Rút gọn biểu thức $P=\dfrac{\sqrt{a}+\sqrt[4]{a b}}{\sqrt[4]a+\sqrt[4]b}-\dfrac{\sqrt{a}-\sqrt{b}}{\sqrt[4]a-\sqrt[4]b}$, $(a > 0, b > 0)$ ta thu được kết quả  có dạng $-m \sqrt{b}$ với $m; n$ là số tự nhiên. Tính giá trị $m+2n$.
	 \shortans{9}
	\loigiai{ 
	Ta  có
	 \allowdisplaybreaks
	\begin{eqnarray*}
		P&=&\dfrac{\sqrt{a}+\sqrt[4]{a}}{\sqrt[4]{a}+\sqrt[4]{b}}-\dfrac{\sqrt{a}-\sqrt{b}}{\sqrt[4]{a}-\sqrt[4]{b}}\\ 
		 &=&\dfrac{(\sqrt[4]{a})^2+\sqrt[4]{ab}}{\sqrt[4]{a}+\sqrt[4]{b}}-\dfrac{(\sqrt[4]{a})^2-(\sqrt[4]{b})^2}{\sqrt[4]{a}-\sqrt[4]{b}} \\
		 &=&\dfrac{\sqrt[4]{a}(\sqrt[4]{a}+\sqrt[4]{b})}{\sqrt[4]{a}+\sqrt[4]{b}}-\dfrac{(\sqrt[4]{a}-\sqrt[4]{b})(\sqrt[4]{a}+\sqrt[4]{b})}{\sqrt[4]{a}-\sqrt[4]{b}}\\
		&=&\sqrt[4]{a}-(\sqrt[4]{a}+\sqrt[4]{b})=-\sqrt[4]{b}. 
	\end{eqnarray*}
	Do đó, $m=1$ và $n=4$. Vậy  $m+2n=9$.
	 }
\end{ex}

% nguồn: Ok_TraLoiNgan_Lop11_Lan1
\begin{ex}%[1D6V1-2]%[ma:34cce992]
	Rút gọn biểu thức sau $Q=\dfrac{a^{\sqrt{3}+1}  \cdot a^{2-\sqrt{3}}}{\left(a^{\sqrt{2}-2}\right)^{\sqrt{2}+2}}$ với điều kiện $a > 0$ ta thu được kết quả có dạng $a^m$ với $m$ là số tự nhiên. Tính giá trị $m$.
	 \shortans{5}
	 \loigiai{
		Ta có 
		\begin{itemize}
			 \item $a^{\sqrt{3}+1} \cdot a^{2-\sqrt{3}}=a^{\sqrt{3}+1+2-\sqrt{3}}=a^3$.
			\item  $\left(a^{\sqrt{2}-2}\right)^{\sqrt{2}+2}=a^{(\sqrt{2}+2)(\sqrt{2}-2)}=a^{-2}$.
		\end{itemize}
		 Vậy $Q=\dfrac{a^3}{a^{-2}}=a^5$, do đó $m=5$.
	 }
\end{ex}

% nguồn: Ok_TraLoiNgan_Lop11_Lan1
\begin{ex}%[1D6V1-2]%[ma:f472aabe]
	Công ty FTK về mua bán xe ô tô đã qua sử dụng, sau khi khảo sát thị  trường $6$ tháng đã đưa ra công thức chung về giá trị còn lại của ô tô $4$ chỗ kể từ khi đưa vào sử đựng (các loại xe $4$ chỗ không sử dụng mục đích kinh doanh) được tính
	$$P(t)=A \cdot\left(\dfrac{3}{4}\right)^{\tfrac{t}{4}}.$$ 
	Trong đó $A$ là giá tiền ban đầu mua xe, $t$ là số năm kể từ khi đưa vào sử dụng. Giá trị còn lại của xe ô tô sau $30$ tháng đưa vào sử dụng có dạng $768\,601\,abc$, với $a$; $b$; $c$ là các số nguyên, tính giá trị $S=a+b+c$? Biết giá trị mua xe  ban đầu là $920$ triệu.
	\shortans{7} 
	 \loigiai{
		Ta  có $A=920$ triệu; $t=2{,}5$ năm.\\
		Vậy  giá trị còn lại của xe ô tô sau $30$ tháng đưa vào sử dụng là
		$$ 
		P(2{,}5)=920 \cdot  \left(\dfrac{3}{4}\right)^{\tfrac{2{,}5}{4}}=768\,601\,304.$$
		Suy ra $a=3$, $b=0$, $c=4$.  Vậy $S=7$.
	 }
\end{ex}

% nguồn: Ok_TraLoiNgan_Lop11_Lan1
\begin{ex}%[1D6V1-2]%[ma:21a1f754]
	Cho  biểu thức
	 $$
	P=\left[\dfrac{\sqrt[3]{a^2 b}-\sqrt[3]{a b^2}}{\sqrt[3]{a^2}-2 \sqrt[3]{a b}+\sqrt[3]{b^2}}-\dfrac{a+b}{\sqrt[3]{a^2}-\sqrt[3]{b^2}}\right](\sqrt[6]{a}+\sqrt[6]{b})^{-1}+\sqrt[6]{a}, \text{  với }a>0, b>0, a \neq b.
	$$ 
	Tính $\dfrac{P^6}{b}$. 
	\shortans{1} 
	 \loigiai{
		Ta có: 
		 $$
		\begin{aligned}
			& P=\left[\dfrac{\sqrt[3]{a^2 b}-\sqrt[3]{a b^2}}{\sqrt[3]{a^2}-2 \sqrt[3]{a b}+\sqrt[3]{b^2}}-\dfrac{a+b}{\sqrt[3]{a^2}-\sqrt[3]{b^2}}\right](\sqrt[6]{a}+\sqrt[6]{b})^{-1}+\sqrt[6]{a}\\ 
			& =\left[\dfrac{\sqrt[3]{a  b}(\sqrt[3]{a}-\sqrt[3]{b})}{(\sqrt[3]{a}-\sqrt[3]{b})^2}-\dfrac{(\sqrt[3]{a}+\sqrt[3]{b})\left(\sqrt[3]{a^2}-\sqrt[3]{a b}+\sqrt[3]{b^2}\right)}{(\sqrt[3]{a}-\sqrt[3]{b})(\sqrt[3]{a}+\sqrt[3]{b})}\right](\sqrt[6]{a}+\sqrt[6]{b})^{-1}+\sqrt[6]{a}\\
			& =\left[\dfrac{\sqrt[3]{a b}}{\sqrt[3]{a}-\sqrt[3]{b}}-\dfrac{\sqrt[3]{a^2}-\sqrt[3]{a b}+\sqrt[3]{b^2}}{\sqrt[3]{a}-\sqrt[3]{b}}\right]  \dfrac{1}{\sqrt[6]{a}+\sqrt[6]{b}}+\sqrt[6]{a}\\
			&  =-\dfrac{(\sqrt[3]{a}-\sqrt[3]{b})^2}{\sqrt[3]{a}-\sqrt[3]{b}}\dfrac{1}{(\sqrt[6]{a}+\sqrt[6]{b})}+\sqrt[6]{a}\\
			& =\dfrac{\sqrt[3]{b}-\sqrt[3]{a}}{\sqrt[6]{a}+\sqrt[6]{b}}+\sqrt[6]{a}-\sqrt[6]{a}+\sqrt[6]{a}=\sqrt[6]{b} 
		\end{aligned}
		$$ 
		Vậy  $\dfrac{P^6}{b}=1$.
	 }
\end{ex}

% nguồn: Ok_TraLoiNgan_Lop11_Lan2
\begin{ex}%[1D6V1-2]%[ma:949c6c9b]
	Cho hàm số $f(a)=\dfrac{a^{\tfrac{-1}{3}}\left(\sqrt[3]{a}-\sqrt[3]{a^4}\right)}{a^{\tfrac{1}{8}}\left(\sqrt[8]{a^3}-\sqrt[8]{a^{-1}}\right)}$, với $0<a\neq  1$. Tính giá trị $M=f(2025)$.
	 \shortans{-46}
	\loigiai{ 
		Ta  có\allowdisplaybreaks
		\begin{align*}
			 f(a)&=~\dfrac{a^{\tfrac{-1}{3}}\left(\sqrt[3]{a}-\sqrt[3]{a^4}\right)}{a^{\tfrac{1}{8}}\left(\sqrt[8]{a^3}-\sqrt[8]{a^{-1}}\right)}\\
			 &=~\dfrac{a^{\tfrac{-1}{3}}\left(a^{\tfrac{1}{3}}-a^{\tfrac{4}{3}}\right)}{a^{\tfrac{1}{8}}\left(a^{\tfrac{3}{8}}-a^{\tfrac{-1}{8}}\right)}\\
			&=~\dfrac{a^{\tfrac{-1}{3}}\cdot a^{\tfrac{1}{3}}(1-a)}{a^{\tfrac{1}{8}}\cdot  a^{\tfrac{-1}{8}}(a^{\tfrac{1}{2}}-1)}\\
			 &=~\dfrac{-(a^{\tfrac{1}{2}}-1)(a^{\tfrac{1}{2}}+1)}{(a^{\tfrac{1}{2}}-1)}\\
			 &=~-(a^{\tfrac{1}{2}}+1)=-a^{\tfrac{1}{2}}-1=-\sqrt{a}-1.
		\end{align*}
		Do đó, $M=f(2025)=-\sqrt{2025}-1=-46$. 
	} 
\end{ex}

% nguồn: Ok_TraLoiNgan_Lop11_Lan2
\begin{ex}%[1D6V1-2]%[ma:77bd350f]
	 Cho $9^x+9^{-x}=14$. Khi đó, giá trị của biểu thức $\dfrac{6+3\left(3^x+3^{-x}\right)}{3^{x+1}+3^{1-x}-2}=\dfrac{a}{b}$ ($a,b\in\mathbb{N}$; $\dfrac{a}{b}$ là phân số tối giản). Tính $P=a+b$.
	 \shortans{14}
	 \loigiai{
		Ta có $9^x+9^{-x}=14$ $\Leftrightarrow$  $\left(3^x+3^{-x}\right)^2=16\Rightarrow 3^x+3^{-x}=4$. Do đó,
		\begin{align*}
			\dfrac{6+3\left(3^x+3^{-x}\right)}{3^{x+1}+3^{1-x}-2}=\dfrac{6+3\left(3^x+3^{-x}\right)}{3\cdot \left(3^x+3^{-x}\right)-2}=\dfrac{6+3\cdot  4}{3\cdot 4-2}=\dfrac{9}{5}.
		\end{align*}
		Suy ra $a=9$  và $b=5$.\\
		Vậy $P=a+b=14$. 
	 }
\end{ex}

% nguồn: Ok_TraLoiNgan_Lop11_Lan2
\begin{ex}%[1D6V1-2]%[ma:e4828f45]
	Tìm tất cả các số thực $m$  sao cho $\dfrac{4^a}{4^a+m}+\dfrac{4^b}{4^b+m}=1$ với mọi $a+b=1$.
	\shortans{2} 
	\loigiai{ 
		Ta có $a+b=1\Leftrightarrow  b=1-a$.\\
		 Thay vào $\dfrac{4^a}{4^a+m}+\dfrac{4^b}{4^b+m}=1$ ta được
		\begin{align*}
			\dfrac{4^a}{4^a+m}+\dfrac{4^{1-a}}{4^{1-a}+m}=1\Leftrightarrow\dfrac{4+m\cdot 4^a+4+m\cdot 4^{1-a}}{4+m\cdot 4^a+m\cdot 4^{1-a}+m^2}=1\Leftrightarrow  m^2=4\Leftrightarrow m=\pm 2.
		\end{align*}
		Thử lại, với $m=2$, thỏa mãn; với $m=-2$ thì phương trình không xác  định với $a=b=\dfrac{1}{2}$.\\
		Vậy $m=2$ là giá trị  cần tìm.
	} 
\end{ex}

% nguồn: Ok_TraLoiNgan_Lop11_Lan2
\begin{ex}%[1D6V1-2]%[ma:1ec25d3d]
	Cho $4^x+4^{-x}=7$. Khi đó, biểu thức $P=\dfrac{5-2^x-2^{-x}}{8+4\cdot 2^x+4\cdot  2^{-x}}=\dfrac{a}{b}$ với $\dfrac{a}{b}$ là phân số tối giản và $a\in\mathbb{Z}$, $b\in\mathbb{N}^*$. Tính tích $a\cdot b$.
	 \shortans{10}
	 \loigiai{
		Ta  có
		\begin{align*}
			4^x+4^{-x}=7\Leftrightarrow\left(2^x\right)^2+\left(2^{-x}\right)^2=7\Leftrightarrow\left(2^x+2^{-x}\right)^2-2\cdot  2^x\cdot 2^{-x}=7\Leftrightarrow\left(2^x+2^{-x}\right)^2=9.
		\end{align*}
		Suy ra  $2^x+2^{-x}=3$.\\
		Do  đó, $P=\dfrac{5-\left(2^x+2^{-x}\right)}{8+4\cdot \left(2^x+2^{-x}\right)}=\dfrac{5-3}{8+4\cdot 3}=\dfrac{2}{20}=\dfrac{1}{10}$.\\
		Suy ra $a=1$,  $b=10$.\\
		Vậy $a\cdot  b=10$.
	} 
\end{ex}

% nguồn: Ok_TraLoiNgan_Lop11_Lan2
\begin{ex}%[1D6V1-2]%[ma:cc3d778d]
	Cho số  thực $x$ thỏa mãn $9^x+9^{-x}=14$. Tính giá trị biểu thức $P=\dfrac{6+3\left(3^x+3^{-x}\right)}{2-3^{x+1}-3^{1-x}}$.
	 \shortans{-1,8}
	 \loigiai{
		Theo  bài ra ta có
		\begin{align*}
			 9^x+9^{-x}=14\Leftrightarrow\left(3^x\right)^2+\left(3^{-x}\right)^2=14\Leftrightarrow(3^x)^2+2\cdot 3^x\cdot 3^{-x}+\left(3^{-x}\right)^2=16\Leftrightarrow\left(3^x+3^{-x}\right)^2=16.
		\end{align*}
		Suy ra $3^x+3^{-x}=4$.\\ 
		Do  đó
		\begin{align*}
			P=\dfrac{6+3\left(3^x+3^{-x}\right)}{2-3\cdot 3^x-3\cdot 3^{-x}}=\dfrac{6+3\left(3^x+3^{-x}\right)}{2-3\left(3^x+3^{-x}\right)}=\dfrac{6+3\cdot 4}{2-3\cdot 4}=\dfrac{18}{-10}=-\dfrac{9}{5}. 
		\end{align*}
		Vậy $P=-\dfrac{9}{5}=-1{,}8$. 
	} 
\end{ex}

% nguồn: Ok_TraLoiNgan_Lop11_Lan2
\begin{ex}%[1D6V1-2]%[ma:adb4b7c4]
	Cho $a$, $b$ là các số thực thỏa mãn $3\cdot a^a+2^b=7\sqrt{2}$ và $5\cdot  2^a-2^b=9\sqrt{2}$. Tính tổng $S=a+b$.
	\shortans{2} 
	\loigiai{ 
		 Ta có
		\begin{align*}
			 \heva{
			& 3\cdot 2^a+2^b=7\sqrt{2}\\ 
			& 5\cdot  2^a-2^b=9\sqrt{2}
			 }\Leftrightarrow\heva{
			&  3\cdot 2^a+2^b=7\sqrt{2}\\
			 & 8\cdot 2^a=16\sqrt{2}
			 }\Leftrightarrow\heva{
			& 2^b=7\sqrt{2}-3\cdot  2^a\\
			& 2^a=2\sqrt{2} 
			}\Leftrightarrow\heva{ 
			& 2^b=\sqrt{2}\\ 
			& 2^a=2\sqrt{2}. 
			} 
		\end{align*}
		Từ đó suy ra  $2^{a+b}=2^a\cdot 2^b=2\sqrt{2}\cdot \sqrt{2}=4=2^2$ $\Rightarrow a+b=2$.\\
		Vậy  $S=a+b=2$.
	} 
\end{ex}

% nguồn: Ok_TraLoiNgan_Lop11_Lan2
\begin{ex}%[1D6V1-2]%[ma:893fc092]
	Cho $9^x+9^{-x}=23$.  Tính giá trị của biểu thức $P=\dfrac{5+3^x+3^{-x}}{1-3^x-3^{-x}}$.
	 \shortans{-2,5}
	\loigiai{ 
		 Ta có $9^x+9^{-x}=23\Leftrightarrow\left(3^x+3^{-x}\right)^2=25\Leftrightarrow\hoac{
		& 3^x+3^{-x}=5\\ 
		& 3^x+3^{-x}=-5\quad~(\text{vô  nghiệm}).
		 }$\\
		 Từ đó, ta được $P=\dfrac{5+\left(3^x+3^{-x}\right)}{1-\left(3^x+3^{-x}\right)}=\dfrac{5+5}{1-5}=-\dfrac{5}{2}=-2{,}5$.
	} 
\end{ex}

% nguồn: Ok_TraLoiNgan_Lop11_Lan2
\begin{ex}%[1D6V1-2]%[ma:089e6a4a]
	Cho  $\sqrt{\alpha-2+\sqrt{2\alpha-5}}+\sqrt{\alpha+2+3\sqrt{2\alpha-5}}=7\sqrt{2}$, $\alpha\in\left[\dfrac{5}{2};+\infty\right)$. Tìm $\alpha$.
	\shortans{15} 
	 \loigiai{
		Đặt  $t=\sqrt{2\alpha-5}$.\\
		$\forall\alpha\in\left[\dfrac{5}{2};+\infty\right)$, ta có $t\ge 0$ và $\alpha-2=\dfrac{t^2+1}{2}$ và $\alpha+2=\dfrac{t^2+9}{2}$.\\ 
		Khi đó, đẳng thức $\sqrt{\alpha-2+\sqrt{2\alpha-5}}+\sqrt{\alpha+2+3\sqrt{2\alpha-5}}=7\sqrt{2}$ trở  thành
		\begin{align*}
			\sqrt{(t+1)^2}+\sqrt{(t+3)^2}=14\Leftrightarrow\left|t+1\right|+\left|t+3\right|=14\Leftrightarrow  t+1+t+3=14\Leftrightarrow t=5.
		\end{align*}
		Do đó, $\sqrt{2\alpha-5}=5\Rightarrow 2\alpha-5=25\Leftrightarrow\alpha=15$. 
	 }
\end{ex}

% nguồn: Ok_TraLoiNgan_Lop11_Lan1
\begin{ex}%[1D6C1-2]%[ma:64c4e865]
	Cho $f(x)=\mathrm{e}^{\sqrt{\sqrt{1+\tfrac{1}{x^2}+\tfrac{1}{(x+1)^2}}}}$. Biết rằng $f(1) \cdot f(2) \cdot f(3) \cdots f(2\,025)=\mathrm{e}^{\tfrac{m}{n}}$ với $m$, $n$ là các số tự nhiên và $\dfrac{m}{n}$ là phân số tối giản.  Tính $m-n^2$.
	\shortans{-1} 
	 \loigiai{
		Đặt  $g(x)=\sqrt{1+\dfrac{1}{x^2}+\dfrac{1}{(1+x)^2}}$.\\
		 Với $x>0$ ta có
		$$\begin{aligned}
		g(x) & =\sqrt{1+\frac{1}{x^2}+\frac{1}{(1+x)^2}}=\frac{\sqrt{x^2+(x+1)^2+x^2 \cdot(x+1)^2}}{x(x+1)}=\frac{\sqrt{\left(x^2+x+1\right)^2}}{x(x+1)} \\ 
		& =\frac{x^2+x+1}{x(x+1)}=1+\frac{1}{x(x+1)}=1+\frac{1}{x}-\frac{1}{x+1} 
	\end{aligned}
	 $$
	Suy  ra
	\begin{eqnarray*}
		 &&g(1)+g(2)+g(3)+\cdots+g(2\,025)\\
		&=&\left(1+\frac{1}{1}-\frac{1}{2}\right)+\left(1+\frac{1}{2}-\frac{1}{3}\right)+\left(1+\frac{1}{3}-\frac{1}{4}\right)+\cdots+\left(1+\frac{1}{2\,025}-\frac{1}{2\,026}\right)\\ 
		 &=&2\,026-\frac{1}{2\,026}.
	\end{eqnarray*}
	Khi đó  $f(1) \cdot f(2) \cdot f(3) \cdots f(2\,025)=\mathrm{e}^{g(1)+g(2)+g(3)+\cdots+g(2\,025)}=\mathrm{e}^{2\,026-\tfrac{1}{2\,026}}=\mathrm{e}^{\tfrac{2\,026^2-1}{2\,026}}=\mathrm{e}^{\tfrac{m}{n}}$.\\
	Do đó  $m=2\,026^2-1$, $n=2\,026$.\\
	Vậy  $m-n^2=2\,026^2-1-2\,026^2=-1$.
} 
\end{ex}

% nguồn: Ok_TraLoiNgan_Lop11_Lan2
\begin{ex}%[1D6C1-2]%[ma:0e80ae23]
	Cho hàm số $f(x)=\dfrac{2^x}{2^x+2}$. Tính tổng $f(0)+f\left(\dfrac{1}{10}\right)+\cdots+f\left(\dfrac{18}{10}\right)+f\left(\dfrac{19}{10}\right)$  (kết quả làm tròn đến hàng phần trăm).
	\shortans{9,83} 
	\loigiai{ 
		Với $a+b=2$, ta có $f(a)+f(b)=\dfrac{2^a}{2^a+2}+\dfrac{2^b}{2^b+2}=\dfrac{2\cdot 2^{a+b}+2\cdot  2^a+2\cdot 2^b}{2^{a+b}+2\cdot 2^a+2\cdot 2^b+4}=1$.\\
		Do  đó,\allowdisplaybreaks
		\begin{align*}
			 &f(0)+f\left(\dfrac{1}{10}\right)+\cdots+f\left(\dfrac{18}{10}\right)+f\left(\dfrac{19}{10}\right)\\
			 =&f(0)+\left[f\left(\dfrac{1}{10}\right)+f\left(\dfrac{19}{10}\right)\right]+\cdots+\left[f\left(\dfrac{9}{10}\right)+f\left(\dfrac{11}{10}\right)\right]+f(1)\\
			=&f(0)+9\cdot 1+f(1)=\dfrac{1}{3}+9+\dfrac{1}{2}=\dfrac{59}{6}\approx9{,}83 
		\end{align*}
	} 
\end{ex}

% nguồn: Ok_TraLoiNgan_Lop11_Lan2
\begin{ex}%[1D6C1-2]%[ma:2d80c6ea]
	Biết rằng $P=2023!\cdot \left(1+\dfrac{1}{1}\right)^1\left(1+\dfrac{1}{2}\right)^2\cdots \left(1+\dfrac{1}{2022}\right)^{2024}\left(1+\dfrac{1}{2023}\right)^{2025}=2024^m$.  Tìm $m$.
	 \shortans{2025}
	 \loigiai{
		Ta có\allowdisplaybreaks 
		\begin{align*}
			P&=2023!\cdot \left(1+\dfrac{1}{1}\right)^1\left(1+\dfrac{1}{2}\right)^2\cdots\left(1+\dfrac{1}{2022}\right)^{2024}\left(1+\dfrac{1}{2023}\right)^{2025}\\ 
			&=2023!\cdot \dfrac{2^1\cdot 3^2\cdots 2023^{2024}\cdot 2024^{2025}}{1^1\cdot 2^2\cdots2022^{2024}\cdot 2023^{2025}}\\ 
			&=2023!\cdot \dfrac{2024^{2025}}{2\cdot 3\cdots  2022\cdot 2023}\\
			&=2023!\cdot  \dfrac{2024^{2025}}{2023!}=2024^{2025}.
		\end{align*}
		Vậy $m=2025$. 
	} 
\end{ex}

% nguồn: Ok_TraLoiNgan_Lop11_Lan2
\begin{ex}%[1D6C1-2]%[ma:69d2d1df]
	Cho $f(x)=\dfrac{2^x}{2^x+2}$. Tính tổng 
	\begin{align*}
		T=f(0)+f\left(\dfrac{1}{2024}\right)+f\left(\dfrac{2}{2024}\right)+\cdots+f\left(\dfrac{4046}{2024}\right)+f\left(\dfrac{4047}{2024}\right) 
	\end{align*}
	(Kết quả  làm tròn đến hàng đơn vị).
	\shortans{2024} 
	\loigiai{ 
		Với $a,b\in\mathbb{R}$ thỏa mãn  $a+b=1$ ta có
		\begin{align*}
			 f(a)+f(b)=\dfrac{2^a}{2^a+2}+\dfrac{2^b}{2^b+2}=\dfrac{2^{a+b}+2\cdot 2^a+2\cdot 2^b+2^{a+b}}{2^{a+b}+2(2^a+2^b)+2}=\dfrac{2+22(2^a+2^b)}{2+2(2^a+2^b)}=1.
		\end{align*}
		 Ta viết tổng cần tính theo $2$ cách
		\begin{align*}
			 T&=~f(0)+f\left(\dfrac{1}{2024}\right)+f\left(\dfrac{2}{2024}\right)+\cdots+f\left(\dfrac{4047}{2024}\right)\tag{1}\\
			T&=~f\left(\dfrac{4047}{2024}\right)+f\left(\dfrac{4046}{2024}\right)+\cdots+f(0).\tag{2} 
		\end{align*}
		Lấy cộng vế theo  vế $(1)$ và $(2)$ ta được
		\begin{align*}
			 2T&=~2f(0)+\left[f\left(\dfrac{1}{2024}\right)+f\left(\dfrac{4047}{2024}\right)\right]+\cdots+\left[f\left(\dfrac{4047}{2024}\right)+f\left(\dfrac{1}{2024}\right)\right]\\
			&=~\dfrac{2}{3}+4047\cdot 1=\dfrac{12143}{3}. 
		\end{align*}
		Do đó, $T=\dfrac{12143}{6}\approx  2024$.
	 }
\end{ex}

% nguồn: Ok_TraLoiNgan_Lop11_Lan2
\begin{ex}%[1D6C1-2]%[ma:d614baf6]
	Cho $f(x)=5^{\sqrt{1+\tfrac{1}{x^2}+\tfrac{1}{(x+1)^2}}}$. Biết $f(1)\cdot f(2)\cdots f(2023)=5^{\tfrac{m}{n}}$ với $m$, $n$ là các  số nguyên dương và phân số $\dfrac{m}{n}$ tối giản. Tính giá trị của biểu thức $m-n^2$.
	 \shortans{-1}
	 \loigiai{
		Ta có\allowdisplaybreaks 
		\begin{align*}
			 1+\dfrac{1}{x^2}+\dfrac{1}{(x+1)^2}&=~\dfrac{x^2(x+1)^2+(x+1)^2+x^2}{x^2(x+1)^2}\\
			 &=~\dfrac{(x(x+1))^2+2x^2+2x+1}{x^2(x+1)^2}\\
			&=~\dfrac{(x^2+x)^2+2(x^2+x)+1}{x^2(x+1)^2}\\ 
			 &=~\dfrac{(x^2+x+1)^2}{(x(x+1))^2}=\left(1+\dfrac{1}{x}-\dfrac{1}{x+1}\right)^2.
		\end{align*}
		 Do đó, $f(x)=5^{\sqrt{1+\tfrac{1}{x^2}+\tfrac{1}{(x+1)^2}}}=5^{1+\tfrac{1}{x}-\tfrac{1}{x+1}}$.\\
		Suy ra $f(1)\cdot f(2)\cdots f(2023)=5^{\sum_{x=1}^{2023}\left(1+\tfrac{1}{x}-\tfrac{1}{x+1}\right)}$.\\ 
		Bằng  phương pháp quy nạp ta chứng minh được
		\begin{align*}
			\sum_{x=1}^k\left(1+\dfrac{1}{x}-\dfrac{1}{x+1}\right)=k+1-\dfrac{1}{k+1}=\dfrac{(k+1)^2-1}{k+1}. 
		\end{align*}
		 Thay $k=2023$, ta được $\dfrac{m}{n}=\dfrac{2024^2-1}{2024}$.\\
		Suy ra $m=2024^2-1$, $n=2024$.\\ 
		 Vậy $m-n^2=(2024^2-1)-2024^2=-1$.
	} 
\end{ex}


%%%%%%%%%%%%%%%% TỰ LUẬN %%%%%%%%%%%%%%%%

% nguồn: Ok_TuanLuan_Toan11_Lan1
\begin{bt}%[1D6H1-2]%[ma:18ca6c76]
	 Rút gọn biểu thức $A=\dfrac{a^{\frac{7}{3}}-a^{\frac{1}{3}}}{\sqrt[3]{a^4}+\sqrt[3]{a}}$ $(a>0)$.
	\loigiai{ 
	 Ta có $A=\dfrac{a^{\frac{7}{3}}-a^{\frac{1}{3}}}{a^{\frac{4}{3}}+a^{\frac{1}{3}}}=\dfrac{a^{\frac{1}{3}}\left(a^2-1\right)}{a^{\frac{1}{3}}(a+1)}=\dfrac{a^2-1}{a+1}=a-1$.
	} 
\end{bt}

% nguồn: Ok_TuanLuan_Toan11_Lan1
\begin{bt}%[1D6H1-2]%[1D6V2-2]%[ma:fdb96e36]
	\begin{enumerate}
		\item  Rút gọn biểu thức $A=\dfrac{\sqrt[5]{a^{15} b^5}}{\sqrt[3]{a^6 b^3}}$ (với $a, b$ là những số thực dương).
		\item Đặt $a=\log_3 2; b=\log_3  11; c=\log_3 23$, hãy biểu diễn $\log_3 2024$ theo ba số $a, b, c$.
	\end{enumerate}
	\loigiai{ 
		\begin{enumerate}
			\item Ta có  $A=\dfrac{\sqrt[5]{a^{15} b^5}}{\sqrt[3]{a^6 b^3}}=\dfrac{a^{5}b}{a^2b}=a^3$.
			\item Ta có $\log_3  2024=\log_3({2^3\cdot 11\cdot 23})=3\log_3{2}+\log_3{11}+\log_3{23}=3a+b+c$.
		\end{enumerate}
	} 
\end{bt}

% nguồn: Ok_TuanLuan_Toan11_Lan1
\begin{bt}%[1D6H1-2]%[ma:e23e0727]
	Viết biểu thức sau dưới dạng một lũy thừa  ($a>0$)
	\begin{listEX}[2]
		\item $a^{\tfrac{7}{5}}:  \sqrt[5]{a}$.
		 \item $\sqrt{a^7 \cdot \sqrt[3]{a}}$.
	\end{listEX}
	 \loigiai{
		\begin{enumerate}
			\item Ta có $a^{\tfrac{7}{5}}:  \sqrt[5]{a}=a^{\tfrac{7}{5}}: a^{\tfrac{1}{5}}=a^{\tfrac{6}{5}}$.
			\item $\sqrt{a^7 \cdot \sqrt[3]{a}}  =a^{\tfrac{7}{2}}\cdot \sqrt[6]{a}=a^{\tfrac{7}{2}}\cdot a^{\tfrac{1}{6}} =a^{\tfrac{11}{3}} $.
		\end{enumerate}
	} 
\end{bt}

% nguồn: Ok_TuanLuan_Toan11_Lan1
\begin{bt}%[1D6H1-2]%[ma:cc04870d]
	Cho $a>0$, $b>0$, rút gọn biểu thức sau 
	\begin{listEX}[2]
		\item  $a^\frac{1}{3}a^\frac{1}{2}a^\frac{7}{6}$.
		 \item $\sqrt[3]{a^2}a^\frac{1}{4}:a^\frac{1}{6}$.
	\end{listEX}
	\loigiai{ 
		\begin{enumerate}
			\item $a^\frac{1}{3}a^\frac{1}{2}a^\frac{7}{6} =  a^{\frac{1}{3}+\frac{1}{2}+\frac{7}{6}} = a^\frac{2+3+7}{6} = a^2$.
			\item $\sqrt[3]{a^2}a^\frac{1}{4}:a^\frac{1}{6} = a^\frac{2}{3} a^\frac{1}{4}:a^\frac{1}{6} = a^{\frac{2}{3}  +\frac{1}{4}-\frac{1}{6}} = a^\frac{8+3-2}{12} = a^\frac{3}{4}$.
		\end{enumerate}
	} 
\end{bt}

% nguồn: Ok_TuanLuan_Toan11_Lan1
\begin{bt}%[1D6H1-2]%[ma:819c53a7]
	 Đưa các biểu thức sau về dạng lũy thừa:
	\begin{enumerate}
		\item $\sqrt{a\sqrt[3]a}$ ($a>0$). 
		\item  $\dfrac{\sqrt{2}\cdot\sqrt[3]4}{16^{0{,}75}}$.
		 \item $\sqrt{\sqrt[5]{a^{15}\cdot b^{15}}}~(a,b>0)$.
	\end{enumerate}
	\loigiai{ 
		\begin{enumerate}
			\item $\sqrt{a\sqrt[3]a}=\left(a\cdot a^{\tfrac{1}{3}}\right) ^{\tfrac{1}{2}}=\left( a^{\tfrac{4}{3}}\right) ^{\tfrac{1}{2}}=a^{\tfrac{2}{3}}$. 
			 \item $\dfrac{\sqrt{2}\cdot\sqrt[3]4}{16^{0{,}75}}=\dfrac{2^{\tfrac{1}{2}}\cdot2^{\tfrac{2}{3}}}{2^3\cdot\left( 2^4\right) ^{\tfrac{3}{4}}}=\dfrac{2^{\tfrac{7}{6}}}{2^6}=2^{\tfrac{7}{6}-6}=2^{-\tfrac{29}{6}}$.
			\item $\sqrt{\sqrt[5]{a^{15}\cdot b^{15}}}=\left(  a^{15}\cdot b^{15}\right) ^{\tfrac{1}{5}\cdot\tfrac{1}{2}}=\left( a^{15}\cdot b^{15}\right) ^{\tfrac{1}{10}}=a^{\tfrac{3}{2}}\cdot b^{\tfrac{3}{2}}$.
		\end{enumerate}
	} 
\end{bt}

% nguồn: Ok_TuanLuan_Toan11_Lan1
\begin{bt}%[1D6H1-2]%[ma:915355c7]
	 Thực hiện các yêu cầu sau:
	\begin{enumerate}
		\item Cho $a$ là một số thực dương. Viết biểu thức $P=a^{\frac{3}{5}} \cdot \sqrt[3]{a^2}$ dưới dạng lũy  thừa với số mũ hữu tỉ.
		\item  Viết biểu thức $P=\sqrt[3]{x \cdot \sqrt[4]{x}}$ ($x>0$) dưới dạng luỹ thừa với số mũ hữu tỷ.
	\end{enumerate}
	 \loigiai{
		\begin{enumerate}
			\item Ta có $P=a^{\frac{3}{5}} \cdot  \sqrt[3]{a^2}=a^{\frac{3}{5}} \cdot a^{\frac{2}{3}}=a^{\frac{3}{5}+\frac{2}{3}}=a^{\frac{19}{15}}$.
			\item $P=\sqrt[3]{x  \cdot \sqrt[4]{x}}=\sqrt[3]{x \cdot x^{\frac{1}{4}}}=\left(x^{\frac{5}{4}}\right)^{\frac{1}{3}}=x^{\frac{5}{12}}$.
		\end{enumerate}
	 }
\end{bt}

% nguồn: Ok_TuanLuan_Toan11_Lan1
\begin{bt}%[1D6H1-2]%[ma:9f186d00]
	Cho  $x>0$, rút gọn biểu thức $A=\dfrac{x^{\tfrac{4}{5}} \cdot \sqrt[6]{x^5}}{\sqrt[3]{x \cdot \sqrt{x}}}$.
	\loigiai{Ta có 
		$A=\dfrac{x^{\tfrac{4}{5}}  \cdot x^{\tfrac{5}{6}}}{\left( x\cdot x^{\tfrac{1}{2}}\right)^{\tfrac{1}{3}}}
		=\dfrac{x^{\tfrac{4}{5}+\tfrac{5}{6}}}{x^{\tfrac{1}{3}}\cdot x^{\tfrac{1}{6}}} 
		=x^{\tfrac{4}{5}+\tfrac{5}{6}-\tfrac{1}{3}-\tfrac{1}{6}}=x^{\tfrac{17}{15}}$.\\ 
	 }
\end{bt}

% nguồn: Ok_TuanLuan_Toan11_Lan1
\begin{bt}%[1D6H1-2]%[ma:7520624f]
	Rút gọn các  biểu thức:
	\begin{enumerate}
		 \item $P=a^{\sqrt{2}}\cdot\left( \dfrac{1}{a}\right)^{\sqrt{2}-1}$ với $a>0$.
		\item $Q=a^{\tfrac{7}{3}}:\sqrt[3]a$ với  $a>0$.
		\item $K=\dfrac{x^{\tfrac{1}{3}}\cdot \sqrt[6]{x}}{\sqrt[4]{x}}$ với $x>0$. 
		\item $M=x^{\tfrac{1}{3}}\cdot \sqrt[6]{x}$ với $x>0$. 
	\end{enumerate}
	\loigiai{ 
		\begin{enumerate}
			\item $P=a^{\sqrt{2}}\cdot\left(  \dfrac{1}{a}\right)^{\sqrt{2}-1}=a^{\sqrt{2}}\cdot a^{-\sqrt{2}+1}=a^{\sqrt{2}-\sqrt{2}+1}=a$.
			\item  $Q=a^{\tfrac{7}{3}}:\sqrt[3]a=a^{\tfrac{7}{3}}:a^{\tfrac{1}{3}}=a^{\tfrac{7}{3}-\tfrac{1}{3}}=a^2$.
			\item $K=\dfrac{x^{\tfrac{1}{3}}\cdot \sqrt[6]{x}}{\sqrt[4]{x}}=\dfrac{x^{\tfrac{1}{3}}\cdot x^{\tfrac{1}{6}}}{x^{\tfrac{1}{4}}}=x^{\tfrac{1}{3}+\tfrac{1}{6}-\tfrac{1}{4}}=x^{\tfrac{1}{4}}=\sqrt[4]{x}$. 
			\item $M=x^{\tfrac{1}{3}}\cdot  \sqrt[6]{x}=x^{\tfrac{1}{3}}\cdot x^{\tfrac{1}{6}}=x^{\tfrac{1}{2}}=\sqrt{x}$.
		\end{enumerate}
	} 
\end{bt}

% nguồn: Ok_TuanLuan_Toan11_Lan2
\begin{ex}%[1D6H1-2]%[ma:fae301d6]
	Rút gọn biểu thức $A=\dfrac{a^{\frac{7}{3}}-a^{\frac{1}{3}}}{\sqrt[3]{a^4}+\sqrt[3]{a}}$  $(a>0)$.
	 \loigiai{
	Ta có  $A=\dfrac{a^{\frac{7}{3}}-a^{\frac{1}{3}}}{a^{\frac{4}{3}}+a^{\frac{1}{3}}}=\dfrac{a^{\frac{1}{3}}\left(a^2-1\right)}{a^{\frac{1}{3}}(a+1)}=\dfrac{a^2-1}{a+1}=a-1$.
	} 
\end{ex}

% nguồn: Ok_TuanLuan_Toan11_Lan2
\begin{ex}%[1D6H1-2]%[1D6V2-2]%[ma:ee9f30a5]
	\begin{enumerate}
		\item Rút gọn biểu thức $A=\dfrac{\sqrt[5]{a^{15} b^5}}{\sqrt[3]{a^6 b^3}}$ (với $a, b$ là những  số thực dương).
		\item Đặt  $a=\log_3 2; b=\log_3 11; c=\log_3 23$, hãy biểu diễn $\log_3 2024$ theo ba số $a, b, c$.
	\end{enumerate}
	\loigiai{ 
		\begin{enumerate}
			\item Ta có  $A=\dfrac{\sqrt[5]{a^{15} b^5}}{\sqrt[3]{a^6 b^3}}=\dfrac{a^{5}b}{a^2b}=a^3$.
			\item Ta  có $\log_3 2024=\log_3({2^3\cdot 11\cdot 23})=3\log_3{2}+\log_3{11}+\log_3{23}=3a+b+c$.
		\end{enumerate}
	 }
\end{ex}

% nguồn: Ok_TuanLuan_Toan11_Lan2
\begin{ex}%[1D6H1-2]%[ma:49a1c7f2]
	Viết biểu thức sau  dưới dạng một lũy thừa ($a>0$)
	\begin{listEX}[2]
		 \item $a^{\tfrac{7}{5}}: \sqrt[5]{a}$.
		\item $\sqrt{a^7 \cdot  \sqrt[3]{a}}$.
	\end{listEX}
	\loigiai{ 
		\begin{enumerate}
			\item Ta  có $a^{\tfrac{7}{5}}: \sqrt[5]{a}=a^{\tfrac{7}{5}}: a^{\tfrac{1}{5}}=a^{\tfrac{6}{5}}$.
			\item $\sqrt{a^7 \cdot \sqrt[3]{a}}  =a^{\tfrac{7}{2}}\cdot \sqrt[6]{a}=a^{\tfrac{7}{2}}\cdot a^{\tfrac{1}{6}} =a^{\tfrac{11}{3}} $.
		\end{enumerate}
	 }
\end{ex}

% nguồn: Ok_TuanLuan_Toan11_Lan2
\begin{ex}%[1D6H1-2]%[ma:9dc914db]
	Cho $a>0$, $b>0$, rút gọn biểu thức sau 
	\begin{listEX}[2]
		 \item $a^\frac{1}{3}a^\frac{1}{2}a^\frac{7}{6}$.
		\item $\sqrt[3]{a^2}a^\frac{1}{4}:a^\frac{1}{6}$. 
	\end{listEX}
	 \loigiai{
		\begin{enumerate}
			\item $a^\frac{1}{3}a^\frac{1}{2}a^\frac{7}{6} = a^{\frac{1}{3}+\frac{1}{2}+\frac{7}{6}}  = a^\frac{2+3+7}{6} = a^2$.
			\item $\sqrt[3]{a^2}a^\frac{1}{4}:a^\frac{1}{6} = a^\frac{2}{3} a^\frac{1}{4}:a^\frac{1}{6} = a^{\frac{2}{3} +\frac{1}{4}-\frac{1}{6}} =  a^\frac{8+3-2}{12} = a^\frac{3}{4}$.
		\end{enumerate}
	 }
\end{ex}

% nguồn: Ok_TuanLuan_Toan11_Lan2
\begin{ex}%[1D6H1-2]%[ma:96d79f1a]
	Đưa  các biểu thức sau về dạng lũy thừa:
	\begin{enumerate}
		\item $\sqrt{a\sqrt[3]a}$ ($a>0$). 
		\item  $\dfrac{\sqrt{2}\cdot\sqrt[3]4}{16^{0{,}75}}$.
		\item $\sqrt{\sqrt[5]{a^{15}\cdot b^{15}}}~(a,b>0)$. 
	\end{enumerate}
	\loigiai{ 
		\begin{enumerate}
			\item $\sqrt{a\sqrt[3]a}=\left(a\cdot a^{\tfrac{1}{3}}\right) ^{\tfrac{1}{2}}=\left( a^{\tfrac{4}{3}}\right)  ^{\tfrac{1}{2}}=a^{\tfrac{2}{3}}$.
			\item $\dfrac{\sqrt{2}\cdot\sqrt[3]4}{16^{0{,}75}}=\dfrac{2^{\tfrac{1}{2}}\cdot2^{\tfrac{2}{3}}}{2^3\cdot\left( 2^4\right) ^{\tfrac{3}{4}}}=\dfrac{2^{\tfrac{7}{6}}}{2^6}=2^{\tfrac{7}{6}-6}=2^{-\tfrac{29}{6}}$. 
			\item $\sqrt{\sqrt[5]{a^{15}\cdot  b^{15}}}=\left( a^{15}\cdot b^{15}\right) ^{\tfrac{1}{5}\cdot\tfrac{1}{2}}=\left( a^{15}\cdot b^{15}\right) ^{\tfrac{1}{10}}=a^{\tfrac{3}{2}}\cdot b^{\tfrac{3}{2}}$.
		\end{enumerate}
	} 
\end{ex}

% nguồn: Ok_TuanLuan_Toan11_Lan2
\begin{ex}%[1D6H1-2]%[ma:5080e80f]
	Thực hiện các yêu cầu  sau:
	\begin{enumerate}
		\item Cho $a$ là một số thực dương. Viết biểu thức $P=a^{\frac{3}{5}} \cdot \sqrt[3]{a^2}$ dưới dạng lũy thừa với  số mũ hữu tỉ.
		\item Viết biểu thức $P=\sqrt[3]{x \cdot  \sqrt[4]{x}}$ ($x>0$) dưới dạng luỹ thừa với số mũ hữu tỷ.
	\end{enumerate}
	 \loigiai{
		\begin{enumerate}
			\item  Ta có $P=a^{\frac{3}{5}} \cdot \sqrt[3]{a^2}=a^{\frac{3}{5}} \cdot a^{\frac{2}{3}}=a^{\frac{3}{5}+\frac{2}{3}}=a^{\frac{19}{15}}$.
			\item $P=\sqrt[3]{x  \cdot \sqrt[4]{x}}=\sqrt[3]{x \cdot x^{\frac{1}{4}}}=\left(x^{\frac{5}{4}}\right)^{\frac{1}{3}}=x^{\frac{5}{12}}$.
		\end{enumerate}
	 }
\end{ex}

% nguồn: Ok_TuanLuan_Toan11_Lan2
\begin{ex}%[1D6H1-2]%[ma:f32d6654]
	Cho $x>0$, rút gọn biểu thức $A=\dfrac{x^{\tfrac{4}{5}} \cdot \sqrt[6]{x^5}}{\sqrt[3]{x \cdot  \sqrt{x}}}$.
	 \loigiai{Ta có
		$A=\dfrac{x^{\tfrac{4}{5}} \cdot x^{\tfrac{5}{6}}}{\left(  x\cdot x^{\tfrac{1}{2}}\right)^{\tfrac{1}{3}}}
		=\dfrac{x^{\tfrac{4}{5}+\tfrac{5}{6}}}{x^{\tfrac{1}{3}}\cdot  x^{\tfrac{1}{6}}}
		=x^{\tfrac{4}{5}+\tfrac{5}{6}-\tfrac{1}{3}-\tfrac{1}{6}}=x^{\tfrac{17}{15}}$.\\ 
	} 
\end{ex}

% nguồn: Ok_TuanLuan_Toan11_Lan2
\begin{ex}%[1D6H1-2]%[ma:db9a022a]
	Rút gọn các biểu  thức:
	\begin{enumerate}
		\item  $P=a^{\sqrt{2}}\cdot\left( \dfrac{1}{a}\right)^{\sqrt{2}-1}$ với $a>0$.
		\item  $Q=a^{\tfrac{7}{3}}:\sqrt[3]a$ với $a>0$.
		\item $K=\dfrac{x^{\tfrac{1}{3}}\cdot \sqrt[6]{x}}{\sqrt[4]{x}}$ với $x>0$. 
		\item $M=x^{\tfrac{1}{3}}\cdot \sqrt[6]{x}$ với $x>0$. 
	\end{enumerate}
	\loigiai{ 
		\begin{enumerate}
			\item $P=a^{\sqrt{2}}\cdot\left( \dfrac{1}{a}\right)^{\sqrt{2}-1}=a^{\sqrt{2}}\cdot a^{-\sqrt{2}+1}=a^{\sqrt{2}-\sqrt{2}+1}=a$. 
			\item $Q=a^{\tfrac{7}{3}}:\sqrt[3]a=a^{\tfrac{7}{3}}:a^{\tfrac{1}{3}}=a^{\tfrac{7}{3}-\tfrac{1}{3}}=a^2$. 
			\item $K=\dfrac{x^{\tfrac{1}{3}}\cdot  \sqrt[6]{x}}{\sqrt[4]{x}}=\dfrac{x^{\tfrac{1}{3}}\cdot x^{\tfrac{1}{6}}}{x^{\tfrac{1}{4}}}=x^{\tfrac{1}{3}+\tfrac{1}{6}-\tfrac{1}{4}}=x^{\tfrac{1}{4}}=\sqrt[4]{x}$.
			\item  $M=x^{\tfrac{1}{3}}\cdot \sqrt[6]{x}=x^{\tfrac{1}{3}}\cdot x^{\tfrac{1}{6}}=x^{\tfrac{1}{2}}=\sqrt{x}$.
		\end{enumerate}
	} 
\end{ex}

% nguồn: Ok_TuanLuan_Toan11_Lan2
\begin{ex}%[1D6H1-2]%[ma:5f376f8f]
	 Tính giá trị biểu thức $A=\left(\dfrac{1}{81}\right)^{-0.75}+\left(\dfrac{1}{625}\right)^{-0{,}25}-\left(\dfrac{1}{32}\right)^{-0{,}6}$.
	\loigiai{ 
		Ta  có \allowdisplaybreaks
		\begin{eqnarray*}
			A&=&\left(\dfrac{1}{81}\right)^{-0{,}75}+\left(\dfrac{1}{625}\right)^{-0{,}25}-\left(\dfrac{1}{32}\right)^{-0{,}6}\\  &=&\left(\dfrac{1}{81}\right)^{-\tfrac{3}{4}}+\left(\dfrac{1}{625}\right)^{-\tfrac{1}{4}}-\left(\dfrac{1}{32}\right)^{-\tfrac{3}{5}}\\
			&=&\left(3^{-4}\right)^{-\tfrac{3}{4}}+\left(5^{-4}\right)^{-\tfrac{1}{4}}-\left(2^{-5}\right)^{-\tfrac{3}{5}}\\  &=&3^3+5-2^3=24.
		\end{eqnarray*}
	 }
\end{ex}

% nguồn: Ok_TuanLuan_Toan11_Lan2
\begin{ex}%[1D6H1-2]%[ma:e31ba585]
	Giá trị của biểu  thức $A=\left(a+1\right)^{-1}+\left(b+1\right)^{-1}$ và $a=\left(2+\sqrt{3}\right)^{-1}$ và $b=\left(2-\sqrt{3}\right)^{-1}$ là bao nhiêu?
	 \loigiai{
	$A=\left(a+1\right)^{-1}+\left(b+1\right)^{-1}=\left(2+\sqrt{3}+1\right)^{-1}+\left(2-\sqrt{3}+1\right)^{-1}$ $=\dfrac{1}{3+\sqrt{3}}+\dfrac{1}{3-\sqrt{3}}$  $=1$.}
\end{ex}

% nguồn: Ok_TuanLuan_Toan11_Lan2
\begin{ex}%[1D6H1-2]%[ma:7d8ce06f]
	Cho $ 0<a<1$ và $ b>1$. Biết rằng biểu thức$\sqrt{\left(a^{\pi}+b^{\pi}\right)^2-\left(4^{\tfrac{1}{\pi}}ab\right)^{\pi}}=m{a^{\pi}}+n{b^{\pi}}$, với $m$, $n\in\mathbb{Z}$. Giá  trị của $ m+n$ bằng bao nhiêu?
	 \loigiai{
		Ta có $\sqrt{\left(a^{\pi}+b^{\pi}\right)^2-\left(4^{\tfrac{1}{\pi}}ab\right)^{\pi}}=\sqrt{a^{2\pi}+b^{2\pi}+2a^{\pi}{b^{\pi}}-4a^{\pi}{b^{\pi}}}=\sqrt{a^{2\pi}+b^{2\pi}-2a^{\pi}{b^{\pi}}}=\left|a^{\pi}-b^{\pi}\right|$.\\ 
		Vì  $ 0<a<1$ và $ b>1$ nên $a^{\pi}<b^{\pi}$ do đó $\left|a^{\pi}-b^{\pi}\right|=b^{\pi}-a^{\pi}$.\\
		Suy ra $\sqrt{\left(a^{\pi}+b^{\pi}\right)^2-\left(4^{\tfrac{1}{\pi}}ab\right)^{\pi}}=b^{\pi}-a^{\pi}$ suy ra $  m=-1$; $n=1$.\\
	 Vậy $ m+n=0$.}
\end{ex}

% nguồn: Ok_TuanLuan_Toan11_Lan1
\begin{bt}%[1D6V1-2]%[ma:4b9fb83f]
	Rút  gọn biểu thức sau:
	$$P=\left[\dfrac{\sqrt[3]{a^2b}-\sqrt[3]{ab^2}}{\sqrt[3]{a^2}-2\sqrt[3]{ab}+\sqrt[3]{b^2}}-\dfrac{a+b}{\sqrt[3]{a^2}-\sqrt[3]{b^2}}\right]\big(\sqrt[6]{a}+\sqrt[6]{b}\big)^{-1}+\sqrt[6]{a},\text{ với }a>0,\,b>0,\,a\ne b$$ 
	 \loigiai{
		\begin{eqnarray*}
			&P&=\left[\dfrac{\sqrt[3]{a^2b}-\sqrt[3]{ab^2}}{\sqrt[3]{a^2}-2\sqrt[3]{ab}+\sqrt[3]{b^2}}-\dfrac{a+b}{\sqrt[3]{a^2}-\sqrt[3]{b^2}}\right]\big(\sqrt[6]{a}+\sqrt[6]{b}\big)^{-1}+\sqrt[6]{a}\\ 
			 &&=\left[\dfrac{\sqrt[3]{ab}\big(\sqrt[3]{a}-\sqrt[3]{b}\big)}{\big(\sqrt[3]{a}-\sqrt[3]{b}\big)^2}-\dfrac{\big(\sqrt[3]{a}+\sqrt[3]{b}\big)\left(\sqrt[3]{a^2}-\sqrt[3]{ab}+\sqrt[3]{b^2}\right)}{\big(\sqrt[3]{a}-\sqrt[3]{b}\big)\big(\sqrt[3]{a}+\sqrt[3]{b}\big)}\right]\big(\sqrt[6]{a}+\sqrt[6]{b}\big)^{-1}+\sqrt[6]{a}\\
			&&=\left[\dfrac{\sqrt[3]{ab}}{\sqrt[3]{a}-\sqrt[3]{b}}-\dfrac{\sqrt[3]{a^2}-\sqrt[3]{ab}+\sqrt[3]{b^2}}{\sqrt[3]{a}-\sqrt[3]{b}}\right]\dfrac{1}{\sqrt[6]{a}+\sqrt[6]{b}}+\sqrt[6]{a}\\ 
			&&=-\dfrac{\big(\sqrt[3]{a}-\sqrt[3]{b}\big)^2}{\sqrt[3]{a}-\sqrt[3]{b}}\cdot\dfrac{1}{\sqrt[6]{a}+\sqrt[6]{b}}+\sqrt[6]{a}\\ 
			 &&=\dfrac{\sqrt[3]{b}-\sqrt[3]{a}}{\sqrt[6]{a}+\sqrt[6]{b}}+\sqrt[6]{a}=\dfrac{\big(\sqrt[6]{a}+\sqrt[6]{b}\big)\sqrt[6]{a}-\sqrt[6]{b}\big)}{\sqrt[6]{a}+\sqrt[6]{b}}+\sqrt[6]{a}=\sqrt[6]{b}-\sqrt[6]{a}+\sqrt[6]{a}=\sqrt[6]{b}.
		\end{eqnarray*}
	 }
\end{bt}

% nguồn: Ok_TuanLuan_Toan11_Lan1
\begin{bt}%[1D6V1-2]%[ma:46775355]
	Rút gọn  biểu thức sau:\\
	$P=\left[\dfrac{\sqrt[{3}]{a^2 b}-\sqrt[{3}]{ab^2}}{\sqrt[{3}]{a^2}-2\sqrt[{3}]{ab}+\sqrt[{3}]{b^2}}-\dfrac{a+b}{\sqrt[{3}]{a^2}-\sqrt[{3}]{b^2}} \right](\sqrt[{6}]{a}+\sqrt[{6}]{b})^{-1}+\sqrt[{6}]{a}, \;$ với $\; a > 0,b > 0,a\ne b$. 
	\loigiai{ 
		 Ta có:
		\begin{eqnarray*}
			P&=&\left[\dfrac{\sqrt[{3}]{a^2 b}-\sqrt[{3}]{ab^2}}{\sqrt[{3}]{a^2}-2\sqrt[{3}]{ab}+\sqrt[{3}]{b^2}}-\dfrac{a+b}{\sqrt[{3}]{a^2}-\sqrt[{3}]{b^2}} \right](\sqrt[{6}]{a}+\sqrt[{6}]{b})^{-1}+\sqrt[{6}]{a} \\ 
			&=&\left[\dfrac{\sqrt[{3}]{ab} (\sqrt[{3}]{a}-\sqrt[{3}]{b})}{(\sqrt[{3}]{a}-\sqrt[{3}]{b})^2}-\dfrac{(\sqrt[{3}]{a}+\sqrt[{3}]{b})\left(\sqrt[{3}]{a^2}-\sqrt[{3}]{ab}+\sqrt[{3}]{b^2}  \right)}{(\sqrt[{3}]{a}-\sqrt[{3}]{b})(\sqrt[{3}]{a}+\sqrt[{3}]{b})} \right](\sqrt[{6}]{a}+\sqrt[{6}]{b})^{-1}+\sqrt[{6}]{a} \\
			&=&\left[\dfrac{\sqrt[{3}]{ab}}{\sqrt[{3}]{a}-\sqrt[{3}]{b}}-\dfrac{\sqrt[{3}]{a^2}-\sqrt[{3}]{ab}+\sqrt[{3}]{b^2}}{\sqrt[{3}]{a}-\sqrt[{3}]{b}} \right]\dfrac{1}{\sqrt[{6}]{a}+\sqrt[{6}]{b}}+\sqrt[{6}]{a}\\ 
			&=&-\dfrac{(\sqrt[{3}]{a}-\sqrt[{3}]{b})^2}{\sqrt[{3}]{a}-\sqrt[{3}]{b}} \dfrac{1}{(\sqrt[{6}]{a}+\sqrt[{6}]{b})}+\sqrt[{6}]{a}\\ 
			&=&\dfrac{\sqrt[{3}]{b}-\sqrt[{3}]{a}}{\sqrt[{6}]{a}+\sqrt[{6}]{b}}+\sqrt[{6}]{a}=\sqrt[{6}]{b}-\sqrt[{6}]{a}+\sqrt[{6}]{a}=\sqrt[{6}]{b}. 
		\end{eqnarray*}
	 }
\end{bt}

% nguồn: Ok_TuanLuan_Toan11_Lan2
\begin{ex}%[1D6V1-2]%[ma:8173ef8d]
	Rút gọn biểu thức sau: 
	$$P=\left[\dfrac{\sqrt[3]{a^2b}-\sqrt[3]{ab^2}}{\sqrt[3]{a^2}-2\sqrt[3]{ab}+\sqrt[3]{b^2}}-\dfrac{a+b}{\sqrt[3]{a^2}-\sqrt[3]{b^2}}\right]\big(\sqrt[6]{a}+\sqrt[6]{b}\big)^{-1}+\sqrt[6]{a},\text{ với  }a>0,\,b>0,\,a\ne b$$
	\loigiai{ 
		\begin{eqnarray*}
			&P&=\left[\dfrac{\sqrt[3]{a^2b}-\sqrt[3]{ab^2}}{\sqrt[3]{a^2}-2\sqrt[3]{ab}+\sqrt[3]{b^2}}-\dfrac{a+b}{\sqrt[3]{a^2}-\sqrt[3]{b^2}}\right]\big(\sqrt[6]{a}+\sqrt[6]{b}\big)^{-1}+\sqrt[6]{a}\\ 
			&&=\left[\dfrac{\sqrt[3]{ab}\big(\sqrt[3]{a}-\sqrt[3]{b}\big)}{\big(\sqrt[3]{a}-\sqrt[3]{b}\big)^2}-\dfrac{\big(\sqrt[3]{a}+\sqrt[3]{b}\big)\left(\sqrt[3]{a^2}-\sqrt[3]{ab}+\sqrt[3]{b^2}\right)}{\big(\sqrt[3]{a}-\sqrt[3]{b}\big)\big(\sqrt[3]{a}+\sqrt[3]{b}\big)}\right]\big(\sqrt[6]{a}+\sqrt[6]{b}\big)^{-1}+\sqrt[6]{a}\\ 
			&&=\left[\dfrac{\sqrt[3]{ab}}{\sqrt[3]{a}-\sqrt[3]{b}}-\dfrac{\sqrt[3]{a^2}-\sqrt[3]{ab}+\sqrt[3]{b^2}}{\sqrt[3]{a}-\sqrt[3]{b}}\right]\dfrac{1}{\sqrt[6]{a}+\sqrt[6]{b}}+\sqrt[6]{a}\\ 
			&&=-\dfrac{\big(\sqrt[3]{a}-\sqrt[3]{b}\big)^2}{\sqrt[3]{a}-\sqrt[3]{b}}\cdot\dfrac{1}{\sqrt[6]{a}+\sqrt[6]{b}}+\sqrt[6]{a}\\ 
			&&=\dfrac{\sqrt[3]{b}-\sqrt[3]{a}}{\sqrt[6]{a}+\sqrt[6]{b}}+\sqrt[6]{a}=\dfrac{\big(\sqrt[6]{a}+\sqrt[6]{b}\big)\sqrt[6]{a}-\sqrt[6]{b}\big)}{\sqrt[6]{a}+\sqrt[6]{b}}+\sqrt[6]{a}=\sqrt[6]{b}-\sqrt[6]{a}+\sqrt[6]{a}=\sqrt[6]{b}. 
		\end{eqnarray*}
	 }
\end{ex}

% nguồn: Ok_TuanLuan_Toan11_Lan2
\begin{ex}%[1D6V1-2]%[ma:9cbad613]
	Rút gọn biểu thức  sau:\\
	$P=\left[\dfrac{\sqrt[{3}]{a^2  b}-\sqrt[{3}]{ab^2}}{\sqrt[{3}]{a^2}-2\sqrt[{3}]{ab}+\sqrt[{3}]{b^2}}-\dfrac{a+b}{\sqrt[{3}]{a^2}-\sqrt[{3}]{b^2}} \right](\sqrt[{6}]{a}+\sqrt[{6}]{b})^{-1}+\sqrt[{6}]{a}, \;$ với $\; a > 0,b > 0,a\ne b$.
	\loigiai{ 
		Ta  có:
		\begin{eqnarray*}
			P&=&\left[\dfrac{\sqrt[{3}]{a^2 b}-\sqrt[{3}]{ab^2}}{\sqrt[{3}]{a^2}-2\sqrt[{3}]{ab}+\sqrt[{3}]{b^2}}-\dfrac{a+b}{\sqrt[{3}]{a^2}-\sqrt[{3}]{b^2}} \right](\sqrt[{6}]{a}+\sqrt[{6}]{b})^{-1}+\sqrt[{6}]{a}  \\
			&=&\left[\dfrac{\sqrt[{3}]{ab} (\sqrt[{3}]{a}-\sqrt[{3}]{b})}{(\sqrt[{3}]{a}-\sqrt[{3}]{b})^2}-\dfrac{(\sqrt[{3}]{a}+\sqrt[{3}]{b})\left(\sqrt[{3}]{a^2}-\sqrt[{3}]{ab}+\sqrt[{3}]{b^2} \right)}{(\sqrt[{3}]{a}-\sqrt[{3}]{b})(\sqrt[{3}]{a}+\sqrt[{3}]{b})} \right](\sqrt[{6}]{a}+\sqrt[{6}]{b})^{-1}+\sqrt[{6}]{a} \\ 
			 &=&\left[\dfrac{\sqrt[{3}]{ab}}{\sqrt[{3}]{a}-\sqrt[{3}]{b}}-\dfrac{\sqrt[{3}]{a^2}-\sqrt[{3}]{ab}+\sqrt[{3}]{b^2}}{\sqrt[{3}]{a}-\sqrt[{3}]{b}} \right]\dfrac{1}{\sqrt[{6}]{a}+\sqrt[{6}]{b}}+\sqrt[{6}]{a}\\
			&=&-\dfrac{(\sqrt[{3}]{a}-\sqrt[{3}]{b})^2}{\sqrt[{3}]{a}-\sqrt[{3}]{b}} \dfrac{1}{(\sqrt[{6}]{a}+\sqrt[{6}]{b})}+\sqrt[{6}]{a}\\ 
			&=&\dfrac{\sqrt[{3}]{b}-\sqrt[{3}]{a}}{\sqrt[{6}]{a}+\sqrt[{6}]{b}}+\sqrt[{6}]{a}=\sqrt[{6}]{b}-\sqrt[{6}]{a}+\sqrt[{6}]{a}=\sqrt[{6}]{b}. 
		\end{eqnarray*}
	} 
\end{ex}
